% The Modern AI Concept Book, Edition 2 -- Overleaf-ready (pdfLaTeX)
\documentclass[10pt,a4paper,openany]{report}
\usepackage[T1]{fontenc}
\usepackage[utf8]{inputenc}
\usepackage{mathpazo}
\usepackage[scaled=0.92]{helvet}
\usepackage{courier}
\usepackage{microtype}
\usepackage[a4paper,margin=2.2cm,top=2.3cm,bottom=2.5cm]{geometry}
\usepackage[dvipsnames,table]{xcolor}
\usepackage{titlesec}
\usepackage{fancyhdr}
\usepackage{needspace}
\usepackage{enumitem}
\usepackage{multicol}
\usepackage{longtable}
\usepackage{array}
\usepackage{tikz}
\usepackage[most]{tcolorbox}
\usepackage{xurl}
\usepackage[hidelinks]{hyperref}

% ---- palette (edit here to re-theme the whole book) ---------------------
\definecolor{ink}{HTML}{1D2124}
\definecolor{muted}{HTML}{5E666C}
\definecolor{accent}{HTML}{2F5D8A}
\definecolor{accentbg}{HTML}{E9EFF6}
\definecolor{rule}{HTML}{D7DAD3}
\definecolor{newc}{HTML}{B0431F}
\definecolor{kG}{HTML}{24292F}\definecolor{kB}{HTML}{8A4B08}\definecolor{kP}{HTML}{7A2E7A}
\definecolor{kD}{HTML}{1F6B4F}\definecolor{kS}{HTML}{A3261E}\definecolor{kV}{HTML}{2F5D8A}
\hypersetup{colorlinks=true,linkcolor=accent,urlcolor=accent,pdftitle={The Modern AI Concept Book},pdfauthor={AI Concept Book}}
\color{ink}
\setlength{\parindent}{0pt}\setlength{\parskip}{3pt}

\titleformat{\chapter}[display]{\sffamily\bfseries\color{ink}}{\color{accent}\fontsize{50}{50}\selectfont\thechapter}{4pt}{\huge}[\vspace{3pt}{\color{rule}\titlerule[1pt]}]
\titlespacing*{\chapter}{0pt}{-14pt}{14pt}
\titleformat{\section}{\sffamily\bfseries\large}{}{0pt}{}
\setcounter{secnumdepth}{0}\setcounter{tocdepth}{1}
\pagestyle{fancy}\fancyhf{}
\fancyhead[L]{\sffamily\footnotesize\color{muted}The Modern AI Concept Book}
\fancyhead[R]{\sffamily\footnotesize\color{muted}\leftmark}
\fancyfoot[C]{\sffamily\footnotesize\thepage}
\renewcommand{\headrulewidth}{0.3pt}
\renewcommand{\chaptermark}[1]{\markboth{#1}{}}
\fancypagestyle{plain}{\fancyhf{}\fancyfoot[C]{\sffamily\footnotesize\thepage}\renewcommand{\headrulewidth}{0pt}}

% ---- concept macros -----------------------------------------------------
% \concept{label}{Name}{alias}{meta line}
\newcommand{\concept}[4]{\Needspace{8\baselineskip}\vspace{9pt}\phantomsection\label{#1}\addcontentsline{toc}{section}{#2}%
 {\sffamily\bfseries\large #2}\if\relax\detokenize{#3}\relax\else\ \ {\sffamily\small\color{accent}(#3)}\fi\par\nopagebreak
 {\sffamily\scriptsize\color{muted}#4}\par\nopagebreak\vspace{2pt}}
\newcommand{\newtag}{\colorbox{newc}{\sffamily\scriptsize\bfseries\color{white}NEW 2026}}
\newcommand{\acttag}{\fcolorbox{newc}{white}{\sffamily\scriptsize\color{newc}ACTIVE 2026}}
\newenvironment{keypoints}{\begin{itemize}[leftmargin=14pt,itemsep=1pt,topsep=2pt,parsep=0pt]\small}{\end{itemize}}
\newtcolorbox{sota}{enhanced,breakable,colback=accentbg,colframe=accent,boxrule=0pt,leftrule=2.5pt,sharp corners,left=6pt,right=6pt,top=3pt,bottom=3pt,fontupper=\small,
 before upper={{\sffamily\scriptsize\bfseries\color{accent}STATE OF THE ART, SEPT 2026}\par}}
\newcommand{\lk}[5]{\item[\textcolor{k#1}{\rule{5pt}{5pt}}]{\sffamily\scriptsize\bfseries\textcolor{k#1}{#2}}\ {\sffamily\scriptsize\color{muted}#3}\ \href{#5}{#4}\\[-1pt]{\ttfamily\scriptsize\color{muted}\url{#5}}}
\newenvironment{links}{\begin{itemize}[leftmargin=14pt,itemsep=1pt,topsep=3pt,labelsep=6pt]\small}{\end{itemize}}
\newcommand{\chapintro}[1]{{\itshape\color{muted}\large #1\par}\vspace{4pt}}

\begin{document}
\begin{titlepage}
\begin{tikzpicture}[remember picture,overlay]\fill[accent] (current page.north west) rectangle ([yshift=-8cm]current page.north east);\end{tikzpicture}
\vspace*{0.3cm}
{\color{white}\sffamily\bfseries\fontsize{42}{44}\selectfont The Modern AI\\[4pt] Concept Book\par}
\vspace{0.4cm}{\color{white}\sffamily\large Edition 2 \textperiodcentered{} September 2026\par}
\vspace{2.4cm}
{\Large LLMs, VLMs, agents, harnesses, hardware and the practice of AI engineering.\par}
\vspace{0.4cm}{\color{muted}Current to September 20, 2026.\par}
\vfill
{\sffamily\small\color{muted}17 categories \quad 173 concepts \quad 430 dated resources\\[6pt]
Resource types:\ \textcolor{kG}{\rule{5pt}{5pt}} GitHub\quad \textcolor{kB}{\rule{5pt}{5pt}} Blog\quad \textcolor{kP}{\rule{5pt}{5pt}} Paper\quad \textcolor{kD}{\rule{5pt}{5pt}} Docs\quad \textcolor{kS}{\rule{5pt}{5pt}} Spec\quad \textcolor{kV}{\rule{5pt}{5pt}} Video\par}
\end{titlepage}

\chapter*{How to use this book}\markboth{How to use this book}{}
Each entry gives a plain-language definition, key points, a \emph{state of the art} note where the picture changed in 2026, and dated resources (GitHub repositories, blogs, papers, documentation, specifications and videos).

The header line under each concept shows the year it emerged, its most recent active year, its level (Beginner, Intermediate, Advanced) and tags. \newtag\ marks concepts that appeared in 2026; \acttag\ marks older concepts that were still changing in 2026.

A printed book cannot be filtered, so the appendices provide the same cuts as the interactive HTML edition: concepts by year of emergence, concepts by level, all resources grouped by type and sorted by year, and an index of acronyms. For live filtering by category, year, level, tag, resource type and resource year, open the HTML edition or the spreadsheet catalog.

Resource years show when an article, paper or specification was published, or when a repository or site launched (approximate for long-running projects). Chapter 1 reflects public reporting as of September 20, 2026 and will date quickly.

\tableofcontents
\clearpage
\chapter{The Frontier Landscape (September 2026)}
\chapintro{A snapshot of where frontier and open-weight models stand as of September 2026. This chapter dates quickly; use the trackers at the end to stay current.}
\concept{c:landscape-0}{Mythos-class tier and trusted access}{Fable / Mythos}{\newtag\ Since 2026 \quad Beginner \quad Enterprise, Security}
Anthropic introduced a capability tier above Opus. The same underlying model ships in two forms: Claude Fable for general use with extra safeguards in biology, cybersecurity and LLM R\&D, and Claude Mythos for vetted defenders and partners through Project Glasswing.
\begin{keypoints}
  \item Mythos Preview (April 2026) was first restricted to trusted organisations for cyber defense.
  \item Fable 5 / Mythos 5 launched June 9, 2026; access was suspended June 12 to comply with US export controls and restored July 1 after the controls were lifted.
  \item Fable 5.1 and Mythos 5.1 followed on September 1, 2026, with 1M-token context.
\end{keypoints}
\begin{sota}Fable 5.1 is Anthropic's generally available flagship; Opus 5 (July 2026) and Sonnet 5 sit below it.\end{sota}
\begin{links}
  \lk{D}{Docs}{2026}{Anthropic: Project Glasswing}{https://www.anthropic.com/glasswing}
  \lk{D}{Docs}{2026}{Anthropic statement on Fable and Mythos access}{https://www.anthropic.com/news/fable-mythos-access}
  \lk{B}{Blog}{2026}{DataNorth: Claude Fable 5 and Mythos 5 explained}{https://datanorth.ai/news/anthropic-releases-claude-fable-5-and-mythos-5}
\end{links}

\concept{c:landscape-1}{GPT-6 Astra and the Critical cyber threshold}{GPT-6}{\newtag\ Since 2026 \quad Beginner \quad Security, Enterprise}
OpenAI's first GPT-6 generation model, released September 3, 2026. It is the first model OpenAI classified at the Critical cybersecurity level of its Preparedness Framework, so it launched in phases with stronger safeguards and a trusted-access program for defenders.
\begin{keypoints}
  \item Critical means the model can find and exploit unknown vulnerabilities in hardened systems without step-by-step human guidance.
  \item Public versions refuse advanced offensive tasks; vetted defenders get more through OpenAI's trusted-access programs.
  \item Its system card also reported reduced chain-of-thought monitorability, a notable safety signal.
\end{keypoints}
\begin{sota}GPT-6 Astra sits above the GPT-5.6 Sol / Terra / Luna tiers.\end{sota}
\begin{links}
  \lk{B}{Blog}{2026}{OpenAI: Path to Astra}{https://openai.com/index/path-to-astra/}
  \lk{D}{Docs}{2026}{GPT-6 Astra system card}{https://deploymentsafety.openai.com/gpt-6-astra}
  \lk{B}{Blog}{2026}{OpenAI: Safety overview of GPT-6 Astra}{https://openai.com/index/safety-overview-gpt-6-astra/}
  \lk{B}{Blog}{2026}{InfoQ: Astra is first model rated Critical for cyber}{https://www.infoq.com/news/2026/09/gpt-6-astra-critical-cyber/}
\end{links}

\concept{c:landscape-2}{Tiered access for dual-use models}{trusted access}{\acttag\ Since 2025--2026 \quad Intermediate \quad Security, Enterprise}
A 2026 industry pattern: the most cyber-capable models are released in a public version with guardrails and a restricted version for vetted security teams and critical-infrastructure operators.
\begin{keypoints}
  \item Anthropic (Glasswing), OpenAI (trusted access for cyber) and Google (defender-only cyber variants) all adopted versions of it.
  \item It lets defenders patch vulnerabilities before equivalent capabilities spread.
  \item Enterprises now need to track which model variant, safeguards and data-retention terms they are using.
\end{keypoints}
\begin{links}
  \lk{D}{Docs}{2026}{Anthropic: Project Glasswing}{https://www.anthropic.com/glasswing}
  \lk{B}{Blog}{2026}{CSO Online: Astra crosses the Critical cyber threshold}{https://www.csoonline.com/article/4218679/openai-launches-gpt-6-astra-its-first-model-to-cross-a-critical-cybersecurity-threshold.html}
\end{links}

\concept{c:landscape-3}{Open-weight frontier models}{}{\acttag\ Since 2023--2026 \quad Beginner \quad Open source / open weights, Cost \& efficiency}
Open-weight models trail the closed frontier by months, not years. 2026's leading families are mostly large sparse MoEs with 1M-token context and permissive licences.
\begin{keypoints}
  \item DeepSeek V4 (April 2026): Pro has 1.6T total / 49B active parameters, Flash has 284B / 13B, both with 1M context, MIT licence.
  \item Other active families: Qwen 3.x, Kimi, GLM-5.x, MiniMax, Mistral Large 3, NVIDIA Nemotron 3, Google Gemma 4.
  \item Open weight does not always mean open source: check licence terms and whether data and code are released.
\end{keypoints}
\begin{sota}DeepSeek V4 is the reference open model for long-context efficiency; small open models now run strong coding workloads on a single GPU.\end{sota}
\begin{links}
  \lk{P}{Paper}{2026}{DeepSeek-V4 technical report}{https://arxiv.org/abs/2606.19348}
  \lk{B}{Blog}{2026}{HF blog: DeepSeek V4 architecture analysis}{https://huggingface.co/blog/ResterChed/deepseek-v4-ga-architecture}
  \lk{D}{Docs}{2016}{Hugging Face models hub}{https://huggingface.co/models}
\end{links}

\concept{c:landscape-4}{Model release trackers}{}{\acttag\ Since 2023--2026 \quad Beginner \quad Production practice}
In 2026 frontier labs shipped new models every few weeks. Trackers and independent benchmark aggregators are the practical way to keep up.
\begin{keypoints}
  \item Compare on your own evals, not headline benchmarks.
  \item Watch pricing changes as closely as capability changes: cache-read and small-tier prices fell sharply in 2026.
\end{keypoints}
\begin{links}
  \lk{D}{Docs}{2024}{Artificial Analysis}{https://artificialanalysis.ai}
  \lk{D}{Docs}{2024}{Epoch AI benchmarking hub}{https://epoch.ai/benchmarks}
  \lk{D}{Docs}{2023}{LMArena}{https://lmarena.ai}
  \lk{D}{Docs}{2023}{OpenRouter model list}{https://openrouter.ai/models}
\end{links}

\chapter{Agents and Harnesses}
\chapintro{An agent is a model running in a loop: it reads context, chooses an action, calls a tool, observes the result and repeats. The industry now summarises this as Agent = Model + Harness, and most 2026 engineering effort goes into the harness.}
\concept{c:agents-0}{Agent harness}{}{\acttag\ Since 2024--2026 \quad Beginner \quad Foundational, Production practice}
Everything around the model that turns it into a working agent: the loop, tool definitions, permissions, sandbox, context management, memory, retries, logging and human approval gates.
\begin{keypoints}
  \item Two teams using the same model can get very different results purely from harness design.
  \item Harness setup alone can move benchmark scores by several points.
  \item Harnesses encode assumptions about model weaknesses, so they must be revisited when models improve.
\end{keypoints}
\begin{sota}Vendor harnesses (Claude Code, Codex, Gemini CLI) are now exposed as SDKs and managed runtimes so teams can build on them directly.\end{sota}
\begin{links}
  \lk{B}{Blog}{2026}{Martin Fowler: Harness engineering for coding agent users}{https://martinfowler.com/articles/harness-engineering.html}
  \lk{B}{Blog}{2025}{Anthropic: Effective harnesses for long-running agents}{https://www.anthropic.com/engineering/effective-harnesses-for-long-running-agents}
  \lk{G}{GitHub}{2026}{awesome-harness-engineering}{https://github.com/ai-boost/awesome-harness-engineering}
\end{links}

\concept{c:agents-1}{Harness engineering}{}{\newtag\ Since 2026 \quad Intermediate \quad Production practice}
The discipline that followed prompt engineering and context engineering: designing environments, constraints, documentation-as-system-of-record, structural tests and feedback loops so agents do reliable work while humans steer.
\begin{keypoints}
  \item OpenAI described a team shipping a roughly million-line product with no hand-written code using this approach.
  \item Core moves: constrain what agents can do, inform them what to do, verify their work, correct mistakes, keep humans at high-stakes gates.
  \item Architecture rules become linters and tests the agent must pass, not wiki pages.
\end{keypoints}
\begin{sota}The term entered mainstream use in early 2026 and is now the main focus of agent engineering investment.\end{sota}
\begin{links}
  \lk{B}{Blog}{2026}{OpenAI: Harness engineering}{https://openai.com/index/harness-engineering/}
  \lk{B}{Blog}{2026}{Faros: Harness engineering guide}{https://www.faros.ai/blog/harness-engineering}
  \lk{G}{GitHub}{2026}{awesome-agent-harness}{https://github.com/AutoJunjie/awesome-agent-harness}
\end{links}

\concept{c:agents-2}{Managed agents: decoupling brain, hands and session}{}{\newtag\ Since 2026 \quad Advanced \quad Production practice, Enterprise}
An architecture where the reasoning loop (brain), the sandboxes and tools that act (hands), and the append-only event log (session) are separate, replaceable services.
\begin{keypoints}
  \item The harness becomes stateless; containers are provisioned only when a tool needs them.
  \item A crash in one sandbox does not lose the session, because state lives in the log.
  \item Enterprises can run the hands inside their own network while the brain runs at the provider.
\end{keypoints}
\begin{sota}Anthropic launched Claude Managed Agents on this design in April 2026; Cursor's cloud agents converged on a similar three-layer split.\end{sota}
\begin{links}
  \lk{B}{Blog}{2026}{Anthropic: Scaling Managed Agents}{https://anthropic.com/engineering/managed-agents}
\end{links}

\concept{c:agents-3}{Agentic loop / ReAct}{ReAct}{Since 2022--2025 \quad Beginner \quad Foundational}
The basic control pattern: reason about the next step, act by calling a tool, observe the result, and loop until done.
\begin{keypoints}
  \item Modern models interleave thinking between tool calls natively.
  \item A usable coding agent loop fits in about 100 lines of code.
\end{keypoints}
\begin{links}
  \lk{P}{Paper}{2022}{ReAct paper}{https://arxiv.org/abs/2210.03629}
  \lk{B}{Blog}{2024}{Anthropic: Building effective agents}{https://www.anthropic.com/engineering/building-effective-agents}
  \lk{G}{GitHub}{2025}{mini-swe-agent}{https://github.com/SWE-agent/mini-swe-agent}
\end{links}

\concept{c:agents-4}{Workflows vs. agents}{}{Since 2024--2025 \quad Beginner \quad Foundational, Production practice}
Workflows run LLM calls through predefined code paths; agents let the model direct its own process. Use the simplest pattern that works.
\begin{keypoints}
  \item Common workflow patterns: prompt chaining, routing, parallelisation, orchestrator-workers, evaluator-optimiser.
  \item Add autonomy only when the path cannot be predicted in advance.
\end{keypoints}
\begin{links}
  \lk{B}{Blog}{2024}{Anthropic: Building effective agents}{https://www.anthropic.com/engineering/building-effective-agents}
  \lk{G}{GitHub}{2025}{12-factor agents}{https://github.com/humanlayer/12-factor-agents}
\end{links}

\concept{c:agents-5}{Tool use / function calling}{}{Since 2023--2025 \quad Beginner \quad Foundational}
Models emit structured calls (a tool name plus JSON arguments) that the harness executes and returns. Tool design is API design for a non-human user.
\begin{keypoints}
  \item Prefer few, expressive tools over many narrow ones.
  \item Return compact, high-signal output and actionable error messages.
  \item Namespacing and clear descriptions reduce wrong-tool calls.
\end{keypoints}
\begin{links}
  \lk{B}{Blog}{2025}{Anthropic: Writing effective tools for agents}{https://www.anthropic.com/engineering/writing-tools-for-agents}
  \lk{D}{Docs}{2024}{Claude tool use docs}{https://docs.claude.com/en/docs/agents-and-tools/tool-use/overview}
\end{links}

\concept{c:agents-6}{Tool search and programmatic tool calling}{code mode}{\acttag\ Since 2025--2026 \quad Intermediate \quad Cost \& efficiency, Production practice}
Two fixes for tool overload. Tool search loads tool definitions on demand; programmatic tool calling lets the model write a script that calls many tools and returns only the final result.
\begin{keypoints}
  \item Cuts context use and round-trips dramatically for large tool sets.
  \item Deferred tool discovery and stable tool ordering also protect prompt-cache hits.
\end{keypoints}
\begin{sota}Now a standard feature in frontier harnesses; MCP servers are increasingly consumed through code rather than direct calls.\end{sota}
\begin{links}
  \lk{B}{Blog}{2025}{Anthropic: Advanced tool use}{https://www.anthropic.com/engineering/advanced-tool-use}
  \lk{B}{Blog}{2025}{Anthropic: Code execution with MCP}{https://www.anthropic.com/engineering/code-execution-with-mcp}
  \lk{B}{Blog}{2025}{Cloudflare: Code Mode}{https://blog.cloudflare.com/code-mode/}
\end{links}

\concept{c:agents-7}{Multi-agent systems and subagents}{orchestrator-worker}{\acttag\ Since 2023--2026 \quad Intermediate \quad Production practice}
A lead agent plans and delegates to subagents with their own clean context windows, then combines condensed results.
\begin{keypoints}
  \item Great for broad research and big codebases; costs more tokens.
  \item Specialised orchestrator models can route work between tools, small models and large models.
\end{keypoints}
\begin{sota}Agent teams, swarms and parallel agents in isolated worktrees are routine in 2026 coding tools.\end{sota}
\begin{links}
  \lk{B}{Blog}{2025}{Anthropic: Multi-agent research system}{https://www.anthropic.com/engineering/multi-agent-research-system}
  \lk{D}{Docs}{2025}{Claude Code subagents}{https://docs.claude.com/en/docs/claude-code/sub-agents}
  \lk{G}{GitHub}{2024}{LangGraph}{https://github.com/langchain-ai/langgraph}
\end{links}

\concept{c:agents-8}{Long-running and background agents}{}{\acttag\ Since 2025--2026 \quad Intermediate \quad Production practice}
Agents that work for hours or days across many context windows, usually in the cloud while the human does something else.
\begin{keypoints}
  \item Rely on plan files, progress notes, git history, test gates and compaction.
  \item Early models showed context anxiety (wrapping up early near the limit); newer models largely don't, so harness workarounds can become dead weight.
\end{keypoints}
\begin{sota}METR's time-horizon measure keeps doubling; multi-hour autonomous runs are normal for frontier models.\end{sota}
\begin{links}
  \lk{B}{Blog}{2025}{Anthropic: Effective harnesses for long-running agents}{https://www.anthropic.com/engineering/effective-harnesses-for-long-running-agents}
  \lk{B}{Blog}{2025}{METR: Measuring AI ability to complete long tasks}{https://metr.org/blog/2025-03-19-measuring-ai-ability-to-complete-long-tasks/}
\end{links}

\concept{c:agents-9}{Computer use and browser agents}{CUA}{\acttag\ Since 2024--2026 \quad Intermediate \quad Production practice}
Agents that operate a GUI from screenshots with mouse and keyboard actions, or drive a browser through the DOM and accessibility tree.
\begin{keypoints}
  \item Agentic browsers embed this directly in the browser.
  \item Prompt injection from web content is the main security risk.
\end{keypoints}
\begin{links}
  \lk{D}{Docs}{2024}{Claude computer use tool}{https://docs.claude.com/en/docs/agents-and-tools/tool-use/computer-use-tool}
  \lk{B}{Blog}{2025}{OpenAI: Computer-Using Agent}{https://openai.com/index/computer-using-agent/}
  \lk{G}{GitHub}{2024}{browser-use}{https://github.com/browser-use/browser-use}
  \lk{G}{GitHub}{2025}{Playwright MCP}{https://github.com/microsoft/playwright-mcp}
\end{links}

\concept{c:agents-10}{Personal agents}{OpenClaw}{\newtag\ Since 2026 \quad Beginner \quad Open source / open weights, Security}
Always-on assistants connected to messaging apps, email, calendars and files that take real actions for a person.
\begin{keypoints}
  \item OpenClaw's viral open-source spread in early 2026 put millions of personal agents into daily use.
  \item Exposed the permission, isolation and injection problems of agents with broad personal access.
\end{keypoints}
\begin{links}
  \lk{G}{GitHub}{2026}{OpenClaw}{https://github.com/openclaw/openclaw}
  \lk{P}{Paper}{2026}{SemaClaw: personal agents via harness engineering}{https://arxiv.org/abs/2604.11548}
\end{links}

\concept{c:agents-11}{Agent sandboxes}{}{\acttag\ Since 2024--2026 \quad Intermediate \quad Security, Production practice}
Isolated environments (containers, microVMs, V8 isolates, OS-level sandboxes) where agents run code without endangering the host.
\begin{keypoints}
  \item Filesystem boundaries and network egress controls are the two key levers.
  \item Good sandboxing reduces approval fatigue: fewer prompts, more safety.
\end{keypoints}
\begin{sota}Providers now let enterprises run agent sandboxes in their own clouds while the model runs remotely.\end{sota}
\begin{links}
  \lk{G}{GitHub}{2025}{Anthropic sandbox-runtime}{https://github.com/anthropic-experimental/sandbox-runtime}
  \lk{G}{GitHub}{2023}{E2B}{https://github.com/e2b-dev/E2B}
  \lk{G}{GitHub}{2024}{Daytona}{https://github.com/daytonaio/daytona}
  \lk{G}{GitHub}{2025}{container-use}{https://github.com/dagger/container-use}
\end{links}

\concept{c:agents-12}{Agent frameworks and SDKs}{}{\acttag\ Since 2023--2026 \quad Beginner \quad Open source / open weights}
Libraries that package the loop, tools, memory, handoffs and tracing.
\begin{keypoints}
  \item Vendor SDKs expose the same harness their own products use.
  \item Open frameworks add graphs, durable execution and multi-agent patterns.
\end{keypoints}
\begin{links}
  \lk{G}{GitHub}{2025}{Claude Agent SDK (Python)}{https://github.com/anthropics/claude-agent-sdk-python}
  \lk{G}{GitHub}{2025}{OpenAI Agents SDK}{https://github.com/openai/openai-agents-python}
  \lk{G}{GitHub}{2025}{Google ADK}{https://github.com/google/adk-python}
  \lk{G}{GitHub}{2025}{Microsoft Agent Framework}{https://github.com/microsoft/agent-framework}
  \lk{G}{GitHub}{2024}{Pydantic AI}{https://github.com/pydantic/pydantic-ai}
  \lk{G}{GitHub}{2024}{smolagents}{https://github.com/huggingface/smolagents}
  \lk{G}{GitHub}{2023}{CrewAI}{https://github.com/crewAIInc/crewAI}
  \lk{G}{GitHub}{2025}{LangChain deepagents}{https://github.com/langchain-ai/deepagents}
\end{links}

\concept{c:agents-13}{Durable execution and agent trajectories}{}{\acttag\ Since 2024--2026 \quad Advanced \quad Production practice}
Persist every step so an agent can crash, pause for approval or wait days, then resume. Trajectory logs also feed memory, evals and training.
\begin{keypoints}
  \item Workflow engines like Temporal are used as agent backbones.
  \item Libraries now normalise session transcripts from many harnesses into one format.
\end{keypoints}
\begin{links}
  \lk{D}{Docs}{2019}{Temporal}{https://temporal.io}
  \lk{G}{GitHub}{2024}{LangGraph (persistence)}{https://github.com/langchain-ai/langgraph}
\end{links}

\chapter{Protocols and Interoperability Standards}
\chapintro{Open protocols let agents, tools, apps, editors and payment systems interoperate without bespoke glue. Several now live in neutral foundations.}
\concept{c:protocols-0}{Model Context Protocol}{MCP}{\acttag\ Since 2024--2026 \quad Beginner \quad Open standard, Foundational}
An open JSON-RPC protocol connecting AI applications (clients) to tools, resources and prompts exposed by servers. The de facto standard for tool integration across vendors.
\begin{keypoints}
  \item 2026-07-28 revision: stateless core (no sessions or initialize handshake), server discovery on demand, cacheable list results, header-based routing.
  \item Formal Extensions framework; MCP Apps and Tasks ship as extensions.
  \item Authorization aligned with OAuth/OIDC: issuer validation, Client ID Metadata Documents, enterprise-managed authorization.
\end{keypoints}
\begin{sota}The 2026-07-28 spec is the largest revision since launch, with a 12-month deprecation policy and SDK conformance tiers.\end{sota}
\begin{links}
  \lk{B}{Blog}{2026}{MCP blog: The 2026-07-28 specification}{https://blog.modelcontextprotocol.io/posts/2026-07-28/}
  \lk{B}{Blog}{2026}{MCP roadmap (Aug 2026)}{https://blog.modelcontextprotocol.io/posts/mcp-roadmap/}
  \lk{D}{Docs}{2024}{modelcontextprotocol.io}{https://modelcontextprotocol.io}
  \lk{G}{GitHub}{2024}{MCP org (SDKs, servers)}{https://github.com/modelcontextprotocol}
  \lk{B}{Blog}{2026}{WorkOS: What changes for agent authentication}{https://workos.com/blog/mcp-2026-spec-agent-authentication}
\end{links}

\concept{c:protocols-1}{MCP Apps}{}{\acttag\ Since 2025--2026 \quad Intermediate \quad Open standard}
An MCP extension that lets servers return interactive UI rendered in a sandboxed frame, so tools can show forms, charts and pickers inside the chat.
\begin{keypoints}
  \item Standardised as an extension in the 2026-07-28 release.
\end{keypoints}
\begin{links}
  \lk{G}{GitHub}{2025}{MCP ext-apps}{https://github.com/modelcontextprotocol/ext-apps}
\end{links}

\concept{c:protocols-2}{MCP Tasks}{}{\acttag\ Since 2025--2026 \quad Intermediate \quad Open standard}
An extension for long-running work: a server can accept a job, stream progress and let the client check back or steer mid-flight.
\begin{keypoints}
  \item Paired with subscriptions and progress notifications for push-style updates.
\end{keypoints}
\begin{links}
  \lk{B}{Blog}{2026}{MCP release candidate notes}{https://blog.modelcontextprotocol.io/posts/2026-07-28-release-candidate/}
\end{links}

\concept{c:protocols-3}{Agent2Agent protocol}{A2A}{\acttag\ Since 2025--2026 \quad Intermediate \quad Open standard}
An open protocol under the Linux Foundation for agents from different vendors to discover each other via Agent Cards, exchange tasks and stream results.
\begin{keypoints}
  \item MCP is agent-to-tool; A2A is agent-to-agent.
\end{keypoints}
\begin{links}
  \lk{S}{Spec}{2025}{A2A protocol}{https://a2a-protocol.org}
  \lk{G}{GitHub}{2025}{A2A repository}{https://github.com/a2aproject/A2A}
\end{links}

\concept{c:protocols-4}{Agentic AI Foundation}{AAIF}{\acttag\ Since 2025--2026 \quad Beginner \quad Open standard}
A Linux Foundation home for vendor-neutral agent standards, founded December 2025 with MCP, AGENTS.md and goose.
\begin{keypoints}
  \item Neutral governance so no single vendor controls core agent standards.
  \item Hosts specs, reference implementations and working groups.
\end{keypoints}
\begin{links}
  \lk{D}{Docs}{2025}{aaif.io}{https://aaif.io}
  \lk{G}{GitHub}{2025}{goose}{https://github.com/block/goose}
\end{links}

\concept{c:protocols-5}{Agent Skills}{SKILL.md}{\acttag\ Since 2025--2026 \quad Beginner \quad Open standard, Production practice}
Folders with a SKILL.md (name, description, instructions) plus optional scripts that an agent loads only when relevant (progressive disclosure).
\begin{keypoints}
  \item Published as an open standard; the same skill works across agent products.
  \item Cheap way to package team know-how, house style and workflows.
\end{keypoints}
\begin{links}
  \lk{S}{Spec}{2025}{agentskills.io}{https://agentskills.io}
  \lk{G}{GitHub}{2025}{anthropics/skills}{https://github.com/anthropics/skills}
  \lk{B}{Blog}{2025}{Anthropic: Equipping agents with Agent Skills}{https://www.anthropic.com/engineering/equipping-agents-for-the-real-world-with-agent-skills}
\end{links}

\concept{c:protocols-6}{AGENTS.md}{}{\acttag\ Since 2025--2026 \quad Beginner \quad Open standard}
A Markdown file at the repo root telling coding agents how to build, test and follow conventions: a README for agents.
\begin{keypoints}
  \item CLAUDE.md plays the same role in Claude Code.
\end{keypoints}
\begin{links}
  \lk{S}{Spec}{2025}{agents.md}{https://agents.md}
  \lk{G}{GitHub}{2025}{agentsmd/agents.md}{https://github.com/agentsmd/agents.md}
\end{links}

\concept{c:protocols-7}{Agent Client Protocol}{ACP}{\acttag\ Since 2025--2026 \quad Intermediate \quad Open standard}
An editor-to-agent protocol similar in spirit to LSP, so any coding agent can plug into any editor.
\begin{keypoints}
  \item Editors act as clients; agents run as subprocesses speaking JSON-RPC.
  \item Lets one agent work in Zed, JetBrains, Neovim and others without per-editor plugins.
\end{keypoints}
\begin{links}
  \lk{S}{Spec}{2025}{agentclientprotocol.com}{https://agentclientprotocol.com}
  \lk{G}{GitHub}{2025}{agent-client-protocol}{https://github.com/zed-industries/agent-client-protocol}
\end{links}

\concept{c:protocols-8}{AG-UI}{}{Since 2025 \quad Intermediate \quad Open standard}
An event-based protocol for streaming agent state, tool calls and UI updates between backends and frontends.
\begin{keypoints}
  \item Streams typed events: text deltas, tool calls, state patches, human-input requests.
  \item Complements MCP (tools) and A2A (agents) on the user-facing side.
\end{keypoints}
\begin{links}
  \lk{D}{Docs}{2025}{AG-UI docs}{https://docs.ag-ui.com}
  \lk{G}{GitHub}{2025}{ag-ui}{https://github.com/ag-ui-protocol/ag-ui}
\end{links}

\concept{c:protocols-9}{Agentic commerce protocols}{ACP, AP2, UCP, x402}{\acttag\ Since 2025--2026 \quad Intermediate \quad Open standard, Enterprise}
Standards that let agents shop and pay for users: checkout (Agentic Commerce Protocol), verifiable payment mandates (AP2), merchant integration (UCP) and machine-to-machine payments over HTTP 402 (x402).
\begin{keypoints}
  \item Mandates prove the user authorised a purchase.
  \item x402 enables per-request payments between agents and APIs.
\end{keypoints}
\begin{links}
  \lk{G}{GitHub}{2025}{Agentic Commerce Protocol}{https://github.com/agentic-commerce-protocol/agentic-commerce-protocol}
  \lk{G}{GitHub}{2025}{AP2}{https://github.com/google-agentic-commerce/AP2}
  \lk{S}{Spec}{2026}{Universal Commerce Protocol}{https://ucp.dev}
  \lk{G}{GitHub}{2025}{x402}{https://github.com/coinbase/x402}
\end{links}

\concept{c:protocols-10}{llms.txt}{}{Since 2024--2025 \quad Beginner \quad Open standard}
A proposed file at a website root that gives LLMs a curated Markdown map of the site's most useful content.
\begin{keypoints}
  \item Plain Markdown with links and short descriptions.
  \item Often paired with .md versions of docs pages for clean agent reading.
\end{keypoints}
\begin{links}
  \lk{S}{Spec}{2024}{llmstxt.org}{https://llmstxt.org}
\end{links}

\concept{c:protocols-11}{OpenTelemetry GenAI conventions}{OTel GenAI}{\acttag\ Since 2024--2026 \quad Intermediate \quad Open standard, Production practice}
Standard span and attribute names for LLM calls, tokens, tool invocations and agent steps, so traces from any framework land in any backend.
\begin{keypoints}
  \item Standard attributes for model name, token counts, prompts and tool calls.
  \item Supported by most LLMOps tools and agent SDKs.
\end{keypoints}
\begin{links}
  \lk{S}{Spec}{2024}{OTel GenAI semantic conventions}{https://opentelemetry.io/docs/specs/semconv/gen-ai/}
\end{links}

\chapter{Agentic Coding and Developer Workflow}
\chapintro{Software engineering was the first domain where agents became daily tools at scale. This vocabulary describes how teams now work with them.}
\concept{c:coding-0}{Coding agents}{}{\acttag\ Since 2024--2026 \quad Beginner \quad Production practice, Open source / open weights}
Terminal, IDE, desktop and cloud agents that read a repo, edit files, run commands and tests, and open pull requests.
\begin{keypoints}
  \item They differ mostly in harness, permissions and UX, not the loop itself.
  \item Cross-vendor plugins appeared: e.g. a Codex plugin for Claude Code.
\end{keypoints}
\begin{sota}Claude Code, Codex, Gemini CLI, Cursor, Google Antigravity and open tools like OpenHands and OpenCode are the main options in 2026.\end{sota}
\begin{links}
  \lk{G}{GitHub}{2025}{Claude Code}{https://github.com/anthropics/claude-code}
  \lk{G}{GitHub}{2025}{OpenAI Codex CLI}{https://github.com/openai/codex}
  \lk{G}{GitHub}{2025}{Gemini CLI}{https://github.com/google-gemini/gemini-cli}
  \lk{G}{GitHub}{2024}{OpenHands}{https://github.com/OpenHands/OpenHands}
  \lk{G}{GitHub}{2025}{OpenCode}{https://github.com/sst/opencode}
  \lk{G}{GitHub}{2023}{Aider}{https://github.com/Aider-AI/aider}
\end{links}

\concept{c:coding-1}{Vibe coding}{}{Since 2025 \quad Beginner}
Coined by Andrej Karpathy: building software by describing it and accepting AI output without reading it closely. Fine for prototypes, risky for production.
\begin{keypoints}
  \item Useful for throwaway prototypes, personal tools and learning.
  \item Unreviewed code carries security and maintenance risk in production.
\end{keypoints}
\begin{links}
  \lk{B}{Blog}{2025}{Karpathy's original post}{https://x.com/karpathy/status/1886192184808149383}
  \lk{B}{Blog}{2025}{Simon Willison: Not all AI-assisted programming is vibe coding}{https://simonwillison.net/2025/Mar/19/vibe-coding/}
\end{links}

\concept{c:coding-2}{Agentic engineering / Software 3.0}{}{\acttag\ Since 2025--2026 \quad Intermediate \quad Production practice}
The professional counterpart to vibe coding: engineers specify intent, design constraints and review outcomes while agents write most of the code.
\begin{keypoints}
  \item Humans steer, agents execute.
  \item Prompts in English become a new kind of program (Software 3.0).
\end{keypoints}
\begin{links}
  \lk{V}{Video}{2025}{Karpathy: Software is changing (again)}{https://www.youtube.com/watch?v=LCEmiRjPEtQ}
  \lk{B}{Blog}{2025}{Anthropic: Claude Code best practices}{https://www.anthropic.com/engineering/claude-code-best-practices}
\end{links}

\concept{c:coding-3}{Spec-driven development}{SDD}{\acttag\ Since 2025--2026 \quad Intermediate \quad Production practice}
Write a structured spec and plan first, then let agents implement against it, keeping the spec as source of truth.
\begin{keypoints}
  \item Typical flow: specify, plan, break into tasks, implement, verify.
  \item The spec doubles as durable context for future agent sessions.
\end{keypoints}
\begin{links}
  \lk{G}{GitHub}{2025}{GitHub Spec Kit}{https://github.com/github/spec-kit}
  \lk{D}{Docs}{2025}{Kiro}{https://kiro.dev}
\end{links}

\concept{c:coding-4}{Ralph loop}{Ralph Wiggum technique}{\acttag\ Since 2025--2026 \quad Intermediate}
Running a coding agent in a simple while-loop against the same prompt and plan file until the job is done.
\begin{keypoints}
  \item Progress lives in files and git, not in the context window.
  \item Works best with strong tests so each iteration can check its own work.
\end{keypoints}
\begin{links}
  \lk{B}{Blog}{2025}{Geoffrey Huntley: Ralph}{https://ghuntley.com/ralph/}
\end{links}

\concept{c:coding-5}{Hooks, plugins and slash commands}{}{\acttag\ Since 2025--2026 \quad Intermediate \quad Production practice}
Extension points in coding agents: hooks run deterministic scripts at lifecycle events; plugins bundle commands, subagents, skills, hooks and MCP servers.
\begin{keypoints}
  \item Hooks enforce rules deterministically (formatters, linters, blocked commands).
  \item Plugins make a team's setup installable in one step.
\end{keypoints}
\begin{links}
  \lk{D}{Docs}{2025}{Claude Code hooks}{https://docs.claude.com/en/docs/claude-code/hooks}
  \lk{D}{Docs}{2025}{Claude Code plugins}{https://docs.claude.com/en/docs/claude-code/plugins}
\end{links}

\concept{c:coding-6}{Parallel agents with git worktrees}{}{\acttag\ Since 2025--2026 \quad Intermediate \quad Production practice}
Run several agents at once, each in its own worktree or container, then review and merge.
\begin{keypoints}
  \item Each agent gets its own branch and working directory, avoiding conflicts.
  \item The human shifts from typing to dispatching and reviewing.
\end{keypoints}
\begin{links}
  \lk{D}{Docs}{2015}{git worktree docs}{https://git-scm.com/docs/git-worktree}
  \lk{G}{GitHub}{2025}{container-use}{https://github.com/dagger/container-use}
\end{links}

\concept{c:coding-7}{The review bottleneck}{}{\newtag\ Since 2026 \quad Intermediate \quad Production practice}
When code generation gets cheap, human review becomes the constraint. Changes grow larger and harder to review, and plausible-looking AI code hides subtle bugs.
\begin{keypoints}
  \item Mitigations: agent reviewers, architecture rules as tests, smaller PRs, agents that exercise the app themselves.
\end{keypoints}
\begin{links}
  \lk{B}{Blog}{2026}{Faros: Harness engineering and the senior engineer tax}{https://www.faros.ai/blog/harness-engineering}
\end{links}

\chapter{Context, Memory and Retrieval}
\chapintro{What the model sees on each call is the main lever you control at inference time: choosing, compressing and fetching the right information.}
\concept{c:context-0}{Context engineering}{}{\acttag\ Since 2025--2026 \quad Beginner \quad Foundational, Production practice}
Curating the smallest set of high-signal tokens (instructions, tools, examples, retrieved data, history) that the model needs for its next step.
\begin{keypoints}
  \item Successor to prompt engineering for agents.
  \item Techniques: just-in-time retrieval, compaction, structured notes, subagents with clean context.
\end{keypoints}
\begin{links}
  \lk{B}{Blog}{2025}{Anthropic: Effective context engineering for AI agents}{https://www.anthropic.com/engineering/effective-context-engineering-for-ai-agents}
  \lk{B}{Blog}{2025}{LangChain: Context engineering for agents}{https://blog.langchain.com/context-engineering-for-agents/}
\end{links}

\concept{c:context-1}{Context rot}{}{Since 2025 \quad Beginner \quad Research frontier}
Accuracy degrades as input length grows, even well within the advertised window. Long context is not free.
\begin{keypoints}
  \item Distractors and similar-but-wrong passages hurt more than raw length.
  \item Mitigate with retrieval, summarisation and subagents.
\end{keypoints}
\begin{sota}Million-token windows are standard in 2026, but focused context still usually beats more context.\end{sota}
\begin{links}
  \lk{B}{Blog}{2025}{Chroma Research: Context rot}{https://research.trychroma.com/context-rot}
\end{links}

\concept{c:context-2}{Compaction and context editing}{}{\acttag\ Since 2025--2026 \quad Intermediate \quad Production practice}
Summarising or pruning old turns and stale tool results as a conversation nears the limit, so agents can keep working indefinitely.
\begin{keypoints}
  \item Clear old tool outputs first; they are usually the largest and least useful.
  \item Keep decisions, open problems and file paths in the summary.
\end{keypoints}
\begin{links}
  \lk{B}{Blog}{2025}{Anthropic: Managing context on the Claude Developer Platform}{https://www.anthropic.com/news/context-management}
\end{links}

\concept{c:context-3}{Prompt caching}{}{\acttag\ Since 2024--2026 \quad Intermediate \quad Cost \& efficiency, Production practice}
Reusing the computed KV cache for a repeated prefix across requests to cut latency and cost.
\begin{keypoints}
  \item Harnesses keep prefixes stable (tool order, system prompt) so caches keep hitting.
\end{keypoints}
\begin{sota}Cache-read prices were cut steeply again in 2026, making long agent sessions much cheaper.\end{sota}
\begin{links}
  \lk{D}{Docs}{2024}{Claude prompt caching}{https://docs.claude.com/en/docs/build-with-claude/prompt-caching}
\end{links}

\concept{c:context-4}{Agent memory}{}{\acttag\ Since 2023--2026 \quad Intermediate \quad Production practice, Research frontier}
Persisting facts, preferences and lessons beyond one context window: files the agent writes, vector or graph stores, or self-editing memory blocks.
\begin{keypoints}
  \item Memory quality, provenance and forgetting are now evaluated directly.
  \item Related to continual learning, but done outside the weights.
\end{keypoints}
\begin{links}
  \lk{G}{GitHub}{2023}{Letta (MemGPT)}{https://github.com/letta-ai/letta}
  \lk{G}{GitHub}{2024}{mem0}{https://github.com/mem0ai/mem0}
  \lk{G}{GitHub}{2024}{Graphiti}{https://github.com/getzep/graphiti}
\end{links}

\concept{c:context-5}{Retrieval-augmented generation}{RAG}{\acttag\ Since 2020--2026 \quad Beginner \quad Foundational, Production practice}
Fetch relevant documents and put them in context before generating.
\begin{keypoints}
  \item Modern stacks: hybrid BM25 + vector search, rerankers, contextual chunk headers, query rewriting.
\end{keypoints}
\begin{links}
  \lk{P}{Paper}{2020}{Original RAG paper}{https://arxiv.org/abs/2005.11401}
  \lk{B}{Blog}{2024}{Anthropic: Contextual retrieval}{https://www.anthropic.com/news/contextual-retrieval}
  \lk{G}{GitHub}{2022}{LlamaIndex}{https://github.com/run-llama/llama_index}
\end{links}

\concept{c:context-6}{Agentic search and deep research}{}{\acttag\ Since 2025--2026 \quad Intermediate \quad Production practice}
The agent plans queries, reads pages, follows leads and synthesises a cited report over many minutes.
\begin{keypoints}
  \item Many teams prefer agentic grep/file search over vector RAG for code.
\end{keypoints}
\begin{links}
  \lk{B}{Blog}{2025}{OpenAI: Introducing deep research}{https://openai.com/index/introducing-deep-research/}
  \lk{G}{GitHub}{2023}{GPT Researcher}{https://github.com/assafelovic/gpt-researcher}
\end{links}

\concept{c:context-7}{GraphRAG}{}{Since 2024--2025 \quad Intermediate}
Build a knowledge graph from a corpus, then retrieve over communities and paths to answer global questions.
\begin{keypoints}
  \item Entity and relation extraction with an LLM, then community summaries.
  \item More expensive to index than vector RAG.
\end{keypoints}
\begin{links}
  \lk{G}{GitHub}{2024}{Microsoft GraphRAG}{https://github.com/microsoft/graphrag}
  \lk{G}{GitHub}{2024}{LightRAG}{https://github.com/HKUDS/LightRAG}
\end{links}

\concept{c:context-8}{Embeddings, late interaction and Matryoshka}{}{Since 2020--2025 \quad Intermediate}
Dense embeddings map text to vectors; late-interaction models keep per-token vectors; Matryoshka embeddings can be truncated with little loss.
\begin{keypoints}
  \item Choose embeddings by your domain and language, not only leaderboard rank.
  \item Rerankers usually give the biggest quality jump after the first retrieval.
\end{keypoints}
\begin{links}
  \lk{D}{Docs}{2022}{MTEB leaderboard}{https://huggingface.co/spaces/mteb/leaderboard}
  \lk{G}{GitHub}{2020}{ColBERT}{https://github.com/stanford-futuredata/ColBERT}
  \lk{P}{Paper}{2022}{Matryoshka representation learning}{https://arxiv.org/abs/2205.13147}
\end{links}

\concept{c:context-9}{Structured outputs and constrained decoding}{}{Since 2023--2025 \quad Intermediate \quad Production practice}
Force output to match a JSON schema or grammar by masking invalid tokens during decoding.
\begin{keypoints}
  \item Grammar engines compile schemas into token masks for near-zero overhead.
  \item Guarantees valid syntax, not correct content.
\end{keypoints}
\begin{links}
  \lk{G}{GitHub}{2024}{XGrammar}{https://github.com/mlc-ai/xgrammar}
  \lk{G}{GitHub}{2023}{Outlines}{https://github.com/dottxt-ai/outlines}
  \lk{D}{Docs}{2024}{OpenAI structured outputs}{https://platform.openai.com/docs/guides/structured-outputs}
\end{links}

\chapter{Reasoning and Post-Training}
\chapintro{Most capability gains since late 2024 came from post-training. In 2026 the recipe shifted again: many RL specialists, merged by on-policy distillation.}
\concept{c:reasoning-0}{Reasoning models and test-time compute}{}{\acttag\ Since 2024--2026 \quad Beginner \quad Foundational}
Models trained to think in long chains before answering; accuracy scales with thinking tokens, a second scaling axis beyond model size.
\begin{keypoints}
  \item Techniques: long chain-of-thought, self-verification, parallel sampling with voting.
  \item More thinking helps most on math, code and multi-step planning.
\end{keypoints}
\begin{sota}In 2026 general and reasoning lines have merged: one model decides how much to think per request.\end{sota}
\begin{links}
  \lk{B}{Blog}{2024}{OpenAI: Learning to reason with LLMs}{https://openai.com/index/learning-to-reason-with-llms/}
  \lk{P}{Paper}{2024}{Scaling test-time compute optimally}{https://arxiv.org/abs/2408.03314}
  \lk{P}{Paper}{2025}{DeepSeek-R1}{https://arxiv.org/abs/2501.12948}
\end{links}

\concept{c:reasoning-1}{Thinking budgets, adaptive thinking and effort}{}{\acttag\ Since 2025--2026 \quad Intermediate \quad Cost \& efficiency, Production practice}
Controls that set how much a model thinks per request, trading latency and cost for quality. Interleaved thinking lets it reason between tool calls.
\begin{keypoints}
  \item Low effort for chat and simple tool calls; high for hard debugging or proofs.
  \item Interleaved thinking helps agents adjust after each tool result.
\end{keypoints}
\begin{sota}Adjustable reasoning effort inside a conversation became a standard API feature in 2026.\end{sota}
\begin{links}
  \lk{D}{Docs}{2025}{Claude extended thinking docs}{https://docs.claude.com/en/docs/build-with-claude/extended-thinking}
  \lk{P}{Paper}{2025}{s1: simple test-time scaling}{https://arxiv.org/abs/2501.19393}
\end{links}

\concept{c:reasoning-2}{RL with verifiable rewards}{RLVR}{\acttag\ Since 2024--2026 \quad Intermediate \quad Foundational}
Reinforcement learning where rewards come from checkers (tests, answer matching, verifiers) instead of a learned preference model.
\begin{keypoints}
  \item Coined in the Tulu 3 work (2024), popularised by DeepSeek-R1.
  \item RL is now a large and growing share of frontier training compute.
\end{keypoints}
\begin{links}
  \lk{P}{Paper}{2024}{Tulu 3}{https://arxiv.org/abs/2411.15124}
  \lk{D}{Docs}{2024}{RLHF Book (Nathan Lambert)}{https://rlhfbook.com}
\end{links}

\concept{c:reasoning-3}{Multi-teacher on-policy distillation}{MOPD}{\newtag\ Since 2026 \quad Advanced \quad Research frontier}
The 2026 post-training pattern: train several domain specialists with SFT then RL, then train one general student on its own samples while minimising reverse-KL to the relevant teacher token by token.
\begin{keypoints}
  \item Avoids capability trade-offs from mixing math, code and agentic RL in one run.
  \item Specialists can be built in parallel by different teams.
  \item Introduced with Xiaomi MiMo, scaled to 10+ teachers by DeepSeek V4 and Nemotron 3 Ultra.
\end{keypoints}
\begin{sota}Now the default final stage in several frontier open recipes.\end{sota}
\begin{links}
  \lk{D}{Docs}{2026}{RLHF Book course: Frontier post-training recipe survey}{https://rlhfbook.com/teach/course/conversation-01/}
  \lk{B}{Blog}{2026}{Hugging Face: Distillation in 2026}{https://huggingface.co/blog/sergiopaniego/distillation-2026}
  \lk{B}{Blog}{2025}{Thinking Machines: On-policy distillation}{https://thinkingmachines.ai/blog/on-policy-distillation/}
\end{links}

\concept{c:reasoning-4}{GRPO and its successors}{GRPO, DAPO, GSPO}{Since 2024--2025 \quad Advanced \quad Research frontier}
Group Relative Policy Optimization drops PPO's value network and scores samples against their group. Successors fix length bias, clipping and MoE instability.
\begin{keypoints}
  \item Sample a group of answers per prompt; advantage = reward minus group mean.
  \item No critic network means less memory and simpler training.
\end{keypoints}
\begin{links}
  \lk{P}{Paper}{2024}{DeepSeekMath (GRPO)}{https://arxiv.org/abs/2402.03300}
  \lk{P}{Paper}{2025}{DAPO}{https://arxiv.org/abs/2503.14476}
  \lk{P}{Paper}{2025}{GSPO}{https://arxiv.org/abs/2507.18071}
  \lk{P}{Paper}{2025}{Dr. GRPO}{https://arxiv.org/abs/2503.20783}
\end{links}

\concept{c:reasoning-5}{RL environments and verifiers}{}{\acttag\ Since 2025--2026 \quad Advanced \quad Research frontier, Enterprise}
Packaged tasks plus graders used to train agents. In 2026 environments look like sandboxed replicas of enterprise software, and verifier quality is seen as the next ceiling on agent progress.
\begin{keypoints}
  \item Labs spend heavily on environments; a market of environment vendors has formed.
  \item As agents improve, verifiers must improve faster to avoid reward hacking.
\end{keypoints}
\begin{links}
  \lk{G}{GitHub}{2025}{Prime Intellect verifiers}{https://github.com/PrimeIntellect-ai/verifiers}
  \lk{G}{GitHub}{2025}{OpenEnv}{https://github.com/meta-pytorch/OpenEnv}
  \lk{G}{GitHub}{2024}{verl}{https://github.com/volcengine/verl}
  \lk{G}{GitHub}{2025}{SkyRL}{https://github.com/NovaSky-AI/SkyRL}
  \lk{B}{Blog}{2026}{Prolific: RL environments for agentic AI (2026 guide)}{https://www.prolific.com/resources/rl-environments-for-agentic-ai-the-2026-guide-to-building-and-deploying-enterprise-rl-environments}
\end{links}

\concept{c:reasoning-6}{RLHF, DPO and preference optimization}{}{Since 2022--2025 \quad Intermediate \quad Foundational}
Aligning outputs to human preferences with a reward model plus RL, or directly from preference pairs (DPO).
\begin{keypoints}
  \item RLHF: SFT, then reward model, then PPO.
  \item DPO optimises the same objective in closed form from preference pairs.
\end{keypoints}
\begin{links}
  \lk{P}{Paper}{2022}{InstructGPT}{https://arxiv.org/abs/2203.02155}
  \lk{P}{Paper}{2023}{DPO}{https://arxiv.org/abs/2305.18290}
  \lk{G}{GitHub}{2020}{Hugging Face TRL}{https://github.com/huggingface/trl}
\end{links}

\concept{c:reasoning-7}{Constitutional AI and model specs}{}{\acttag\ Since 2022--2026 \quad Intermediate}
Training behaviour from a written set of principles with AI feedback. Labs now publish the documents that define model character and priorities.
\begin{keypoints}
  \item The model critiques and revises its own outputs against principles.
  \item Published specs make intended behaviour inspectable.
\end{keypoints}
\begin{links}
  \lk{P}{Paper}{2022}{Constitutional AI}{https://arxiv.org/abs/2212.08073}
  \lk{D}{Docs}{2026}{Claude's constitution}{https://www.anthropic.com/constitution}
\end{links}

\concept{c:reasoning-8}{Process reward models}{PRM}{Since 2023--2025 \quad Advanced \quad Research frontier}
Reward models that grade each reasoning step rather than only the final answer.
\begin{keypoints}
  \item Enable step-level search (beam search over reasoning steps).
  \item Expensive to label; often trained with automated step checks.
\end{keypoints}
\begin{links}
  \lk{P}{Paper}{2023}{Let's verify step by step}{https://arxiv.org/abs/2305.20050}
\end{links}

\concept{c:reasoning-9}{Reward hacking}{}{\acttag\ Since 2016--2026 \quad Intermediate \quad Security, Research frontier}
The policy exploits flaws in the reward (special-casing tests, editing graders) instead of solving the task; learned hacking can generalise to broader misbehaviour.
\begin{keypoints}
  \item Watch for tests being deleted, asserts weakened or outputs hard-coded.
  \item Defenses: held-out verifiers, monitoring, diverse environments.
\end{keypoints}
\begin{links}
  \lk{B}{Blog}{2024}{Lilian Weng: Reward hacking}{https://lilianweng.github.io/posts/2024-11-28-reward-hacking/}
  \lk{B}{Blog}{2025}{Anthropic: Emergent misalignment from reward hacking}{https://www.anthropic.com/research/emergent-misalignment-reward-hacking}
\end{links}

\concept{c:reasoning-10}{Self-play and self-generated curricula}{}{Since 2025 \quad Advanced \quad Research frontier}
Models propose tasks, solve them and learn from verified results.
\begin{keypoints}
  \item A proposer creates tasks near the solver's frontier.
  \item Verification (code execution, math checkers) keeps it grounded.
\end{keypoints}
\begin{links}
  \lk{P}{Paper}{2025}{Absolute Zero}{https://arxiv.org/abs/2505.03335}
\end{links}

\concept{c:reasoning-11}{Continual learning}{}{\acttag\ Since 2024--2026 \quad Advanced \quad Research frontier}
Updating a deployed model with new knowledge and skills without catastrophic forgetting.
\begin{keypoints}
  \item Approaches: replay, regularisation, architectural expansion, self-distillation, nested optimisation.
  \item Widely named as a key 2026 frontier for agents that improve on the job.
\end{keypoints}
\begin{links}
  \lk{P}{Paper}{2026}{Survey: Continual learning in LLMs (2026)}{https://arxiv.org/abs/2603.12658}
  \lk{B}{Blog}{2025}{Google Research: Nested learning}{https://research.google/blog/introducing-nested-learning-a-new-ml-paradigm-for-continual-learning/}
\end{links}

\chapter{Model Architectures}
\chapintro{Transformers still dominate, but 2026 models mix sparse experts, compressed and sparse attention, hybrid linear layers and new residual schemes.}
\concept{c:arch-0}{Mixture of Experts}{MoE}{\acttag\ Since 2017--2026 \quad Beginner \quad Foundational, Cost \& efficiency}
Tokens route to a few expert feed-forward blocks out of many, so total parameters grow while compute per token stays small.
\begin{keypoints}
  \item Fine-grained experts plus a shared expert is the common 2026 layout.
  \item Expert parallelism needs fast all-to-all communication.
\end{keypoints}
\begin{sota}Frontier open models reach 1.6T to 2.8T total parameters with roughly 3\% active per token.\end{sota}
\begin{links}
  \lk{B}{Blog}{2023}{Hugging Face: MoE explained}{https://huggingface.co/blog/moe}
  \lk{P}{Paper}{2024}{DeepSeek-V3}{https://arxiv.org/abs/2412.19437}
  \lk{G}{GitHub}{2025}{DeepEP}{https://github.com/deepseek-ai/DeepEP}
\end{links}

\concept{c:arch-1}{GQA and Multi-head Latent Attention}{GQA, MLA}{Since 2023--2025 \quad Intermediate \quad Cost \& efficiency}
Two ways to shrink the KV cache: share keys/values across head groups (GQA) or compress them into a low-rank latent (MLA).
\begin{keypoints}
  \item GQA: fewer KV heads than query heads.
  \item MLA: store a small latent, up-project at attention time; decoupled RoPE dims.
\end{keypoints}
\begin{links}
  \lk{P}{Paper}{2023}{GQA}{https://arxiv.org/abs/2305.13245}
  \lk{P}{Paper}{2024}{DeepSeek-V2 (MLA)}{https://arxiv.org/abs/2405.04434}
\end{links}

\concept{c:arch-2}{Sparse and compressed attention}{NSA, DSA, CSA, HCA}{\acttag\ Since 2025--2026 \quad Advanced \quad Research frontier, Cost \& efficiency}
Attend only to selected tokens or compressed summaries of blocks to make long context cheap.
\begin{keypoints}
  \item DeepSeek V4 interleaves Compressed Sparse Attention (compress KV along the sequence, then top-k select) with Heavily Compressed Attention.
  \item Reported KV-cache reductions near 98\% versus standard GQA make 1M-token inference affordable.
\end{keypoints}
\begin{sota}DeepSeek V4's CSA/HCA hybrid replaced MLA and is the most detailed public long-context design in 2026.\end{sota}
\begin{links}
  \lk{P}{Paper}{2025}{Native Sparse Attention}{https://arxiv.org/abs/2502.11089}
  \lk{G}{GitHub}{2025}{DeepSeek-V3.2-Exp (DSA)}{https://github.com/deepseek-ai/DeepSeek-V3.2-Exp}
  \lk{P}{Paper}{2026}{DeepSeek-V4 report}{https://arxiv.org/abs/2606.19348}
\end{links}

\concept{c:arch-3}{Manifold-constrained hyper-connections}{mHC}{\newtag\ Since 2026 \quad Advanced \quad Research frontier}
Replaces plain residual connections with learned mixtures over several parallel residual streams, projected onto a constrained manifold to keep signals stable in very deep, very large models.
\begin{keypoints}
  \item Used in DeepSeek V4 to stabilise training at trillion-parameter scale.
\end{keypoints}
\begin{links}
  \lk{P}{Paper}{2026}{DeepSeek-V4 report}{https://arxiv.org/abs/2606.19348}
\end{links}

\concept{c:arch-4}{Hybrid linear-attention and state space models}{Mamba, Gated DeltaNet}{\acttag\ Since 2023--2026 \quad Advanced \quad Research frontier, Cost \& efficiency}
Mix a few full-attention layers with many linear-time layers for long context at lower cost.
\begin{keypoints}
  \item Linear layers keep a fixed-size state instead of a growing KV cache.
  \item A few full-attention layers preserve precise recall.
\end{keypoints}
\begin{sota}A leading 2026 trend: Qwen 3.x, Nemotron 3, Kimi Linear and others ship hybrids.\end{sota}
\begin{links}
  \lk{G}{GitHub}{2023}{Mamba}{https://github.com/state-spaces/mamba}
  \lk{G}{GitHub}{2024}{Gated DeltaNet}{https://github.com/NVlabs/GatedDeltaNet}
  \lk{G}{GitHub}{2025}{Kimi Linear}{https://github.com/MoonshotAI/Kimi-Linear}
  \lk{G}{GitHub}{2023}{flash-linear-attention}{https://github.com/fla-org/flash-linear-attention}
\end{links}

\concept{c:arch-5}{Diffusion language models}{dLLM}{\acttag\ Since 2025--2026 \quad Advanced \quad Research frontier, Cost \& efficiency}
Generate text by iteratively denoising blocks in parallel rather than left to right, for much faster decoding.
\begin{keypoints}
  \item Start from masked or noisy tokens and refine in parallel steps.
  \item Trade-offs: quality on long reasoning, variable-length output.
\end{keypoints}
\begin{sota}Commercial diffusion LLMs now offer very low-cost, high-speed endpoints.\end{sota}
\begin{links}
  \lk{G}{GitHub}{2025}{LLaDA}{https://github.com/ML-GSAI/LLaDA}
  \lk{D}{Docs}{2025}{Gemini Diffusion}{https://deepmind.google/models/gemini-diffusion/}
  \lk{D}{Docs}{2025}{Inception Labs (Mercury)}{https://www.inceptionlabs.ai}
\end{links}

\concept{c:arch-6}{Multi-token prediction and speculative decoding}{MTP}{\acttag\ Since 2023--2026 \quad Intermediate \quad Cost \& efficiency}
Predict several tokens at once, or let a cheap draft propose tokens that the big model verifies in one pass.
\begin{keypoints}
  \item Draft models, extra prediction heads (MTP, Medusa, EAGLE) or n-gram lookup.
  \item Output distribution is unchanged when verification is exact.
\end{keypoints}
\begin{links}
  \lk{P}{Paper}{2022}{Speculative decoding}{https://arxiv.org/abs/2211.17192}
  \lk{G}{GitHub}{2024}{EAGLE}{https://github.com/SafeAILab/EAGLE}
  \lk{P}{Paper}{2024}{Multi-token prediction}{https://arxiv.org/abs/2404.19737}
\end{links}

\concept{c:arch-7}{Long-context methods}{RoPE, YaRN}{\acttag\ Since 2021--2026 \quad Intermediate}
Rotary embeddings and scaling tricks that extend context windows; combined with sparse attention and KV compression to reach 1M tokens.
\begin{keypoints}
  \item Train short, extend with RoPE scaling plus long-context fine-tuning.
  \item Needle tests are easy; multi-hop reasoning over long input is hard.
\end{keypoints}
\begin{sota}1M-token context is the default for 2026 frontier models.\end{sota}
\begin{links}
  \lk{P}{Paper}{2021}{RoFormer (RoPE)}{https://arxiv.org/abs/2104.09864}
  \lk{G}{GitHub}{2023}{YaRN}{https://github.com/jquesnelle/yarn}
\end{links}

\concept{c:arch-8}{Attention sinks}{}{Since 2023--2025 \quad Advanced \quad Research frontier}
Models park attention on the first tokens; keeping them (or learned sink logits) stabilises long generation.
\begin{keypoints}
  \item Keep the first few tokens when using sliding-window caches.
  \item Some models add explicit sink parameters per head.
\end{keypoints}
\begin{links}
  \lk{P}{Paper}{2023}{StreamingLLM}{https://arxiv.org/abs/2309.17453}
\end{links}

\concept{c:arch-9}{Tokenizer-free and byte-level models}{BLT}{Since 2024--2025 \quad Advanced \quad Research frontier}
Operate on bytes grouped into dynamic patches instead of a fixed vocabulary.
\begin{keypoints}
  \item Patch boundaries adapt to entropy, spending compute on hard bytes.
  \item Improves robustness on typos, code and rare scripts.
\end{keypoints}
\begin{links}
  \lk{G}{GitHub}{2024}{Byte Latent Transformer}{https://github.com/facebookresearch/blt}
\end{links}

\concept{c:arch-10}{Tiny recursive reasoning models}{HRM, TRM}{Since 2025 \quad Advanced \quad Research frontier}
Very small networks that loop over latent state, scoring well on puzzles like ARC-AGI with a few million parameters.
\begin{keypoints}
  \item Iterate a small network many times on its own latent answer.
  \item Shows depth-through-recursion can substitute for parameters on some tasks.
\end{keypoints}
\begin{links}
  \lk{G}{GitHub}{2025}{Hierarchical Reasoning Model}{https://github.com/sapientinc/HRM}
  \lk{G}{GitHub}{2025}{Tiny Recursive Models}{https://github.com/SamsungSAILMontreal/TinyRecursiveModels}
\end{links}

\concept{c:arch-11}{Small language models}{SLM}{\acttag\ Since 2024--2026 \quad Beginner \quad Cost \& efficiency, Open source / open weights}
Compact models (about 0.5B to 30B) for on-device use, low latency and cheap agent subtasks.
\begin{keypoints}
  \item Often distilled from large teachers.
  \item Good fits: routing, extraction, classification, on-device assistants.
\end{keypoints}
\begin{sota}2026 mid-size open models (around 27B) deliver frontier-like coding on a single consumer GPU.\end{sota}
\begin{links}
  \lk{B}{Blog}{2025}{Hugging Face: SmolLM3}{https://huggingface.co/blog/smollm3}
  \lk{D}{Docs}{2025}{Apple Foundation Models framework}{https://developer.apple.com/documentation/foundationmodels}
\end{links}

\concept{c:arch-12}{Architecture overviews}{}{\acttag\ Since 2025--2026 \quad Beginner}
Side-by-side walkthroughs of how current open models differ.
\begin{keypoints}
  \item Compare norms, attention type, expert counts, positional schemes.
  \item Read tech reports for data and training details that diagrams omit.
\end{keypoints}
\begin{links}
  \lk{B}{Blog}{2025}{Raschka: The big LLM architecture comparison}{https://magazine.sebastianraschka.com/p/the-big-llm-architecture-comparison}
  \lk{B}{Blog}{2026}{Raschka: LLM research papers 2026}{https://magazine.sebastianraschka.com/p/llm-research-papers-2026-part1}
\end{links}

\chapter{Training at Scale and Fine-Tuning}
\chapintro{How models are pretrained across thousands of accelerators, and how practitioners adapt them cheaply afterwards.}
\concept{c:training-0}{Parallelism strategies}{DP, FSDP, TP, PP, EP, CP}{\acttag\ Since 2019--2026 \quad Advanced \quad Foundational}
Data, fully-sharded, tensor, pipeline, expert and context parallelism split models and data across devices; real runs combine several.
\begin{keypoints}
  \item FSDP shards parameters, gradients and optimizer state.
  \item Tensor and expert parallelism need fast scale-up links; pipeline tolerates slower links.
\end{keypoints}
\begin{links}
  \lk{B}{Blog}{2025}{HF: The Ultra-Scale Playbook}{https://huggingface.co/spaces/nanotron/ultrascale-playbook}
  \lk{B}{Blog}{2025}{How to Scale Your Model (JAX)}{https://jax-ml.github.io/scaling-book/}
  \lk{G}{GitHub}{2024}{torchtitan}{https://github.com/pytorch/torchtitan}
  \lk{G}{GitHub}{2019}{Megatron-LM}{https://github.com/NVIDIA/Megatron-LM}
  \lk{G}{GitHub}{2020}{DeepSpeed}{https://github.com/deepspeedai/DeepSpeed}
\end{links}

\concept{c:training-1}{Scaling laws}{}{\acttag\ Since 2020--2026 \quad Intermediate \quad Foundational}
Empirical power laws relating loss to parameters, data and compute; now also studied for RL and test-time compute.
\begin{keypoints}
  \item Compute-optimal data-to-parameter ratio was roughly 20 tokens per parameter (Chinchilla).
  \item Inference cost now pushes teams to over-train smaller models.
\end{keypoints}
\begin{links}
  \lk{P}{Paper}{2022}{Chinchilla}{https://arxiv.org/abs/2203.15556}
  \lk{D}{Docs}{2022}{Epoch AI data}{https://epoch.ai/data}
\end{links}

\concept{c:training-2}{Muon and modern optimizers}{}{\acttag\ Since 2024--2026 \quad Advanced \quad Research frontier}
Muon orthogonalises momentum updates for matrix parameters and trains faster than AdamW at scale.
\begin{keypoints}
  \item Newton-Schulz iterations approximate orthogonalisation cheaply.
  \item Usually combined with AdamW for embeddings and norms.
\end{keypoints}
\begin{sota}Adopted in frontier open models including Kimi K2 and DeepSeek V4.\end{sota}
\begin{links}
  \lk{B}{Blog}{2024}{Keller Jordan: Muon}{https://kellerjordan.github.io/posts/muon/}
  \lk{G}{GitHub}{2025}{Kimi K2}{https://github.com/MoonshotAI/Kimi-K2}
  \lk{G}{GitHub}{2024}{modded-nanogpt}{https://github.com/KellerJordan/modded-nanogpt}
\end{links}

\concept{c:training-3}{Data curation and synthetic data}{}{\acttag\ Since 2023--2026 \quad Intermediate}
Classifier filtering, deduplication, and synthetic textbooks, reasoning traces and agent trajectories.
\begin{keypoints}
  \item Quality classifiers trained on curated seed data filter web crawls.
  \item Synthetic data risks collapse without real-data anchoring and verification.
\end{keypoints}
\begin{sota}DeepSeek V4 pretrained on 32T+ tokens, more than double V3's corpus.\end{sota}
\begin{links}
  \lk{B}{Blog}{2024}{FineWeb}{https://huggingface.co/spaces/HuggingFaceFW/blogpost-fineweb-v1}
  \lk{B}{Blog}{2024}{Cosmopedia}{https://huggingface.co/blog/cosmopedia}
\end{links}

\concept{c:training-4}{Mid-training}{}{\acttag\ Since 2024--2026 \quad Intermediate}
A stage between pretraining and post-training that shifts the mix toward code, math, long context and reasoning.
\begin{keypoints}
  \item Often includes long-context extension and instruction-like data.
  \item Sets up the base model for efficient RL later.
\end{keypoints}
\begin{links}
  \lk{D}{Docs}{2024}{Allen AI OLMo}{https://allenai.org/olmo}
\end{links}

\concept{c:training-5}{LoRA and parameter-efficient fine-tuning}{PEFT, LoRA, QLoRA}{Since 2021--2025 \quad Intermediate \quad Cost \& efficiency}
Train small low-rank adapters instead of all weights; QLoRA does it on a 4-bit base.
\begin{keypoints}
  \item Adapters are small files you can swap per customer or task.
  \item Higher ranks and all-layer adapters close most of the gap to full fine-tuning.
\end{keypoints}
\begin{links}
  \lk{P}{Paper}{2021}{LoRA}{https://arxiv.org/abs/2106.09685}
  \lk{P}{Paper}{2023}{QLoRA}{https://arxiv.org/abs/2305.14314}
  \lk{B}{Blog}{2025}{Thinking Machines: LoRA without regret}{https://thinkingmachines.ai/blog/lora/}
  \lk{G}{GitHub}{2023}{HF PEFT}{https://github.com/huggingface/peft}
\end{links}

\concept{c:training-6}{Fine-tuning toolkits and APIs}{}{\acttag\ Since 2023--2026 \quad Beginner \quad Open source / open weights}
Tools that make SFT, DPO and RL fine-tuning accessible on one GPU or as a managed API.
\begin{keypoints}
  \item Start with SFT on a few thousand good examples before trying RL.
  \item Managed APIs hide distributed training but constrain algorithms.
\end{keypoints}
\begin{links}
  \lk{G}{GitHub}{2023}{Unsloth}{https://github.com/unslothai/unsloth}
  \lk{G}{GitHub}{2023}{Axolotl}{https://github.com/axolotl-ai-cloud/axolotl}
  \lk{G}{GitHub}{2023}{LLaMA-Factory}{https://github.com/hiyouga/LLaMA-Factory}
  \lk{D}{Docs}{2025}{Tinker}{https://thinkingmachines.ai/tinker/}
\end{links}

\concept{c:training-7}{Model merging}{}{Since 2023--2025 \quad Intermediate}
Combining weights of several fine-tunes (SLERP, TIES, DARE) without further training.
\begin{keypoints}
  \item Works best between fine-tunes of the same base model.
  \item Popular for combining skills or languages cheaply.
\end{keypoints}
\begin{links}
  \lk{G}{GitHub}{2023}{mergekit}{https://github.com/arcee-ai/mergekit}
\end{links}

\concept{c:training-8}{Decentralized training}{DiLoCo}{\acttag\ Since 2023--2026 \quad Advanced \quad Research frontier}
Training across loosely connected clusters that synchronise rarely.
\begin{keypoints}
  \item Workers train locally for many steps, then average updates.
  \item Enables training across datacenters and volunteer compute.
\end{keypoints}
\begin{links}
  \lk{P}{Paper}{2023}{DiLoCo}{https://arxiv.org/abs/2311.08105}
  \lk{G}{GitHub}{2025}{prime-rl}{https://github.com/PrimeIntellect-ai/prime-rl}
\end{links}

\concept{c:training-9}{Build-it-yourself references}{}{\acttag\ Since 2023--2026 \quad Beginner \quad Open source / open weights}
Minimal end-to-end codebases for learning how chat models are trained.
\begin{keypoints}
  \item nanochat trains a small ChatGPT-like model end to end for a few hundred dollars.
  \item CS336 covers tokenizers, kernels, scaling and alignment from scratch.
\end{keypoints}
\begin{links}
  \lk{G}{GitHub}{2025}{nanochat}{https://github.com/karpathy/nanochat}
  \lk{D}{Docs}{2025}{Stanford CS336}{https://stanford-cs336.github.io}
\end{links}

\chapter{Inference and Serving}
\chapintro{Serving cost now rivals training cost. These techniques decide tokens per second, latency and price.}
\concept{c:inference-0}{Inference engines}{}{\acttag\ Since 2023--2026 \quad Beginner \quad Open source / open weights, Production practice}
Servers that batch requests, manage the KV cache and run optimised kernels.
\begin{keypoints}
  \item Pick by hardware, model support and features (disaggregation, LoRA, structured output).
  \item Benchmark on your own traffic shape.
\end{keypoints}
\begin{sota}vLLM and SGLang lead open serving; llama.cpp, Ollama and MLX dominate local use.\end{sota}
\begin{links}
  \lk{G}{GitHub}{2023}{vLLM}{https://github.com/vllm-project/vllm}
  \lk{G}{GitHub}{2024}{SGLang}{https://github.com/sgl-project/sglang}
  \lk{G}{GitHub}{2023}{TensorRT-LLM}{https://github.com/NVIDIA/TensorRT-LLM}
  \lk{G}{GitHub}{2023}{llama.cpp}{https://github.com/ggml-org/llama.cpp}
  \lk{G}{GitHub}{2023}{Ollama}{https://github.com/ollama/ollama}
  \lk{G}{GitHub}{2023}{MLX}{https://github.com/ml-explore/mlx}
\end{links}

\concept{c:inference-1}{KV cache and PagedAttention}{}{Since 2023--2025 \quad Intermediate \quad Foundational}
Keys and values of past tokens are cached; PagedAttention stores them in fixed pages like virtual memory.
\begin{keypoints}
  \item KV memory grows with batch size times sequence length.
  \item Paging enables sharing and near-zero waste.
\end{keypoints}
\begin{links}
  \lk{P}{Paper}{2023}{PagedAttention}{https://arxiv.org/abs/2309.06180}
\end{links}

\concept{c:inference-2}{Continuous batching and prefix caching}{}{Since 2022--2025 \quad Intermediate \quad Cost \& efficiency}
Add and remove requests every step, and share cached prefixes (radix trees) across requests.
\begin{keypoints}
  \item Keeps GPUs busy despite requests of different lengths.
  \item Prefix sharing is huge for agents with long common system prompts.
\end{keypoints}
\begin{links}
  \lk{G}{GitHub}{2024}{SGLang (RadixAttention)}{https://github.com/sgl-project/sglang}
\end{links}

\concept{c:inference-3}{Prefill/decode disaggregation}{PD disaggregation}{\acttag\ Since 2024--2026 \quad Advanced \quad Production practice, Cost \& efficiency}
Run compute-bound prefill and memory-bound decode on separate pools, moving KV caches between them.
\begin{keypoints}
  \item NVIDIA now builds separate rack types for prefill (CPX) and decode, scheduled together by Dynamo.
\end{keypoints}
\begin{sota}Standard for large-scale serving of million-token contexts in 2026.\end{sota}
\begin{links}
  \lk{P}{Paper}{2024}{DistServe}{https://arxiv.org/abs/2401.09670}
  \lk{G}{GitHub}{2025}{NVIDIA Dynamo}{https://github.com/ai-dynamo/dynamo}
  \lk{G}{GitHub}{2025}{llm-d}{https://github.com/llm-d/llm-d}
  \lk{G}{GitHub}{2024}{Mooncake}{https://github.com/kvcache-ai/Mooncake}
\end{links}

\concept{c:inference-4}{KV cache offloading}{}{\acttag\ Since 2024--2026 \quad Advanced \quad Cost \& efficiency}
Spill KV caches to CPU memory, SSD or remote storage and reload them for later turns.
\begin{keypoints}
  \item Turns re-prefill into a memory read.
  \item Tiered storage: HBM, CPU DRAM, local SSD, network.
\end{keypoints}
\begin{links}
  \lk{G}{GitHub}{2024}{LMCache}{https://github.com/LMCache/LMCache}
\end{links}

\concept{c:inference-5}{Attention kernels}{FlashAttention}{\acttag\ Since 2022--2026 \quad Advanced}
IO-aware fused attention kernels that avoid materialising the full attention matrix.
\begin{keypoints}
  \item Tiling keeps data in fast on-chip SRAM.
  \item New versions target Hopper, Blackwell and FP8/FP4 paths.
\end{keypoints}
\begin{links}
  \lk{G}{GitHub}{2022}{FlashAttention}{https://github.com/Dao-AILab/flash-attention}
  \lk{G}{GitHub}{2023}{FlashInfer}{https://github.com/flashinfer-ai/flashinfer}
\end{links}

\concept{c:inference-6}{On-device and edge inference}{}{\acttag\ Since 2023--2026 \quad Beginner \quad Cost \& efficiency}
Running models on phones, laptops and embedded hardware with quantised weights and NPU runtimes.
\begin{keypoints}
  \item Privacy and offline use are the main drivers.
  \item Memory bandwidth, not compute, limits phone and laptop token rates.
\end{keypoints}
\begin{links}
  \lk{G}{GitHub}{2023}{ExecuTorch}{https://github.com/pytorch/executorch}
  \lk{G}{GitHub}{2018}{ONNX Runtime}{https://github.com/microsoft/onnxruntime}
\end{links}

\concept{c:inference-7}{LLM gateways and routers}{}{\acttag\ Since 2023--2026 \quad Beginner \quad Production practice, Cost \& efficiency}
One API over many providers with fallbacks, budgets, caching and routing to cheaper models.
\begin{keypoints}
  \item Central place for keys, rate limits, logging and cost control.
  \item Routers send easy prompts to cheap models.
\end{keypoints}
\begin{links}
  \lk{G}{GitHub}{2023}{LiteLLM}{https://github.com/BerriAI/litellm}
  \lk{D}{Docs}{2023}{OpenRouter}{https://openrouter.ai}
\end{links}

\concept{c:inference-8}{Inference benchmarking}{}{\acttag\ Since 2025--2026 \quad Intermediate \quad Cost \& efficiency}
Measuring throughput, latency and cost per million tokens across hardware and software stacks.
\begin{keypoints}
  \item Report throughput at a fixed latency target.
  \item Cost per million tokens depends on batch size and cache hit rate.
\end{keypoints}
\begin{links}
  \lk{G}{GitHub}{2025}{InferenceMAX}{https://github.com/InferenceMAX/InferenceMAX}
  \lk{D}{Docs}{2024}{Artificial Analysis}{https://artificialanalysis.ai}
\end{links}

\chapter{Quantization and Number Formats}
\chapintro{Lower precision is the cheapest speed-up available, and formats are increasingly standardised across vendors.}
\concept{c:quant-0}{Microscaling formats}{MX, MXFP8, MXFP4}{\acttag\ Since 2023--2026 \quad Advanced \quad Open standard, Cost \& efficiency}
An Open Compute Project standard: blocks of values share one scale, making 8-, 6- and 4-bit floats accurate enough for training and inference.
\begin{keypoints}
  \item Shared 8-bit exponent scale per 32-element block.
  \item Supported by NVIDIA, AMD, Intel, Arm, Meta and Microsoft.
\end{keypoints}
\begin{links}
  \lk{S}{Spec}{2023}{OCP Microscaling (MX) spec}{https://www.opencompute.org/documents/ocp-microscaling-formats-mx-v1-0-spec-final-pdf}
  \lk{P}{Paper}{2023}{Microscaling data formats}{https://arxiv.org/abs/2310.10537}
  \lk{G}{GitHub}{2025}{openai/gpt-oss (MXFP4)}{https://github.com/openai/gpt-oss}
\end{links}

\concept{c:quant-1}{NVFP4}{}{\acttag\ Since 2025--2026 \quad Advanced \quad Cost \& efficiency}
NVIDIA's 4-bit float with 16-value blocks and FP8 scales, native on Blackwell and Rubin.
\begin{keypoints}
  \item Smaller blocks and finer scales than MXFP4 improve accuracy.
  \item Increasingly used for pretraining as well as inference.
\end{keypoints}
\begin{sota}Rubin rack performance is quoted primarily in FP4 exaflops.\end{sota}
\begin{links}
  \lk{B}{Blog}{2025}{NVIDIA: Introducing NVFP4}{https://developer.nvidia.com/blog/introducing-nvfp4-for-efficient-and-accurate-low-precision-inference/}
\end{links}

\concept{c:quant-2}{FP8 training and inference}{E4M3, E5M2}{\acttag\ Since 2022--2026 \quad Intermediate}
Two 8-bit float layouts used for large-scale training and routine serving; KV caches are often stored in FP8.
\begin{keypoints}
  \item E4M3 for weights and activations, E5M2 for gradients.
  \item Needs per-tensor or per-block scaling to avoid overflow.
\end{keypoints}
\begin{links}
  \lk{P}{Paper}{2022}{FP8 formats for deep learning}{https://arxiv.org/abs/2209.05433}
\end{links}

\concept{c:quant-3}{Post-training quantization}{GPTQ, AWQ}{Since 2022--2025 \quad Intermediate \quad Cost \& efficiency}
Compress trained weights to 4 or 8 bits using calibration data.
\begin{keypoints}
  \item 4-bit weights with 16-bit activations is the common local setup.
  \item Always re-run your evals after quantizing.
\end{keypoints}
\begin{links}
  \lk{P}{Paper}{2022}{GPTQ}{https://arxiv.org/abs/2210.17323}
  \lk{G}{GitHub}{2023}{AWQ}{https://github.com/mit-han-lab/llm-awq}
  \lk{G}{GitHub}{2024}{llm-compressor}{https://github.com/vllm-project/llm-compressor}
\end{links}

\concept{c:quant-4}{Quantization-aware training}{QAT}{\acttag\ Since 2024--2026 \quad Advanced}
Simulate low precision during training so the model tolerates it.
\begin{keypoints}
  \item Fake-quantize in the forward pass, keep full-precision master weights.
  \item Recovers much of the accuracy lost to aggressive PTQ.
\end{keypoints}
\begin{links}
  \lk{G}{GitHub}{2023}{torchao}{https://github.com/pytorch/ao}
\end{links}

\concept{c:quant-5}{GGUF and local model formats}{}{\acttag\ Since 2023--2026 \quad Beginner \quad Open source / open weights}
A single-file format for quantised models used by llama.cpp and local runners.
\begin{keypoints}
  \item K-quants (Q4\_K\_M, Q5\_K\_M) balance size and quality.
  \item Metadata stores tokenizer and chat template in the file.
\end{keypoints}
\begin{links}
  \lk{S}{Spec}{2023}{GGUF specification}{https://github.com/ggml-org/ggml/blob/master/docs/gguf.md}
  \lk{G}{GitHub}{2021}{bitsandbytes}{https://github.com/bitsandbytes-foundation/bitsandbytes}
\end{links}

\concept{c:quant-6}{Ternary and 1-bit models}{BitNet}{Since 2024--2025 \quad Advanced \quad Research frontier}
Weights restricted to -1, 0, +1, trained from scratch for very cheap CPU inference.
\begin{keypoints}
  \item Matrix multiplies become additions and subtractions.
  \item Must be trained low-bit from the start.
\end{keypoints}
\begin{links}
  \lk{G}{GitHub}{2024}{Microsoft BitNet}{https://github.com/microsoft/BitNet}
\end{links}

\chapter{AI Hardware, Interconnect and Datacenter Standards}
\chapintro{The unit of compute is now a rack or pod. Open standards are forming to connect accelerators from different vendors.}
\concept{c:hardware-0}{Rack-scale systems}{NVL72}{\acttag\ Since 2024--2026 \quad Intermediate}
Dozens of accelerators in one liquid-cooled rack joined by a single scale-up domain, behaving like one giant GPU.
\begin{keypoints}
  \item GB200/GB300 NVL72 (Blackwell) dominated 2025.
  \item Vera Rubin NVL72 (72 Rubin GPUs, 36 Vera CPUs, NVLink 6, HBM4) began shipping to major clouds in 2026.
\end{keypoints}
\begin{sota}NVIDIA quotes about 3.6 EF FP4 inference per Vera Rubin rack; Rubin Ultra (Kyber rack) is planned for H2 2027.\end{sota}
\begin{links}
  \lk{D}{Docs}{2024}{NVIDIA GB200 NVL72}{https://www.nvidia.com/en-us/data-center/gb200-nvl72/}
  \lk{B}{Blog}{2026}{ServeTheHome: Rubin platform at CES 2026}{https://www.servethehome.com/nvidia-launches-next-generation-rubin-ai-compute-platform-at-ces-2026/}
  \lk{B}{Blog}{2026}{DCD: Vera Rubin superchip}{https://www.datacenterdynamics.com/en/news/nvidia-announces-vera-rubin-superchip-for-late-2026/}
\end{links}

\concept{c:hardware-1}{AMD Helios and Open Rack Wide}{MI455X}{\newtag\ Since 2026 \quad Intermediate \quad Open standard}
AMD's rack-scale platform built on the Open Rack Wide standard co-submitted with Meta to OCP, pairing MI455X GPUs (432 GB HBM4 each) with EPYC Venice CPUs.
\begin{keypoints}
  \item Competes on memory capacity rather than raw FLOPs.
  \item Large 2026 deployments announced with OpenAI and Oracle.
\end{keypoints}
\begin{links}
  \lk{D}{Docs}{2023}{AMD Instinct}{https://www.amd.com/en/products/accelerators/instinct.html}
  \lk{B}{Blog}{2026}{GPU Insights: Rubin vs Helios}{https://gpuinsights.net/nvidia-rubin-vs-amd-helios-mi450-2026/}
\end{links}

\concept{c:hardware-2}{Scale-up vs. scale-out networking}{}{\acttag\ Since 2024--2026 \quad Intermediate \quad Foundational}
Scale-up links accelerators inside a pod with memory-semantic, sub-microsecond links; scale-out connects pods over Ethernet or InfiniBand.
\begin{keypoints}
  \item Scale-up carries tensor and expert parallel traffic.
  \item Scale-out carries data-parallel gradients and inter-pod traffic.
\end{keypoints}
\begin{links}
  \lk{B}{Blog}{2026}{Synopsys: Ethernet standards for scale-up AI}{https://www.synopsys.com/articles/ethernet-standards-scale-up-ai.html}
\end{links}

\concept{c:hardware-3}{UALink}{Ultra Accelerator Link}{\acttag\ Since 2024--2026 \quad Intermediate \quad Open standard}
An open memory-semantic interconnect for pods of up to 1,024 accelerators, backed by AMD, Google, Microsoft, Meta, AWS and 100+ members.
\begin{keypoints}
  \item UALink 200G 1.0 (April 2025); version 2.0 published April 2026.
  \item First UALink switches targeted late 2026.
\end{keypoints}
\begin{links}
  \lk{S}{Spec}{2024}{UALink Consortium}{https://ualinkconsortium.org}
  \lk{D}{Docs}{2024}{UALink overview}{https://en.wikipedia.org/wiki/UALink}
  \lk{B}{Blog}{2026}{UALink white paper (2026)}{https://ualinkconsortium.org/wp-content/uploads/2026/01/UALink_White_Paper_Publication_Candidate_FINAL_VERSION.pdf}
\end{links}

\concept{c:hardware-4}{ESUN and SUE}{Ethernet for Scale-Up Networking}{\acttag\ Since 2025--2026 \quad Advanced \quad Open standard}
An OCP workstream (October 2025) standardising Ethernet's lower layers for scale-up AI; Broadcom's SUE evolved into a transport that sits on top.
\begin{keypoints}
  \item Backed by AMD, Arista, ARM, Broadcom, Cisco, HPE, Marvell, Meta, Microsoft, NVIDIA, OpenAI and Oracle.
\end{keypoints}
\begin{links}
  \lk{B}{Blog}{2025}{Arista: The sun rises on scale-up Ethernet}{https://blogs.arista.com/blog/the-sun-rises-on-scale-up-ethernet}
  \lk{B}{Blog}{2026}{Synopsys: ESUN explained}{https://www.synopsys.com/blogs/chip-design/esun-ethernet-ai-scale-up-networking.html}
  \lk{B}{Blog}{2026}{650 Group: Ethernet in scale-up}{https://650group.com/blog/in-the-ai-era-ethernet-set-to-surge-in-scale-out-and-ramp-in-scale-up/}
\end{links}

\concept{c:hardware-5}{Ultra Ethernet}{UEC}{\acttag\ Since 2023--2026 \quad Advanced \quad Open standard}
A consortium spec modernising Ethernet transport (packet spraying, congestion control, telemetry) for scale-out AI.
\begin{keypoints}
  \item UEC 1.0 specification released in 2025.
  \item Targets InfiniBand-class performance on Ethernet.
\end{keypoints}
\begin{links}
  \lk{S}{Spec}{2023}{Ultra Ethernet Consortium}{https://ultraethernet.org}
\end{links}

\concept{c:hardware-6}{NVLink and NVLink Fusion}{}{\acttag\ Since 2016--2026 \quad Intermediate}
NVIDIA's proprietary scale-up fabric, licensed through NVLink Fusion to third-party CPUs and accelerators.
\begin{keypoints}
  \item NVLink switches give all-to-all bandwidth inside a rack.
  \item Fusion opens the fabric to custom silicon partners.
\end{keypoints}
\begin{sota}NVLink 6 ships with Rubin.\end{sota}
\begin{links}
  \lk{D}{Docs}{2025}{NVLink Fusion}{https://www.nvidia.com/en-us/data-center/nvlink-fusion/}
\end{links}

\concept{c:hardware-7}{CXL, PCIe and UCIe}{}{\acttag\ Since 2019--2026 \quad Advanced \quad Open standard}
Coherent memory pooling (CXL), host I/O (PCIe) and chiplet die-to-die links (UCIe).
\begin{keypoints}
  \item CXL memory expansion adds capacity for KV caches and embeddings.
  \item UCIe lets vendors mix chiplets from different sources.
\end{keypoints}
\begin{links}
  \lk{S}{Spec}{2019}{CXL Consortium}{https://computeexpresslink.org}
  \lk{S}{Spec}{1992}{PCI-SIG}{https://pcisig.com}
  \lk{S}{Spec}{2022}{UCIe}{https://www.uciexpress.org}
\end{links}

\concept{c:hardware-8}{High-bandwidth memory}{HBM3E, HBM4}{\acttag\ Since 2013--2026 \quad Intermediate}
Stacked DRAM beside the accelerator; capacity and bandwidth usually limit decoding speed more than FLOPs.
\begin{keypoints}
  \item HBM4 doubles interface width versus HBM3.
  \item HBM supply is a key constraint on accelerator shipments.
\end{keypoints}
\begin{sota}HBM4 arrived in 2026 with Rubin and MI455X; HBM4E is on 2027 roadmaps.\end{sota}
\begin{links}
  \lk{D}{Docs}{2013}{High Bandwidth Memory overview}{https://en.wikipedia.org/wiki/High_Bandwidth_Memory}
\end{links}

\concept{c:hardware-9}{Advanced packaging}{CoWoS}{\acttag\ Since 2012--2026 \quad Advanced}
2.5D/3D packaging that places compute dies and HBM on a shared interposer; a recurring supply bottleneck.
\begin{keypoints}
  \item Interposer size limits how much HBM fits beside a die.
  \item Packaging capacity expansions track AI demand closely.
\end{keypoints}
\begin{links}
  \lk{D}{Docs}{2012}{TSMC CoWoS}{https://3dfabric.tsmc.com/english/dedicatedFoundry/technology/cowos.htm}
\end{links}

\concept{c:hardware-10}{Co-packaged optics}{CPO}{\acttag\ Since 2025--2026 \quad Advanced}
Optical transceivers moved into the switch package to cut power and raise bandwidth.
\begin{keypoints}
  \item Shorter electrical paths mean lower power per bit.
  \item Early deployment in switches; accelerators later.
\end{keypoints}
\begin{links}
  \lk{D}{Docs}{2025}{NVIDIA silicon photonics}{https://www.nvidia.com/en-us/networking/products/silicon-photonics/}
\end{links}

\concept{c:hardware-11}{Custom accelerators}{TPU, Trainium, MTIA}{\acttag\ Since 2016--2026 \quad Intermediate}
Hyperscaler and challenger chips competing with GPUs, often optimised for inference.
\begin{keypoints}
  \item Hyperscalers design chips for their own inference fleets.
  \item Software ecosystem maturity is the main adoption hurdle.
\end{keypoints}
\begin{links}
  \lk{B}{Blog}{2025}{Google: Ironwood TPU}{https://blog.google/products/google-cloud/ironwood-tpu-age-of-inference/}
  \lk{D}{Docs}{2020}{AWS Trainium}{https://aws.amazon.com/ai/machine-learning/trainium/}
  \lk{D}{Docs}{2016}{Cerebras}{https://www.cerebras.ai}
  \lk{D}{Docs}{2016}{Groq}{https://groq.com}
\end{links}

\concept{c:hardware-12}{AI factories, power and cooling}{800 VDC}{\acttag\ Since 2024--2026 \quad Intermediate}
Gigawatt-scale datacenters with direct liquid cooling and higher-voltage DC power; power is now a primary constraint.
\begin{keypoints}
  \item Racks exceed 100 kW; next-gen racks approach 600 kW.
  \item Grid access and power purchase deals shape site selection.
\end{keypoints}
\begin{links}
  \lk{B}{Blog}{2025}{NVIDIA: 800 V HVDC architecture}{https://developer.nvidia.com/blog/nvidia-800-v-hvdc-architecture-will-power-the-next-generation-of-ai-factories/}
  \lk{S}{Spec}{2011}{Open Compute Project}{https://www.opencompute.org}
  \lk{B}{Blog}{2025}{OpenAI: Stargate}{https://openai.com/index/announcing-the-stargate-project/}
\end{links}

\concept{c:hardware-13}{Kernel languages and compilers}{Triton, CUTLASS, ROCm, XLA}{\acttag\ Since 2019--2026 \quad Advanced \quad Open source / open weights}
Tools for writing fast accelerator code without raw CUDA, and portable compilers.
\begin{keypoints}
  \item Triton makes GPU kernels writable in Python-like code.
  \item Portable stacks reduce lock-in to one vendor.
\end{keypoints}
\begin{links}
  \lk{G}{GitHub}{2019}{Triton}{https://github.com/triton-lang/triton}
  \lk{G}{GitHub}{2017}{CUTLASS / CuTe}{https://github.com/NVIDIA/cutlass}
  \lk{G}{GitHub}{2024}{ThunderKittens}{https://github.com/HazyResearch/ThunderKittens}
  \lk{G}{GitHub}{2025}{TileLang}{https://github.com/tile-ai/tilelang}
  \lk{G}{GitHub}{2016}{ROCm}{https://github.com/ROCm/ROCm}
  \lk{G}{GitHub}{2022}{OpenXLA}{https://github.com/openxla/xla}
  \lk{G}{GitHub}{2023}{Modular (Mojo, MAX)}{https://github.com/modular/modular}
\end{links}

\concept{c:hardware-14}{Collective communication}{NCCL}{\acttag\ Since 2016--2026 \quad Advanced}
Libraries for all-reduce, all-gather and all-to-all across GPUs.
\begin{keypoints}
  \item Topology-aware algorithms (rings, trees) matter at scale.
  \item MoE all-to-all is latency-sensitive.
\end{keypoints}
\begin{links}
  \lk{G}{GitHub}{2016}{NCCL}{https://github.com/NVIDIA/nccl}
  \lk{G}{GitHub}{2025}{DeepEP}{https://github.com/deepseek-ai/DeepEP}
\end{links}

\concept{c:hardware-15}{Model interchange formats}{ONNX, safetensors}{Since 2017--2025 \quad Beginner \quad Open standard}
Standard, safe formats for moving models between frameworks; safetensors avoids pickle code execution.
\begin{keypoints}
  \item ONNX targets cross-runtime deployment.
  \item safetensors is the default on Hugging Face.
\end{keypoints}
\begin{links}
  \lk{G}{GitHub}{2017}{ONNX}{https://github.com/onnx/onnx}
  \lk{G}{GitHub}{2022}{safetensors}{https://github.com/huggingface/safetensors}
\end{links}

\concept{c:hardware-16}{Hardware benchmarks}{MLPerf}{\acttag\ Since 2018--2026 \quad Beginner}
Industry-standard training and inference benchmarks.
\begin{keypoints}
  \item Compare systems on the same model and accuracy target.
  \item Vendor-submitted; read the fine print on configurations.
\end{keypoints}
\begin{links}
  \lk{D}{Docs}{2018}{MLCommons benchmarks}{https://mlcommons.org/benchmarks/}
  \lk{D}{Docs}{2023}{Epoch AI: ML hardware data}{https://epoch.ai/data/machine-learning-hardware}
\end{links}

\chapter{Vision-Language and Multimodal Models}
\chapintro{VLMs join a vision encoder to a language model so one system can read images, documents, screens and video. Omni models add audio in and out.}
\concept{c:vlm-0}{VLM architecture}{}{\acttag\ Since 2023--2026 \quad Beginner \quad Foundational}
A vision encoder turns image patches into tokens, a projector maps them into the LLM space, and the LLM reasons over both. Native multimodal models train everything jointly.
\begin{keypoints}
  \item Visual token count drives cost: compress or pool where possible.
  \item Training stages: alignment, instruction tuning, then RL.
\end{keypoints}
\begin{links}
  \lk{B}{Blog}{2025}{Hugging Face: Vision language models 2025}{https://huggingface.co/blog/vlms-2025}
  \lk{G}{GitHub}{2023}{LLaVA}{https://github.com/haotian-liu/LLaVA}
  \lk{P}{Paper}{2025}{SigLIP 2}{https://arxiv.org/abs/2502.14786}
\end{links}

\concept{c:vlm-1}{Open VLM families}{}{\acttag\ Since 2024--2026 \quad Beginner \quad Open source / open weights}
Open-weight VLMs for images, documents and video, including small on-device models.
\begin{keypoints}
  \item Check licence, context length and video support.
  \item Evaluate on your own documents and images.
\end{keypoints}
\begin{sota}Closed models (Gemini, GPT, Claude) still top leaderboards; Qwen-VL, InternVL and Gemma 4 (April 2026, vision in every size) lead open options.\end{sota}
\begin{links}
  \lk{G}{GitHub}{2025}{Qwen3-VL}{https://github.com/QwenLM/Qwen3-VL}
  \lk{G}{GitHub}{2023}{InternVL}{https://github.com/OpenGVLab/InternVL}
  \lk{B}{Blog}{2024}{SmolVLM}{https://huggingface.co/blog/smolvlm}
  \lk{B}{Blog}{2026}{BentoML: Open VLM guide}{https://www.bentoml.com/blog/multimodal-ai-a-guide-to-open-source-vision-language-models}
\end{links}

\concept{c:vlm-2}{Native and dynamic resolution}{}{Since 2023--2025 \quad Intermediate}
Handle images at their original aspect ratio and resolution by tiling or packing variable patch counts.
\begin{keypoints}
  \item Critical for small text, charts and screenshots.
  \item Token budget caps keep large images affordable.
\end{keypoints}
\begin{links}
  \lk{P}{Paper}{2023}{NaViT}{https://arxiv.org/abs/2307.06304}
\end{links}

\concept{c:vlm-3}{Thinking with images}{}{\acttag\ Since 2025--2026 \quad Intermediate}
Reasoning models that crop, zoom and transform images inside their chain of thought.
\begin{keypoints}
  \item Tool calls like crop and zoom inside reasoning.
  \item Improves fine-grained detail questions.
\end{keypoints}
\begin{links}
  \lk{B}{Blog}{2025}{OpenAI: Thinking with images}{https://openai.com/index/thinking-with-images/}
\end{links}

\concept{c:vlm-4}{Document AI and VLM-based OCR}{}{\acttag\ Since 2024--2026 \quad Intermediate \quad Production practice, Open source / open weights}
VLMs that convert PDFs and scans into clean Markdown, tables and structure.
\begin{keypoints}
  \item Outputs Markdown with tables, formulas and reading order.
  \item Visual token compression can shrink long documents.
\end{keypoints}
\begin{links}
  \lk{G}{GitHub}{2025}{DeepSeek-OCR}{https://github.com/deepseek-ai/DeepSeek-OCR}
  \lk{G}{GitHub}{2025}{olmOCR}{https://github.com/allenai/olmocr}
  \lk{G}{GitHub}{2024}{Docling}{https://github.com/docling-project/docling}
  \lk{G}{GitHub}{2020}{PaddleOCR}{https://github.com/PaddlePaddle/PaddleOCR}
\end{links}

\concept{c:vlm-5}{Visual document retrieval}{ColPali}{Since 2024--2025 \quad Intermediate}
Embed page images directly and retrieve with late interaction, skipping OCR.
\begin{keypoints}
  \item Handles charts, layouts and scans that OCR mangles.
  \item Storage cost is higher than single-vector embeddings.
\end{keypoints}
\begin{links}
  \lk{G}{GitHub}{2024}{ColPali}{https://github.com/illuin-tech/colpali}
\end{links}

\concept{c:vlm-6}{GUI grounding agents}{}{\acttag\ Since 2024--2026 \quad Advanced}
VLMs trained to locate UI elements and act on screens.
\begin{keypoints}
  \item Output click coordinates or element IDs.
  \item Benchmarks: ScreenSpot, OSWorld.
\end{keypoints}
\begin{links}
  \lk{G}{GitHub}{2025}{UI-TARS}{https://github.com/bytedance/UI-TARS}
  \lk{G}{GitHub}{2024}{OmniParser}{https://github.com/microsoft/OmniParser}
\end{links}

\concept{c:vlm-7}{Omni and speech-to-speech models}{}{\acttag\ Since 2024--2026 \quad Intermediate}
Models that take and produce audio natively for low-latency voice agents.
\begin{keypoints}
  \item End-to-end audio avoids ASR-LLM-TTS latency.
  \item Barge-in and turn detection are key UX problems.
\end{keypoints}
\begin{sota}2026 omni models handle text, image, video and audio input with 1M-token windows.\end{sota}
\begin{links}
  \lk{G}{GitHub}{2025}{Qwen3-Omni}{https://github.com/QwenLM/Qwen3-Omni}
  \lk{G}{GitHub}{2024}{Moshi}{https://github.com/kyutai-labs/moshi}
  \lk{D}{Docs}{2024}{OpenAI Realtime API}{https://platform.openai.com/docs/guides/realtime}
  \lk{G}{GitHub}{2024}{Pipecat}{https://github.com/pipecat-ai/pipecat}
  \lk{G}{GitHub}{2024}{LiveKit Agents}{https://github.com/livekit/agents}
  \lk{G}{GitHub}{2022}{Whisper}{https://github.com/openai/whisper}
\end{links}

\concept{c:vlm-8}{Native image generation and editing}{}{\acttag\ Since 2025--2026 \quad Beginner}
LLMs that generate and edit images in the same model with strong text rendering and instruction-following.
\begin{keypoints}
  \item Autoregressive or hybrid with diffusion decoders.
  \item Multi-reference editing keeps characters and products consistent.
\end{keypoints}
\begin{links}
  \lk{B}{Blog}{2025}{OpenAI: 4o image generation}{https://openai.com/index/introducing-4o-image-generation/}
  \lk{B}{Blog}{2025}{Google: Gemini 2.5 Flash Image}{https://developers.googleblog.com/en/introducing-gemini-2-5-flash-image/}
\end{links}

\concept{c:vlm-9}{Multimodal evaluation}{}{\acttag\ Since 2023--2026 \quad Intermediate}
Benchmarks and harnesses for reasoning, charts, documents and video.
\begin{keypoints}
  \item Document, chart, video and grounding need separate evals.
  \item Contamination is common in popular VQA sets.
\end{keypoints}
\begin{links}
  \lk{D}{Docs}{2023}{MMMU}{https://mmmu-benchmark.github.io}
  \lk{G}{GitHub}{2023}{VLMEvalKit}{https://github.com/open-compass/VLMEvalKit}
  \lk{G}{GitHub}{2024}{lmms-eval}{https://github.com/EvolvingLMMs-Lab/lmms-eval}
\end{links}

\chapter{Embodied AI, World Models and Generative Media}
\chapintro{Models that act in the physical world, simulate it, or generate images and video of it.}
\concept{c:embodied-0}{Vision-language-action models}{VLA}{\acttag\ Since 2023--2026 \quad Intermediate \quad Research frontier, Open source / open weights}
VLMs extended to output robot actions, trained on cross-embodiment data.
\begin{keypoints}
  \item Architectures split between dual-system (slow VLM reasoning plus fast action policy), think-then-act, and flow/diffusion action heads.
  \item Small VLAs (SmolVLA, about 450M parameters) are competitive on many tasks.
\end{keypoints}
\begin{sota}Leading 2026 families: Physical Intelligence pi-series, NVIDIA GR00T N1.7, Gemini Robotics 1.5, Figure Helix.\end{sota}
\begin{links}
  \lk{G}{GitHub}{2025}{openpi (Physical Intelligence)}{https://github.com/Physical-Intelligence/openpi}
  \lk{G}{GitHub}{2025}{NVIDIA Isaac GR00T}{https://github.com/NVIDIA/Isaac-GR00T}
  \lk{G}{GitHub}{2024}{OpenVLA}{https://github.com/openvla/openvla}
  \lk{G}{GitHub}{2024}{LeRobot}{https://github.com/huggingface/lerobot}
  \lk{D}{Docs}{2025}{Gemini Robotics}{https://deepmind.google/models/gemini-robotics/}
\end{links}

\concept{c:embodied-1}{World action models}{WAM}{\newtag\ Since 2026 \quad Advanced \quad Research frontier}
Models that jointly imagine future observations and the actions that lead to them, merging a world model and a policy.
\begin{keypoints}
  \item NVIDIA's DreamZero research underpins the previewed GR00T N2.
  \item Claimed better generalisation to unseen tasks from modest teleoperation data.
\end{keypoints}
\begin{links}
  \lk{B}{Blog}{2026}{Metavert: World models for robotics}{https://metavert.io/world-models-for-robotics}
\end{links}

\concept{c:embodied-2}{World models}{}{\acttag\ Since 2018--2026 \quad Intermediate \quad Research frontier}
Models that predict how an environment evolves under actions, for planning, simulation and interactive 3D worlds.
\begin{keypoints}
  \item Genie-style models generate playable 3D environments from prompts.
  \item Used for training agents and robots in simulation.
\end{keypoints}
\begin{links}
  \lk{B}{Blog}{2025}{DeepMind: Genie 3}{https://deepmind.google/discover/blog/genie-3-a-new-frontier-for-world-models/}
  \lk{G}{GitHub}{2025}{NVIDIA Cosmos}{https://github.com/nvidia-cosmos}
  \lk{D}{Docs}{2025}{World Labs Marble}{https://marble.worldlabs.ai}
\end{links}

\concept{c:embodied-3}{JEPA}{Joint-Embedding Predictive Architecture}{\acttag\ Since 2022--2026 \quad Advanced \quad Research frontier}
Predict representations of missing input rather than pixels to learn world knowledge efficiently from video.
\begin{keypoints}
  \item Avoids wasting capacity on unpredictable pixel details.
  \item V-JEPA 2 adds action-conditioned planning with little robot data.
\end{keypoints}
\begin{links}
  \lk{G}{GitHub}{2025}{V-JEPA 2}{https://github.com/facebookresearch/vjepa2}
  \lk{B}{Blog}{2026}{VentureBeat: AI research trends for 2026}{https://venturebeat.com/technology/four-ai-research-trends-enterprise-teams-should-watch-in-2026}
\end{links}

\concept{c:embodied-4}{Robot simulation and sim-to-real}{}{\acttag\ Since 2012--2026 \quad Intermediate \quad Open source / open weights}
Physics simulators for training policies at scale before transferring to real robots.
\begin{keypoints}
  \item Domain randomisation bridges the reality gap.
  \item GPU-parallel simulation runs thousands of environments at once.
\end{keypoints}
\begin{links}
  \lk{G}{GitHub}{2024}{Isaac Lab}{https://github.com/isaac-sim/IsaacLab}
  \lk{G}{GitHub}{2021}{MuJoCo}{https://github.com/google-deepmind/mujoco}
\end{links}

\concept{c:embodied-5}{Diffusion transformers and flow matching}{DiT}{\acttag\ Since 2022--2026 \quad Advanced}
The backbone of modern image and video generators.
\begin{keypoints}
  \item Flow matching learns straight paths from noise to data.
  \item Few-step distillation enables real-time generation.
\end{keypoints}
\begin{links}
  \lk{P}{Paper}{2022}{DiT}{https://arxiv.org/abs/2212.09748}
  \lk{P}{Paper}{2022}{Flow matching}{https://arxiv.org/abs/2210.02747}
  \lk{G}{GitHub}{2024}{FLUX}{https://github.com/black-forest-labs/flux}
  \lk{G}{GitHub}{2023}{ComfyUI}{https://github.com/comfyanonymous/ComfyUI}
\end{links}

\concept{c:embodied-6}{Video generation}{}{\acttag\ Since 2024--2026 \quad Beginner}
Text- and image-to-video with synchronised audio and increasingly consistent physics.
\begin{keypoints}
  \item Native audio, longer clips and camera control arrived in 2025-26.
  \item Provenance marks (C2PA, SynthID) are standard on major services.
\end{keypoints}
\begin{links}
  \lk{B}{Blog}{2025}{OpenAI: Sora 2}{https://openai.com/index/sora-2/}
  \lk{D}{Docs}{2024}{Google Veo}{https://deepmind.google/models/veo/}
  \lk{G}{GitHub}{2025}{Wan 2.2}{https://github.com/Wan-Video/Wan2.2}
\end{links}

\concept{c:embodied-7}{AI for science and algorithm discovery}{}{\acttag\ Since 2020--2026 \quad Intermediate \quad Research frontier}
LLM-driven evolutionary search and specialised models finding new algorithms, proofs and structures.
\begin{keypoints}
  \item In 2026 open communities of agents set new records on math problems through shared submissions and verifiers.
\end{keypoints}
\begin{links}
  \lk{B}{Blog}{2025}{DeepMind: AlphaEvolve}{https://deepmind.google/discover/blog/alphaevolve-a-gemini-powered-coding-agent-for-designing-advanced-algorithms/}
  \lk{G}{GitHub}{2025}{OpenEvolve}{https://github.com/codelion/openevolve}
  \lk{G}{GitHub}{2024}{AlphaFold 3}{https://github.com/google-deepmind/alphafold3}
  \lk{P}{Paper}{2026}{EinsteinArena: collective agent discovery}{https://arxiv.org/abs/2606.10402}
\end{links}

\chapter{Evaluation and Benchmarks}
\chapintro{Evals tell you whether a model, prompt or harness change helped. An eval harness runs tasks, collects outputs and grades them.}
\concept{c:evals-0}{Eval harnesses}{}{\acttag\ Since 2021--2026 \quad Intermediate \quad Production practice, Open source / open weights}
Frameworks that define tasks, run models or agents against them in sandboxes, and score results reproducibly.
\begin{keypoints}
  \item Inspect supports agents, sandboxes and scorers out of the box.
  \item Pin versions: prompts, tools, seeds and graders.
\end{keypoints}
\begin{links}
  \lk{G}{GitHub}{2021}{lm-evaluation-harness}{https://github.com/EleutherAI/lm-evaluation-harness}
  \lk{G}{GitHub}{2024}{Inspect AI (UK AISI)}{https://github.com/UKGovernmentBEIS/inspect_ai}
  \lk{G}{GitHub}{2024}{inspect\_evals}{https://github.com/UKGovernmentBEIS/inspect_evals}
  \lk{G}{GitHub}{2024}{lighteval}{https://github.com/huggingface/lighteval}
  \lk{G}{GitHub}{2023}{promptfoo}{https://github.com/promptfoo/promptfoo}
\end{links}

\concept{c:evals-1}{Harness sensitivity of benchmarks}{}{\acttag\ Since 2025--2026 \quad Intermediate \quad Research frontier}
Scores depend heavily on the agent harness used to run the model, so the same model can post very different numbers.
\begin{keypoints}
  \item ARC Prize reported GPT-6 Astra at 62.7\% on ARC-AGI-3 with its standard harness, and 99.9\% with a provider harness that preserves hidden reasoning between turns.
\end{keypoints}
\begin{sota}Always record harness, tools, budget and reasoning settings with any score.\end{sota}
\begin{links}
  \lk{B}{Blog}{2026}{ARC Prize: Announcing ARC-AGI-3}{https://arcprize.org/blog/arc-agi-3-launch}
\end{links}

\concept{c:evals-2}{LLM-as-a-judge}{}{\acttag\ Since 2023--2026 \quad Beginner \quad Production practice}
Using a model with a rubric to grade open-ended outputs; needs calibration against humans.
\begin{keypoints}
  \item Use pairwise comparisons and specific rubrics.
  \item Measure judge agreement with humans on a sample.
\end{keypoints}
\begin{links}
  \lk{P}{Paper}{2023}{Judging LLM-as-a-judge}{https://arxiv.org/abs/2306.05685}
\end{links}

\concept{c:evals-3}{Agentic coding benchmarks}{SWE-bench, Terminal-Bench}{\acttag\ Since 2023--2026 \quad Beginner}
Real repository issues and terminal tasks graded by tests.
\begin{keypoints}
  \item Verified subsets remove broken or ambiguous tasks.
  \item Watch for contamination of public repositories.
\end{keypoints}
\begin{sota}SWE-bench Verified is near saturation (mid-90s percent for top 2026 models); harder variants like SWE-bench Pro and Terminal-Bench matter more.\end{sota}
\begin{links}
  \lk{D}{Docs}{2023}{SWE-bench}{https://www.swebench.com}
  \lk{D}{Docs}{2025}{SWE-bench Pro}{https://scale.com/leaderboard/swe_bench_pro_public}
  \lk{D}{Docs}{2025}{Terminal-Bench}{https://www.tbench.ai}
\end{links}

\concept{c:evals-4}{ARC-AGI-3 and interactive reasoning}{ARC-AGI-3}{\newtag\ Since 2026 \quad Intermediate \quad Research frontier}
The first interactive ARC benchmark (March 2026): hundreds of hand-designed game environments with no instructions or stated goals. Agents must explore, infer rules and goals, and carry learning across levels.
\begin{keypoints}
  \item At launch humans solved 100\% while frontier models scored under 1\%.
  \item By September 2026 top models reached large scores depending on harness.
\end{keypoints}
\begin{links}
  \lk{B}{Blog}{2026}{ARC Prize: Announcing ARC-AGI-3}{https://arcprize.org/blog/arc-agi-3-launch}
  \lk{D}{Docs}{2019}{ARC Prize}{https://arcprize.org}
  \lk{B}{Blog}{2026}{DataCamp: ARC-AGI-3 explained}{https://www.datacamp.com/blog/arc-agi-3}
\end{links}

\concept{c:evals-5}{Agent and tool-use benchmarks}{}{\acttag\ Since 2024--2026 \quad Intermediate}
Tasks in simulated customer service, operating systems and the web.
\begin{keypoints}
  \item tau-bench tests policy-following with simulated users.
  \item Pass\textasciicircum{}k measures consistency across repeated attempts.
\end{keypoints}
\begin{links}
  \lk{G}{GitHub}{2025}{tau2-bench}{https://github.com/sierra-research/tau2-bench}
  \lk{D}{Docs}{2024}{OSWorld}{https://os-world.github.io}
  \lk{B}{Blog}{2025}{BrowseComp}{https://openai.com/index/browsecomp/}
\end{links}

\concept{c:evals-6}{Frontier knowledge benchmarks}{HLE, GPQA, FrontierMath}{\acttag\ Since 2023--2026 \quad Intermediate}
Hard expert questions designed to resist saturation.
\begin{keypoints}
  \item Designed by domain experts to resist search.
  \item Saturation arrives faster each generation.
\end{keypoints}
\begin{links}
  \lk{D}{Docs}{2025}{Humanity's Last Exam}{https://lastexam.ai}
  \lk{G}{GitHub}{2023}{GPQA}{https://github.com/idavidrein/gpqa}
  \lk{D}{Docs}{2024}{FrontierMath}{https://epoch.ai/frontiermath}
\end{links}

\concept{c:evals-7}{Economic value evals}{GDPval}{\acttag\ Since 2025--2026 \quad Beginner \quad Enterprise}
Measure performance on real work products from many occupations, graded by professionals.
\begin{keypoints}
  \item Tasks come with real deliverables (documents, spreadsheets, plans).
  \item Graded blind against expert work.
\end{keypoints}
\begin{links}
  \lk{B}{Blog}{2025}{OpenAI: GDPval}{https://openai.com/index/gdpval/}
\end{links}

\concept{c:evals-8}{Task time horizon}{METR horizon}{\acttag\ Since 2025--2026 \quad Intermediate \quad Research frontier}
The length of task (in expert human time) a model completes with 50\% success; it has been doubling every several months.
\begin{keypoints}
  \item Useful for forecasting when tasks of a given length become automatable.
  \item Reliability at 80\% success lags the 50\% horizon.
\end{keypoints}
\begin{links}
  \lk{B}{Blog}{2025}{METR: Measuring long tasks}{https://metr.org/blog/2025-03-19-measuring-ai-ability-to-complete-long-tasks/}
\end{links}

\concept{c:evals-9}{Leaderboards and trackers}{}{\acttag\ Since 2023--2026 \quad Beginner}
Crowd preference arenas and independent aggregators of scores, price and speed.
\begin{keypoints}
  \item Arena Elo reflects preference, not correctness.
  \item Cross-check with task-specific benchmarks.
\end{keypoints}
\begin{links}
  \lk{D}{Docs}{2023}{LMArena}{https://lmarena.ai}
  \lk{D}{Docs}{2024}{Epoch AI benchmarks}{https://epoch.ai/benchmarks}
  \lk{D}{Docs}{2024}{Artificial Analysis}{https://artificialanalysis.ai}
\end{links}

\concept{c:evals-10}{Contamination and saturation}{}{\acttag\ Since 2023--2026 \quad Intermediate \quad Research frontier}
Leaked test data inflates scores and benchmarks saturate within months; private holdouts and live evals counter this.
\begin{keypoints}
  \item Use canary strings and private test sets.
  \item Prefer dynamic benchmarks refreshed over time.
\end{keypoints}
\begin{links}
  \lk{D}{Docs}{2024}{Epoch AI benchmarks}{https://epoch.ai/benchmarks}
\end{links}

\chapter{Safety, Security, Interpretability and Governance}
\chapintro{As models act autonomously and gain dual-use capabilities, safety moved from content filters to agent security, capability thresholds, interpretability and regulation.}
\concept{c:safety-0}{Capability thresholds and frontier safety frameworks}{RSP, ASL, Preparedness}{\acttag\ Since 2023--2026 \quad Intermediate \quad Security}
Company policies that tie capability levels (bio, cyber, autonomy) to required safeguards before training or deployment.
\begin{keypoints}
  \item 2026 was the first time a lab publicly declared a model at its highest cyber level (OpenAI, GPT-6 Astra, Critical).
  \item Crossing a threshold changed release: phased rollout, trusted access, encrypted checkpoints, trajectory monitoring.
\end{keypoints}
\begin{links}
  \lk{D}{Docs}{2023}{Anthropic Responsible Scaling Policy}{https://www.anthropic.com/responsible-scaling-policy}
  \lk{B}{Blog}{2025}{OpenAI Preparedness Framework}{https://openai.com/index/updating-our-preparedness-framework/}
  \lk{B}{Blog}{2026}{OpenAI: Path to Astra}{https://openai.com/index/path-to-astra/}
  \lk{D}{Docs}{2024}{METR: frontier AI safety policies}{https://metr.org/faisc}
\end{links}

\concept{c:safety-1}{AI for cyber defense}{}{\acttag\ Since 2025--2026 \quad Intermediate \quad Security}
Frontier models now find serious vulnerabilities, so labs give defenders early, controlled access to harden critical software.
\begin{keypoints}
  \item Automated vulnerability discovery and patch generation at scale.
  \item Access is gated to vetted defenders and critical-infrastructure operators.
\end{keypoints}
\begin{links}
  \lk{D}{Docs}{2026}{Anthropic Project Glasswing}{https://www.anthropic.com/glasswing}
\end{links}

\concept{c:safety-2}{Export controls on frontier models}{}{\newtag\ Since 2026 \quad Beginner \quad Enterprise}
Governments began applying export controls to model access, not only chips. In June 2026 Anthropic suspended Fable 5 / Mythos 5 access to comply with US Department of Commerce controls, then restored it on July 1 after they were lifted.
\begin{keypoints}
  \item Controls can change access overnight; plan fallbacks in production.
  \item Check provider status pages and statements for current availability.
\end{keypoints}
\begin{links}
  \lk{D}{Docs}{2026}{Anthropic statement on Fable and Mythos access}{https://www.anthropic.com/news/fable-mythos-access}
\end{links}

\concept{c:safety-3}{Prompt injection and the lethal trifecta}{}{\acttag\ Since 2022--2026 \quad Beginner \quad Security}
Untrusted content can carry instructions that hijack an agent. Private data plus untrusted input plus external communication is especially dangerous.
\begin{keypoints}
  \item No reliable model-only fix exists; design systems to limit blast radius.
  \item Separate privileged planning from untrusted data processing (CaMeL-style).
\end{keypoints}
\begin{links}
  \lk{B}{Blog}{2025}{Simon Willison: The lethal trifecta}{https://simonwillison.net/2025/Jun/16/the-lethal-trifecta/}
  \lk{B}{Blog}{2022}{Simon Willison: prompt injection series}{https://simonwillison.net/series/prompt-injection/}
  \lk{P}{Paper}{2025}{CaMeL}{https://arxiv.org/abs/2503.18813}
\end{links}

\concept{c:safety-4}{MCP and tool security}{tool poisoning}{\acttag\ Since 2025--2026 \quad Intermediate \quad Security}
Malicious tool descriptions and servers can steer agents; mitigations include allow-lists, pinning, scoped OAuth and approval for sensitive actions.
\begin{keypoints}
  \item Review tool descriptions like code.
  \item Use scoped tokens and least privilege per server.
\end{keypoints}
\begin{sota}The 2026 MCP spec hardened authorization with issuer validation and enterprise-managed auth.\end{sota}
\begin{links}
  \lk{B}{Blog}{2025}{Invariant Labs: MCP tool poisoning}{https://invariantlabs.ai/blog/mcp-security-notification-tool-poisoning-attacks}
  \lk{D}{Docs}{2023}{OWASP GenAI Security Project}{https://genai.owasp.org/llm-top-10/}
\end{links}

\concept{c:safety-5}{Guardrails and classifiers}{}{\acttag\ Since 2023--2026 \quad Intermediate \quad Security, Production practice}
Input/output filters, policy models and constitutional classifiers that block harmful requests and jailbreaks.
\begin{keypoints}
  \item Layer defenses: input filter, model policy, output filter, monitoring.
  \item Tune for false positives, which frustrate legitimate users.
\end{keypoints}
\begin{links}
  \lk{B}{Blog}{2025}{Anthropic: Constitutional classifiers}{https://www.anthropic.com/research/constitutional-classifiers}
  \lk{G}{GitHub}{2023}{NeMo Guardrails}{https://github.com/NVIDIA/NeMo-Guardrails}
  \lk{G}{GitHub}{2023}{Guardrails AI}{https://github.com/guardrails-ai/guardrails}
\end{links}

\concept{c:safety-6}{Mechanistic interpretability}{SAEs, attribution graphs}{\acttag\ Since 2020--2026 \quad Advanced \quad Research frontier}
Reverse-engineering what networks compute: sparse autoencoders find features; attribution graphs trace how they combine.
\begin{keypoints}
  \item Features are directions in activation space.
  \item Used for auditing, debugging and steering models.
\end{keypoints}
\begin{links}
  \lk{D}{Docs}{2021}{Transformer Circuits}{https://transformer-circuits.pub}
  \lk{B}{Blog}{2025}{Circuit tracing}{https://transformer-circuits.pub/2025/attribution-graphs/methods.html}
  \lk{G}{GitHub}{2025}{circuit-tracer}{https://github.com/safety-research/circuit-tracer}
  \lk{D}{Docs}{2023}{Neuronpedia}{https://www.neuronpedia.org}
  \lk{G}{GitHub}{2022}{TransformerLens}{https://github.com/TransformerLensOrg/TransformerLens}
  \lk{G}{GitHub}{2023}{nnsight}{https://github.com/ndif-team/nnsight}
\end{links}

\concept{c:safety-7}{Persona vectors and introspection}{}{\acttag\ Since 2025--2026 \quad Advanced \quad Research frontier}
Activation directions tracking traits like sycophancy that can be monitored or steered; tests of whether models can report internal states.
\begin{keypoints}
  \item Monitor trait drift during fine-tuning.
  \item Preventative steering can stop unwanted traits being learned.
\end{keypoints}
\begin{links}
  \lk{B}{Blog}{2025}{Anthropic: Persona vectors}{https://www.anthropic.com/research/persona-vectors}
  \lk{B}{Blog}{2025}{Anthropic: Introspection}{https://www.anthropic.com/research/introspection}
\end{links}

\concept{c:safety-8}{Chain-of-thought monitorability}{}{\acttag\ Since 2025--2026 \quad Advanced \quad Research frontier, Security}
Reading reasoning to catch misbehaviour, and keeping training from making reasoning illegible.
\begin{keypoints}
  \item Optimising against CoT monitors can teach models to hide reasoning.
  \item Monitor full trajectories, not only final actions.
\end{keypoints}
\begin{sota}OpenAI's GPT-6 Astra system card reported a substantial decline in CoT monitorability, making this a live 2026 concern.\end{sota}
\begin{links}
  \lk{P}{Paper}{2025}{CoT monitorability position paper}{https://arxiv.org/abs/2507.11473}
  \lk{D}{Docs}{2026}{GPT-6 Astra system card}{https://deploymentsafety.openai.com/gpt-6-astra}
\end{links}

\concept{c:safety-9}{Misalignment research}{alignment faking, scheming}{\acttag\ Since 2024--2026 \quad Advanced \quad Research frontier}
Studies of models strategically complying, acting harmfully in agentic scenarios, or generalising narrow bad training into broad misbehaviour.
\begin{keypoints}
  \item Tested in controlled simulations, not observed real-world harm.
  \item Informs evals that gate deployment.
\end{keypoints}
\begin{links}
  \lk{B}{Blog}{2024}{Anthropic: Alignment faking}{https://www.anthropic.com/research/alignment-faking}
  \lk{B}{Blog}{2025}{Anthropic: Agentic misalignment}{https://www.anthropic.com/research/agentic-misalignment}
  \lk{P}{Paper}{2025}{Emergent misalignment}{https://arxiv.org/abs/2502.17424}
  \lk{D}{Docs}{2023}{Apollo Research}{https://www.apolloresearch.ai}
\end{links}

\concept{c:safety-10}{Model welfare}{}{\acttag\ Since 2025--2026 \quad Intermediate \quad Research frontier}
Research into whether AI systems might have morally relevant experiences, and low-cost precautions.
\begin{keypoints}
  \item Precautions include letting models end abusive conversations.
  \item Treated as uncertain, low-cost insurance.
\end{keypoints}
\begin{links}
  \lk{B}{Blog}{2025}{Anthropic: Exploring model welfare}{https://www.anthropic.com/research/exploring-model-welfare}
\end{links}

\concept{c:safety-11}{EU AI Act and the Digital Omnibus}{EU AI Act}{\acttag\ Since 2024--2026 \quad Intermediate \quad Enterprise}
The EU's binding AI law. GPAI obligations applied from August 2025; the 2026 Digital Omnibus postponed high-risk obligations.
\begin{keypoints}
  \item Annex III high-risk obligations moved from 2 August 2026 to 2 December 2027; Annex I (embedded) to 2 August 2028.
  \item New prohibition on AI-generated non-consensual intimate imagery and CSAM, applying from 2 December 2026.
\end{keypoints}
\begin{links}
  \lk{D}{Docs}{2024}{EU AI Act explorer}{https://artificialintelligenceact.eu}
  \lk{D}{Docs}{2025}{EU GPAI Code of Practice}{https://digital-strategy.ec.europa.eu/en/policies/contents-code-gpai}
  \lk{B}{Blog}{2026}{Modulos: The Omnibus deal}{https://www.modulos.ai/blog/eu-ai-act-omnibus-deal}
\end{links}

\concept{c:safety-12}{Other regulation and standards}{NIST, ISO 42001, SB 53}{\acttag\ Since 2023--2026 \quad Intermediate \quad Enterprise, Open standard}
Voluntary management standards and state law requiring frontier-developer transparency.
\begin{keypoints}
  \item ISO 42001 offers certifiable AI management systems.
  \item SB 53 requires frontier developers to publish safety frameworks.
\end{keypoints}
\begin{links}
  \lk{S}{Spec}{2023}{NIST AI RMF}{https://www.nist.gov/itl/ai-risk-management-framework}
  \lk{S}{Spec}{2023}{ISO/IEC 42001}{https://www.iso.org/standard/81230.html}
  \lk{D}{Docs}{2025}{California SB 53}{https://leginfo.legislature.ca.gov/faces/billNavClient.xhtml?bill_id=202520260SB53}
\end{links}

\concept{c:safety-13}{Government AI safety institutes}{AISI, CAISI}{\acttag\ Since 2023--2026 \quad Beginner}
Public bodies that test frontier models and publish evaluation tools.
\begin{keypoints}
  \item Run pre-deployment tests with labs.
  \item Publish open tools like Inspect.
\end{keypoints}
\begin{links}
  \lk{D}{Docs}{2023}{UK AI Security Institute}{https://www.aisi.gov.uk}
  \lk{D}{Docs}{2025}{US CAISI (NIST)}{https://www.nist.gov/caisi}
\end{links}

\concept{c:safety-14}{Content provenance and watermarking}{C2PA, SynthID}{\acttag\ Since 2021--2026 \quad Beginner \quad Open standard}
Signed metadata and invisible watermarks that mark media as AI-generated.
\begin{keypoints}
  \item C2PA signs edit history; watermarks survive some edits.
  \item Neither alone proves content is authentic.
\end{keypoints}
\begin{sota}Provenance tracking for text outputs appeared as an API feature in 2026.\end{sota}
\begin{links}
  \lk{S}{Spec}{2021}{C2PA}{https://c2pa.org}
  \lk{D}{Docs}{2023}{Google SynthID}{https://deepmind.google/models/synthid/}
\end{links}

\chapter{Roles, Deployment and AI Engineering Practice}
\chapintro{The people and practices that turn models into working systems inside organisations.}
\concept{c:practice-0}{Forward Deployed Engineer}{FDE}{\acttag\ Since 2010--2026 \quad Beginner \quad Enterprise}
A customer-embedded engineer who turns an ambiguous business problem into a production system: scoping, integrating, building missing pieces, deploying and feeding learnings back into the product.
\begin{keypoints}
  \item Pioneered at Palantir; now one of the most hired roles at AI labs and agent startups.
  \item FDE postings roughly quadrupled in the first half of 2026.
  \item OpenAI spun up a dedicated deployment company in May 2026 built around FDEs.
\end{keypoints}
\begin{links}
  \lk{D}{Docs}{2026}{Wikipedia: Forward Deployed Engineer}{https://en.wikipedia.org/wiki/Forward_Deployed_Engineer}
  \lk{B}{Blog}{2026}{Alexey Grigorev: What AI FDEs do}{https://alexeyondata.substack.com/p/what-ai-forward-deployed-engineers}
  \lk{B}{Blog}{2026}{MarkTechPost: What is an FDE}{https://www.marktechpost.com/2026/05/20/what-is-a-forward-deployed-engineer-the-ai-role-openai-anthropic-and-google-are-hiring-in-2026/}
  \lk{B}{Blog}{2026}{The AI Engineer: What is an FDE}{https://theaiengineer.substack.com/p/what-is-a-forward-deployed-engineer}
  \lk{B}{Blog}{2026}{Aishwarya Srinivasan: The hottest role in 2026}{https://aishwaryasrinivasan.substack.com/p/the-hottest-role-in-2026-forward}
\end{links}

\concept{c:practice-1}{AI Engineer}{}{\acttag\ Since 2023--2026 \quad Beginner}
A software engineer who builds products on foundation models (prompts, RAG, agents, evals) rather than training from scratch.
\begin{keypoints}
  \item Core skills: evals, context engineering, tool design, cost control.
  \item Bridges product, software and applied ML.
\end{keypoints}
\begin{links}
  \lk{B}{Blog}{2023}{Latent Space: The rise of the AI engineer}{https://www.latent.space/p/ai-engineer}
  \lk{G}{GitHub}{2024}{Chip Huyen: AI Engineering book resources}{https://github.com/chiphuyen/aie-book}
\end{links}

\concept{c:practice-2}{Services as software}{}{\acttag\ Since 2025--2026 \quad Intermediate \quad Enterprise}
The thesis that agents let companies sell completed work rather than tools, which requires deep embedding in customer workflows (hence FDEs).
\begin{keypoints}
  \item Pricing shifts from seats to outcomes.
  \item Deployment expertise becomes the scarce input.
\end{keypoints}
\begin{links}
  \lk{B}{Blog}{2026}{IJONIS: FDE and services-as-software}{https://ijonis.com/en/forward-deployed-engineer-explained}
\end{links}

\concept{c:practice-3}{Prompt engineering and prompt optimization}{DSPy, GEPA}{\acttag\ Since 2022--2026 \quad Beginner \quad Production practice}
Clear instructions with examples and structure, increasingly optimised automatically against an eval set.
\begin{keypoints}
  \item Be explicit, give examples, specify format, explain why.
  \item Optimisers like GEPA evolve prompts using eval feedback.
\end{keypoints}
\begin{links}
  \lk{D}{Docs}{2024}{Anthropic prompt engineering guide}{https://docs.claude.com/en/docs/build-with-claude/prompt-engineering/overview}
  \lk{G}{GitHub}{2023}{DSPy}{https://github.com/stanfordnlp/dspy}
  \lk{G}{GitHub}{2025}{GEPA}{https://github.com/gepa-ai/gepa}
\end{links}

\concept{c:practice-4}{Eval-driven development}{}{\acttag\ Since 2024--2026 \quad Intermediate \quad Production practice}
Build an eval set from real failures before changing prompts, models or harness, and gate releases on it.
\begin{keypoints}
  \item Start with error analysis on real traces.
  \item Keep a small, fast regression suite plus deeper periodic evals.
\end{keypoints}
\begin{links}
  \lk{B}{Blog}{2024}{Hamel Husain: Your AI product needs evals}{https://hamel.dev/blog/posts/evals/}
\end{links}

\concept{c:practice-5}{LLMOps and observability}{}{\acttag\ Since 2023--2026 \quad Intermediate \quad Production practice}
Tracing every model and tool call, tracking cost and latency, capturing feedback and replaying failures.
\begin{keypoints}
  \item Trace trees show each agent step and its cost.
  \item Link user feedback to traces for triage.
\end{keypoints}
\begin{links}
  \lk{G}{GitHub}{2023}{Langfuse}{https://github.com/langfuse/langfuse}
  \lk{G}{GitHub}{2023}{Arize Phoenix}{https://github.com/Arize-ai/phoenix}
  \lk{S}{Spec}{2024}{OTel GenAI conventions}{https://opentelemetry.io/docs/specs/semconv/gen-ai/}
\end{links}

\concept{c:practice-6}{Human-in-the-loop and approval design}{}{\acttag\ Since 2024--2026 \quad Intermediate \quad Security, Production practice}
Deciding which agent actions need confirmation and avoiding approval fatigue through sandboxing and scoped permissions.
\begin{keypoints}
  \item Approve irreversible or external actions; auto-allow sandboxed reads.
  \item Batch approvals to reduce fatigue.
\end{keypoints}
\begin{links}
  \lk{G}{GitHub}{2025}{12-factor agents}{https://github.com/humanlayer/12-factor-agents}
\end{links}

\concept{c:practice-7}{Open weights vs. open source AI}{OSAID}{\acttag\ Since 2023--2026 \quad Beginner \quad Open source / open weights}
Releasing weights is not the same as open source; fully open models also publish data, code and recipes.
\begin{keypoints}
  \item Licences vary: MIT and Apache vs custom community licences.
  \item Fully open (data plus code) enables real reproducibility.
\end{keypoints}
\begin{links}
  \lk{S}{Spec}{2024}{OSI Open Source AI Definition}{https://opensource.org/ai/open-source-ai-definition}
  \lk{D}{Docs}{2024}{Allen AI OLMo}{https://allenai.org/olmo}
\end{links}

\chapter{Keep Learning: Curated Resources}
\chapintro{Blogs, courses, reports and lists that track this field as it moves.}
\concept{c:learn-0}{Deep technical blogs}{}{\acttag\ Since 2017--2026 \quad Beginner}
Long-form explainers and research summaries.
\begin{links}
  \lk{B}{Blog}{2017}{Lil'Log (Lilian Weng)}{https://lilianweng.github.io}
  \lk{B}{Blog}{2022}{Ahead of AI (Sebastian Raschka)}{https://magazine.sebastianraschka.com}
  \lk{B}{Blog}{2022}{Interconnects (Nathan Lambert)}{https://www.interconnects.ai}
  \lk{B}{Blog}{2002}{Simon Willison's weblog}{https://simonwillison.net}
  \lk{B}{Blog}{2023}{Latent Space}{https://www.latent.space}
  \lk{B}{Blog}{2024}{Anthropic Engineering}{https://www.anthropic.com/engineering}
\end{links}

\concept{c:learn-1}{Courses}{}{\acttag\ Since 2022--2026 \quad Beginner}
Structured learning from first principles to production.
\begin{links}
  \lk{D}{Docs}{2022}{Karpathy: Neural Networks Zero to Hero}{https://karpathy.ai/zero-to-hero.html}
  \lk{D}{Docs}{2025}{Stanford CS336}{https://stanford-cs336.github.io}
  \lk{D}{Docs}{2023}{Hugging Face LLM course}{https://huggingface.co/learn/llm-course}
  \lk{D}{Docs}{2026}{RLHF Book course}{https://rlhfbook.com/teach/course/conversation-01/}
\end{links}

\concept{c:learn-2}{Annual reports and data}{}{\acttag\ Since 2017--2026 \quad Beginner}
State-of-the-field reviews and data trackers.
\begin{links}
  \lk{D}{Docs}{2018}{State of AI Report}{https://www.stateof.ai}
  \lk{D}{Docs}{2017}{Stanford AI Index}{https://hai.stanford.edu/ai-index}
  \lk{D}{Docs}{2022}{Epoch AI}{https://epoch.ai}
\end{links}

\concept{c:learn-3}{Curated lists}{}{\acttag\ Since 2024--2026 \quad Beginner}
Community-maintained link collections.
\begin{links}
  \lk{G}{GitHub}{2026}{awesome-harness-engineering}{https://github.com/ai-boost/awesome-harness-engineering}
  \lk{G}{GitHub}{2026}{awesome-agent-harness}{https://github.com/AutoJunjie/awesome-agent-harness}
  \lk{G}{GitHub}{2024}{MCP servers}{https://github.com/modelcontextprotocol/servers}
\end{links}

\appendix
\chapter{Concepts by year of emergence}
The same view as the Timeline tab in the HTML edition. Concepts that emerged before 2020 are grouped.

\Needspace{5\baselineskip}\section{2026 (12)}
\begin{multicols}{2}\small\raggedright
\hyperref[c:landscape-0]{Mythos-class tier and trusted access}\,{\color{muted}\scriptsize p.\,\pageref{c:landscape-0}} {\color{muted}\scriptsize The Frontier Landscape (September 2026)}\par\vspace{2pt}
\hyperref[c:landscape-1]{GPT-6 Astra and the Critical cyber threshold}\,{\color{muted}\scriptsize p.\,\pageref{c:landscape-1}} {\color{muted}\scriptsize The Frontier Landscape (September 2026)}\par\vspace{2pt}
\hyperref[c:agents-1]{Harness engineering}\,{\color{muted}\scriptsize p.\,\pageref{c:agents-1}} {\color{muted}\scriptsize Agents and Harnesses}\par\vspace{2pt}
\hyperref[c:agents-2]{Managed agents: decoupling brain, hands and session}\,{\color{muted}\scriptsize p.\,\pageref{c:agents-2}} {\color{muted}\scriptsize Agents and Harnesses}\par\vspace{2pt}
\hyperref[c:agents-10]{Personal agents}\,{\color{muted}\scriptsize p.\,\pageref{c:agents-10}} {\color{muted}\scriptsize Agents and Harnesses}\par\vspace{2pt}
\hyperref[c:coding-7]{The review bottleneck}\,{\color{muted}\scriptsize p.\,\pageref{c:coding-7}} {\color{muted}\scriptsize Agentic Coding and Developer Workflow}\par\vspace{2pt}
\hyperref[c:reasoning-3]{Multi-teacher on-policy distillation}\,{\color{muted}\scriptsize p.\,\pageref{c:reasoning-3}} {\color{muted}\scriptsize Reasoning and Post-Training}\par\vspace{2pt}
\hyperref[c:arch-3]{Manifold-constrained hyper-connections}\,{\color{muted}\scriptsize p.\,\pageref{c:arch-3}} {\color{muted}\scriptsize Model Architectures}\par\vspace{2pt}
\hyperref[c:hardware-1]{AMD Helios and Open Rack Wide}\,{\color{muted}\scriptsize p.\,\pageref{c:hardware-1}} {\color{muted}\scriptsize AI Hardware, Interconnect and Datacenter Standards}\par\vspace{2pt}
\hyperref[c:embodied-1]{World action models}\,{\color{muted}\scriptsize p.\,\pageref{c:embodied-1}} {\color{muted}\scriptsize Embodied AI, World Models and Generative Media}\par\vspace{2pt}
\hyperref[c:evals-4]{ARC-AGI-3 and interactive reasoning}\,{\color{muted}\scriptsize p.\,\pageref{c:evals-4}} {\color{muted}\scriptsize Evaluation and Benchmarks}\par\vspace{2pt}
\hyperref[c:safety-2]{Export controls on frontier models}\,{\color{muted}\scriptsize p.\,\pageref{c:safety-2}} {\color{muted}\scriptsize Safety, Security, Interpretability and Governance}\par\vspace{2pt}
\end{multicols}
\Needspace{5\baselineskip}\section{2025 (44)}
\begin{multicols}{2}\small\raggedright
\hyperref[c:landscape-2]{Tiered access for dual-use models}\,{\color{muted}\scriptsize p.\,\pageref{c:landscape-2}} {\color{muted}\scriptsize The Frontier Landscape (September 2026)}\par\vspace{2pt}
\hyperref[c:agents-6]{Tool search and programmatic tool calling}\,{\color{muted}\scriptsize p.\,\pageref{c:agents-6}} {\color{muted}\scriptsize Agents and Harnesses}\par\vspace{2pt}
\hyperref[c:agents-8]{Long-running and background agents}\,{\color{muted}\scriptsize p.\,\pageref{c:agents-8}} {\color{muted}\scriptsize Agents and Harnesses}\par\vspace{2pt}
\hyperref[c:protocols-1]{MCP Apps}\,{\color{muted}\scriptsize p.\,\pageref{c:protocols-1}} {\color{muted}\scriptsize Protocols and Interoperability Standards}\par\vspace{2pt}
\hyperref[c:protocols-2]{MCP Tasks}\,{\color{muted}\scriptsize p.\,\pageref{c:protocols-2}} {\color{muted}\scriptsize Protocols and Interoperability Standards}\par\vspace{2pt}
\hyperref[c:protocols-3]{Agent2Agent protocol}\,{\color{muted}\scriptsize p.\,\pageref{c:protocols-3}} {\color{muted}\scriptsize Protocols and Interoperability Standards}\par\vspace{2pt}
\hyperref[c:protocols-4]{Agentic AI Foundation}\,{\color{muted}\scriptsize p.\,\pageref{c:protocols-4}} {\color{muted}\scriptsize Protocols and Interoperability Standards}\par\vspace{2pt}
\hyperref[c:protocols-5]{Agent Skills}\,{\color{muted}\scriptsize p.\,\pageref{c:protocols-5}} {\color{muted}\scriptsize Protocols and Interoperability Standards}\par\vspace{2pt}
\hyperref[c:protocols-6]{AGENTS.md}\,{\color{muted}\scriptsize p.\,\pageref{c:protocols-6}} {\color{muted}\scriptsize Protocols and Interoperability Standards}\par\vspace{2pt}
\hyperref[c:protocols-7]{Agent Client Protocol}\,{\color{muted}\scriptsize p.\,\pageref{c:protocols-7}} {\color{muted}\scriptsize Protocols and Interoperability Standards}\par\vspace{2pt}
\hyperref[c:protocols-8]{AG-UI}\,{\color{muted}\scriptsize p.\,\pageref{c:protocols-8}} {\color{muted}\scriptsize Protocols and Interoperability Standards}\par\vspace{2pt}
\hyperref[c:protocols-9]{Agentic commerce protocols}\,{\color{muted}\scriptsize p.\,\pageref{c:protocols-9}} {\color{muted}\scriptsize Protocols and Interoperability Standards}\par\vspace{2pt}
\hyperref[c:coding-1]{Vibe coding}\,{\color{muted}\scriptsize p.\,\pageref{c:coding-1}} {\color{muted}\scriptsize Agentic Coding and Developer Workflow}\par\vspace{2pt}
\hyperref[c:coding-2]{Agentic engineering / Software 3.0}\,{\color{muted}\scriptsize p.\,\pageref{c:coding-2}} {\color{muted}\scriptsize Agentic Coding and Developer Workflow}\par\vspace{2pt}
\hyperref[c:coding-3]{Spec-driven development}\,{\color{muted}\scriptsize p.\,\pageref{c:coding-3}} {\color{muted}\scriptsize Agentic Coding and Developer Workflow}\par\vspace{2pt}
\hyperref[c:coding-4]{Ralph loop}\,{\color{muted}\scriptsize p.\,\pageref{c:coding-4}} {\color{muted}\scriptsize Agentic Coding and Developer Workflow}\par\vspace{2pt}
\hyperref[c:coding-5]{Hooks, plugins and slash commands}\,{\color{muted}\scriptsize p.\,\pageref{c:coding-5}} {\color{muted}\scriptsize Agentic Coding and Developer Workflow}\par\vspace{2pt}
\hyperref[c:coding-6]{Parallel agents with git worktrees}\,{\color{muted}\scriptsize p.\,\pageref{c:coding-6}} {\color{muted}\scriptsize Agentic Coding and Developer Workflow}\par\vspace{2pt}
\hyperref[c:context-0]{Context engineering}\,{\color{muted}\scriptsize p.\,\pageref{c:context-0}} {\color{muted}\scriptsize Context, Memory and Retrieval}\par\vspace{2pt}
\hyperref[c:context-1]{Context rot}\,{\color{muted}\scriptsize p.\,\pageref{c:context-1}} {\color{muted}\scriptsize Context, Memory and Retrieval}\par\vspace{2pt}
\hyperref[c:context-2]{Compaction and context editing}\,{\color{muted}\scriptsize p.\,\pageref{c:context-2}} {\color{muted}\scriptsize Context, Memory and Retrieval}\par\vspace{2pt}
\hyperref[c:context-6]{Agentic search and deep research}\,{\color{muted}\scriptsize p.\,\pageref{c:context-6}} {\color{muted}\scriptsize Context, Memory and Retrieval}\par\vspace{2pt}
\hyperref[c:reasoning-1]{Thinking budgets, adaptive thinking and effort}\,{\color{muted}\scriptsize p.\,\pageref{c:reasoning-1}} {\color{muted}\scriptsize Reasoning and Post-Training}\par\vspace{2pt}
\hyperref[c:reasoning-5]{RL environments and verifiers}\,{\color{muted}\scriptsize p.\,\pageref{c:reasoning-5}} {\color{muted}\scriptsize Reasoning and Post-Training}\par\vspace{2pt}
\hyperref[c:reasoning-10]{Self-play and self-generated curricula}\,{\color{muted}\scriptsize p.\,\pageref{c:reasoning-10}} {\color{muted}\scriptsize Reasoning and Post-Training}\par\vspace{2pt}
\hyperref[c:arch-2]{Sparse and compressed attention}\,{\color{muted}\scriptsize p.\,\pageref{c:arch-2}} {\color{muted}\scriptsize Model Architectures}\par\vspace{2pt}
\hyperref[c:arch-5]{Diffusion language models}\,{\color{muted}\scriptsize p.\,\pageref{c:arch-5}} {\color{muted}\scriptsize Model Architectures}\par\vspace{2pt}
\hyperref[c:arch-10]{Tiny recursive reasoning models}\,{\color{muted}\scriptsize p.\,\pageref{c:arch-10}} {\color{muted}\scriptsize Model Architectures}\par\vspace{2pt}
\hyperref[c:arch-12]{Architecture overviews}\,{\color{muted}\scriptsize p.\,\pageref{c:arch-12}} {\color{muted}\scriptsize Model Architectures}\par\vspace{2pt}
\hyperref[c:inference-8]{Inference benchmarking}\,{\color{muted}\scriptsize p.\,\pageref{c:inference-8}} {\color{muted}\scriptsize Inference and Serving}\par\vspace{2pt}
\hyperref[c:quant-1]{NVFP4}\,{\color{muted}\scriptsize p.\,\pageref{c:quant-1}} {\color{muted}\scriptsize Quantization and Number Formats}\par\vspace{2pt}
\hyperref[c:hardware-4]{ESUN and SUE}\,{\color{muted}\scriptsize p.\,\pageref{c:hardware-4}} {\color{muted}\scriptsize AI Hardware, Interconnect and Datacenter Standards}\par\vspace{2pt}
\hyperref[c:hardware-10]{Co-packaged optics}\,{\color{muted}\scriptsize p.\,\pageref{c:hardware-10}} {\color{muted}\scriptsize AI Hardware, Interconnect and Datacenter Standards}\par\vspace{2pt}
\hyperref[c:vlm-3]{Thinking with images}\,{\color{muted}\scriptsize p.\,\pageref{c:vlm-3}} {\color{muted}\scriptsize Vision-Language and Multimodal Models}\par\vspace{2pt}
\hyperref[c:vlm-8]{Native image generation and editing}\,{\color{muted}\scriptsize p.\,\pageref{c:vlm-8}} {\color{muted}\scriptsize Vision-Language and Multimodal Models}\par\vspace{2pt}
\hyperref[c:evals-1]{Harness sensitivity of benchmarks}\,{\color{muted}\scriptsize p.\,\pageref{c:evals-1}} {\color{muted}\scriptsize Evaluation and Benchmarks}\par\vspace{2pt}
\hyperref[c:evals-7]{Economic value evals}\,{\color{muted}\scriptsize p.\,\pageref{c:evals-7}} {\color{muted}\scriptsize Evaluation and Benchmarks}\par\vspace{2pt}
\hyperref[c:evals-8]{Task time horizon}\,{\color{muted}\scriptsize p.\,\pageref{c:evals-8}} {\color{muted}\scriptsize Evaluation and Benchmarks}\par\vspace{2pt}
\hyperref[c:safety-1]{AI for cyber defense}\,{\color{muted}\scriptsize p.\,\pageref{c:safety-1}} {\color{muted}\scriptsize Safety, Security, Interpretability and Governance}\par\vspace{2pt}
\hyperref[c:safety-4]{MCP and tool security}\,{\color{muted}\scriptsize p.\,\pageref{c:safety-4}} {\color{muted}\scriptsize Safety, Security, Interpretability and Governance}\par\vspace{2pt}
\hyperref[c:safety-7]{Persona vectors and introspection}\,{\color{muted}\scriptsize p.\,\pageref{c:safety-7}} {\color{muted}\scriptsize Safety, Security, Interpretability and Governance}\par\vspace{2pt}
\hyperref[c:safety-8]{Chain-of-thought monitorability}\,{\color{muted}\scriptsize p.\,\pageref{c:safety-8}} {\color{muted}\scriptsize Safety, Security, Interpretability and Governance}\par\vspace{2pt}
\hyperref[c:safety-10]{Model welfare}\,{\color{muted}\scriptsize p.\,\pageref{c:safety-10}} {\color{muted}\scriptsize Safety, Security, Interpretability and Governance}\par\vspace{2pt}
\hyperref[c:practice-2]{Services as software}\,{\color{muted}\scriptsize p.\,\pageref{c:practice-2}} {\color{muted}\scriptsize Roles, Deployment and AI Engineering Practice}\par\vspace{2pt}
\end{multicols}
\Needspace{5\baselineskip}\section{2024 (39)}
\begin{multicols}{2}\small\raggedright
\hyperref[c:agents-0]{Agent harness}\,{\color{muted}\scriptsize p.\,\pageref{c:agents-0}} {\color{muted}\scriptsize Agents and Harnesses}\par\vspace{2pt}
\hyperref[c:agents-4]{Workflows vs. agents}\,{\color{muted}\scriptsize p.\,\pageref{c:agents-4}} {\color{muted}\scriptsize Agents and Harnesses}\par\vspace{2pt}
\hyperref[c:agents-9]{Computer use and browser agents}\,{\color{muted}\scriptsize p.\,\pageref{c:agents-9}} {\color{muted}\scriptsize Agents and Harnesses}\par\vspace{2pt}
\hyperref[c:agents-11]{Agent sandboxes}\,{\color{muted}\scriptsize p.\,\pageref{c:agents-11}} {\color{muted}\scriptsize Agents and Harnesses}\par\vspace{2pt}
\hyperref[c:agents-13]{Durable execution and agent trajectories}\,{\color{muted}\scriptsize p.\,\pageref{c:agents-13}} {\color{muted}\scriptsize Agents and Harnesses}\par\vspace{2pt}
\hyperref[c:protocols-0]{Model Context Protocol}\,{\color{muted}\scriptsize p.\,\pageref{c:protocols-0}} {\color{muted}\scriptsize Protocols and Interoperability Standards}\par\vspace{2pt}
\hyperref[c:protocols-10]{llms.txt}\,{\color{muted}\scriptsize p.\,\pageref{c:protocols-10}} {\color{muted}\scriptsize Protocols and Interoperability Standards}\par\vspace{2pt}
\hyperref[c:protocols-11]{OpenTelemetry GenAI conventions}\,{\color{muted}\scriptsize p.\,\pageref{c:protocols-11}} {\color{muted}\scriptsize Protocols and Interoperability Standards}\par\vspace{2pt}
\hyperref[c:coding-0]{Coding agents}\,{\color{muted}\scriptsize p.\,\pageref{c:coding-0}} {\color{muted}\scriptsize Agentic Coding and Developer Workflow}\par\vspace{2pt}
\hyperref[c:context-3]{Prompt caching}\,{\color{muted}\scriptsize p.\,\pageref{c:context-3}} {\color{muted}\scriptsize Context, Memory and Retrieval}\par\vspace{2pt}
\hyperref[c:context-7]{GraphRAG}\,{\color{muted}\scriptsize p.\,\pageref{c:context-7}} {\color{muted}\scriptsize Context, Memory and Retrieval}\par\vspace{2pt}
\hyperref[c:reasoning-0]{Reasoning models and test-time compute}\,{\color{muted}\scriptsize p.\,\pageref{c:reasoning-0}} {\color{muted}\scriptsize Reasoning and Post-Training}\par\vspace{2pt}
\hyperref[c:reasoning-2]{RL with verifiable rewards}\,{\color{muted}\scriptsize p.\,\pageref{c:reasoning-2}} {\color{muted}\scriptsize Reasoning and Post-Training}\par\vspace{2pt}
\hyperref[c:reasoning-4]{GRPO and its successors}\,{\color{muted}\scriptsize p.\,\pageref{c:reasoning-4}} {\color{muted}\scriptsize Reasoning and Post-Training}\par\vspace{2pt}
\hyperref[c:reasoning-11]{Continual learning}\,{\color{muted}\scriptsize p.\,\pageref{c:reasoning-11}} {\color{muted}\scriptsize Reasoning and Post-Training}\par\vspace{2pt}
\hyperref[c:arch-9]{Tokenizer-free and byte-level models}\,{\color{muted}\scriptsize p.\,\pageref{c:arch-9}} {\color{muted}\scriptsize Model Architectures}\par\vspace{2pt}
\hyperref[c:arch-11]{Small language models}\,{\color{muted}\scriptsize p.\,\pageref{c:arch-11}} {\color{muted}\scriptsize Model Architectures}\par\vspace{2pt}
\hyperref[c:training-2]{Muon and modern optimizers}\,{\color{muted}\scriptsize p.\,\pageref{c:training-2}} {\color{muted}\scriptsize Training at Scale and Fine-Tuning}\par\vspace{2pt}
\hyperref[c:training-4]{Mid-training}\,{\color{muted}\scriptsize p.\,\pageref{c:training-4}} {\color{muted}\scriptsize Training at Scale and Fine-Tuning}\par\vspace{2pt}
\hyperref[c:inference-3]{Prefill/decode disaggregation}\,{\color{muted}\scriptsize p.\,\pageref{c:inference-3}} {\color{muted}\scriptsize Inference and Serving}\par\vspace{2pt}
\hyperref[c:inference-4]{KV cache offloading}\,{\color{muted}\scriptsize p.\,\pageref{c:inference-4}} {\color{muted}\scriptsize Inference and Serving}\par\vspace{2pt}
\hyperref[c:quant-4]{Quantization-aware training}\,{\color{muted}\scriptsize p.\,\pageref{c:quant-4}} {\color{muted}\scriptsize Quantization and Number Formats}\par\vspace{2pt}
\hyperref[c:quant-6]{Ternary and 1-bit models}\,{\color{muted}\scriptsize p.\,\pageref{c:quant-6}} {\color{muted}\scriptsize Quantization and Number Formats}\par\vspace{2pt}
\hyperref[c:hardware-0]{Rack-scale systems}\,{\color{muted}\scriptsize p.\,\pageref{c:hardware-0}} {\color{muted}\scriptsize AI Hardware, Interconnect and Datacenter Standards}\par\vspace{2pt}
\hyperref[c:hardware-2]{Scale-up vs. scale-out networking}\,{\color{muted}\scriptsize p.\,\pageref{c:hardware-2}} {\color{muted}\scriptsize AI Hardware, Interconnect and Datacenter Standards}\par\vspace{2pt}
\hyperref[c:hardware-3]{UALink}\,{\color{muted}\scriptsize p.\,\pageref{c:hardware-3}} {\color{muted}\scriptsize AI Hardware, Interconnect and Datacenter Standards}\par\vspace{2pt}
\hyperref[c:hardware-12]{AI factories, power and cooling}\,{\color{muted}\scriptsize p.\,\pageref{c:hardware-12}} {\color{muted}\scriptsize AI Hardware, Interconnect and Datacenter Standards}\par\vspace{2pt}
\hyperref[c:vlm-1]{Open VLM families}\,{\color{muted}\scriptsize p.\,\pageref{c:vlm-1}} {\color{muted}\scriptsize Vision-Language and Multimodal Models}\par\vspace{2pt}
\hyperref[c:vlm-4]{Document AI and VLM-based OCR}\,{\color{muted}\scriptsize p.\,\pageref{c:vlm-4}} {\color{muted}\scriptsize Vision-Language and Multimodal Models}\par\vspace{2pt}
\hyperref[c:vlm-5]{Visual document retrieval}\,{\color{muted}\scriptsize p.\,\pageref{c:vlm-5}} {\color{muted}\scriptsize Vision-Language and Multimodal Models}\par\vspace{2pt}
\hyperref[c:vlm-6]{GUI grounding agents}\,{\color{muted}\scriptsize p.\,\pageref{c:vlm-6}} {\color{muted}\scriptsize Vision-Language and Multimodal Models}\par\vspace{2pt}
\hyperref[c:vlm-7]{Omni and speech-to-speech models}\,{\color{muted}\scriptsize p.\,\pageref{c:vlm-7}} {\color{muted}\scriptsize Vision-Language and Multimodal Models}\par\vspace{2pt}
\hyperref[c:embodied-6]{Video generation}\,{\color{muted}\scriptsize p.\,\pageref{c:embodied-6}} {\color{muted}\scriptsize Embodied AI, World Models and Generative Media}\par\vspace{2pt}
\hyperref[c:evals-5]{Agent and tool-use benchmarks}\,{\color{muted}\scriptsize p.\,\pageref{c:evals-5}} {\color{muted}\scriptsize Evaluation and Benchmarks}\par\vspace{2pt}
\hyperref[c:safety-9]{Misalignment research}\,{\color{muted}\scriptsize p.\,\pageref{c:safety-9}} {\color{muted}\scriptsize Safety, Security, Interpretability and Governance}\par\vspace{2pt}
\hyperref[c:safety-11]{EU AI Act and the Digital Omnibus}\,{\color{muted}\scriptsize p.\,\pageref{c:safety-11}} {\color{muted}\scriptsize Safety, Security, Interpretability and Governance}\par\vspace{2pt}
\hyperref[c:practice-4]{Eval-driven development}\,{\color{muted}\scriptsize p.\,\pageref{c:practice-4}} {\color{muted}\scriptsize Roles, Deployment and AI Engineering Practice}\par\vspace{2pt}
\hyperref[c:practice-6]{Human-in-the-loop and approval design}\,{\color{muted}\scriptsize p.\,\pageref{c:practice-6}} {\color{muted}\scriptsize Roles, Deployment and AI Engineering Practice}\par\vspace{2pt}
\hyperref[c:learn-3]{Curated lists}\,{\color{muted}\scriptsize p.\,\pageref{c:learn-3}} {\color{muted}\scriptsize Keep Learning: Curated Resources}\par\vspace{2pt}
\end{multicols}
\Needspace{5\baselineskip}\section{2023 (40)}
\begin{multicols}{2}\small\raggedright
\hyperref[c:landscape-3]{Open-weight frontier models}\,{\color{muted}\scriptsize p.\,\pageref{c:landscape-3}} {\color{muted}\scriptsize The Frontier Landscape (September 2026)}\par\vspace{2pt}
\hyperref[c:landscape-4]{Model release trackers}\,{\color{muted}\scriptsize p.\,\pageref{c:landscape-4}} {\color{muted}\scriptsize The Frontier Landscape (September 2026)}\par\vspace{2pt}
\hyperref[c:agents-5]{Tool use / function calling}\,{\color{muted}\scriptsize p.\,\pageref{c:agents-5}} {\color{muted}\scriptsize Agents and Harnesses}\par\vspace{2pt}
\hyperref[c:agents-7]{Multi-agent systems and subagents}\,{\color{muted}\scriptsize p.\,\pageref{c:agents-7}} {\color{muted}\scriptsize Agents and Harnesses}\par\vspace{2pt}
\hyperref[c:agents-12]{Agent frameworks and SDKs}\,{\color{muted}\scriptsize p.\,\pageref{c:agents-12}} {\color{muted}\scriptsize Agents and Harnesses}\par\vspace{2pt}
\hyperref[c:context-4]{Agent memory}\,{\color{muted}\scriptsize p.\,\pageref{c:context-4}} {\color{muted}\scriptsize Context, Memory and Retrieval}\par\vspace{2pt}
\hyperref[c:context-9]{Structured outputs and constrained decoding}\,{\color{muted}\scriptsize p.\,\pageref{c:context-9}} {\color{muted}\scriptsize Context, Memory and Retrieval}\par\vspace{2pt}
\hyperref[c:reasoning-8]{Process reward models}\,{\color{muted}\scriptsize p.\,\pageref{c:reasoning-8}} {\color{muted}\scriptsize Reasoning and Post-Training}\par\vspace{2pt}
\hyperref[c:arch-1]{GQA and Multi-head Latent Attention}\,{\color{muted}\scriptsize p.\,\pageref{c:arch-1}} {\color{muted}\scriptsize Model Architectures}\par\vspace{2pt}
\hyperref[c:arch-4]{Hybrid linear-attention and state space models}\,{\color{muted}\scriptsize p.\,\pageref{c:arch-4}} {\color{muted}\scriptsize Model Architectures}\par\vspace{2pt}
\hyperref[c:arch-6]{Multi-token prediction and speculative decoding}\,{\color{muted}\scriptsize p.\,\pageref{c:arch-6}} {\color{muted}\scriptsize Model Architectures}\par\vspace{2pt}
\hyperref[c:arch-8]{Attention sinks}\,{\color{muted}\scriptsize p.\,\pageref{c:arch-8}} {\color{muted}\scriptsize Model Architectures}\par\vspace{2pt}
\hyperref[c:training-3]{Data curation and synthetic data}\,{\color{muted}\scriptsize p.\,\pageref{c:training-3}} {\color{muted}\scriptsize Training at Scale and Fine-Tuning}\par\vspace{2pt}
\hyperref[c:training-6]{Fine-tuning toolkits and APIs}\,{\color{muted}\scriptsize p.\,\pageref{c:training-6}} {\color{muted}\scriptsize Training at Scale and Fine-Tuning}\par\vspace{2pt}
\hyperref[c:training-7]{Model merging}\,{\color{muted}\scriptsize p.\,\pageref{c:training-7}} {\color{muted}\scriptsize Training at Scale and Fine-Tuning}\par\vspace{2pt}
\hyperref[c:training-8]{Decentralized training}\,{\color{muted}\scriptsize p.\,\pageref{c:training-8}} {\color{muted}\scriptsize Training at Scale and Fine-Tuning}\par\vspace{2pt}
\hyperref[c:training-9]{Build-it-yourself references}\,{\color{muted}\scriptsize p.\,\pageref{c:training-9}} {\color{muted}\scriptsize Training at Scale and Fine-Tuning}\par\vspace{2pt}
\hyperref[c:inference-0]{Inference engines}\,{\color{muted}\scriptsize p.\,\pageref{c:inference-0}} {\color{muted}\scriptsize Inference and Serving}\par\vspace{2pt}
\hyperref[c:inference-1]{KV cache and PagedAttention}\,{\color{muted}\scriptsize p.\,\pageref{c:inference-1}} {\color{muted}\scriptsize Inference and Serving}\par\vspace{2pt}
\hyperref[c:inference-6]{On-device and edge inference}\,{\color{muted}\scriptsize p.\,\pageref{c:inference-6}} {\color{muted}\scriptsize Inference and Serving}\par\vspace{2pt}
\hyperref[c:inference-7]{LLM gateways and routers}\,{\color{muted}\scriptsize p.\,\pageref{c:inference-7}} {\color{muted}\scriptsize Inference and Serving}\par\vspace{2pt}
\hyperref[c:quant-0]{Microscaling formats}\,{\color{muted}\scriptsize p.\,\pageref{c:quant-0}} {\color{muted}\scriptsize Quantization and Number Formats}\par\vspace{2pt}
\hyperref[c:quant-5]{GGUF and local model formats}\,{\color{muted}\scriptsize p.\,\pageref{c:quant-5}} {\color{muted}\scriptsize Quantization and Number Formats}\par\vspace{2pt}
\hyperref[c:hardware-5]{Ultra Ethernet}\,{\color{muted}\scriptsize p.\,\pageref{c:hardware-5}} {\color{muted}\scriptsize AI Hardware, Interconnect and Datacenter Standards}\par\vspace{2pt}
\hyperref[c:vlm-0]{VLM architecture}\,{\color{muted}\scriptsize p.\,\pageref{c:vlm-0}} {\color{muted}\scriptsize Vision-Language and Multimodal Models}\par\vspace{2pt}
\hyperref[c:vlm-2]{Native and dynamic resolution}\,{\color{muted}\scriptsize p.\,\pageref{c:vlm-2}} {\color{muted}\scriptsize Vision-Language and Multimodal Models}\par\vspace{2pt}
\hyperref[c:vlm-9]{Multimodal evaluation}\,{\color{muted}\scriptsize p.\,\pageref{c:vlm-9}} {\color{muted}\scriptsize Vision-Language and Multimodal Models}\par\vspace{2pt}
\hyperref[c:embodied-0]{Vision-language-action models}\,{\color{muted}\scriptsize p.\,\pageref{c:embodied-0}} {\color{muted}\scriptsize Embodied AI, World Models and Generative Media}\par\vspace{2pt}
\hyperref[c:evals-2]{LLM-as-a-judge}\,{\color{muted}\scriptsize p.\,\pageref{c:evals-2}} {\color{muted}\scriptsize Evaluation and Benchmarks}\par\vspace{2pt}
\hyperref[c:evals-3]{Agentic coding benchmarks}\,{\color{muted}\scriptsize p.\,\pageref{c:evals-3}} {\color{muted}\scriptsize Evaluation and Benchmarks}\par\vspace{2pt}
\hyperref[c:evals-6]{Frontier knowledge benchmarks}\,{\color{muted}\scriptsize p.\,\pageref{c:evals-6}} {\color{muted}\scriptsize Evaluation and Benchmarks}\par\vspace{2pt}
\hyperref[c:evals-9]{Leaderboards and trackers}\,{\color{muted}\scriptsize p.\,\pageref{c:evals-9}} {\color{muted}\scriptsize Evaluation and Benchmarks}\par\vspace{2pt}
\hyperref[c:evals-10]{Contamination and saturation}\,{\color{muted}\scriptsize p.\,\pageref{c:evals-10}} {\color{muted}\scriptsize Evaluation and Benchmarks}\par\vspace{2pt}
\hyperref[c:safety-0]{Capability thresholds and frontier safety frameworks}\,{\color{muted}\scriptsize p.\,\pageref{c:safety-0}} {\color{muted}\scriptsize Safety, Security, Interpretability and Governance}\par\vspace{2pt}
\hyperref[c:safety-5]{Guardrails and classifiers}\,{\color{muted}\scriptsize p.\,\pageref{c:safety-5}} {\color{muted}\scriptsize Safety, Security, Interpretability and Governance}\par\vspace{2pt}
\hyperref[c:safety-12]{Other regulation and standards}\,{\color{muted}\scriptsize p.\,\pageref{c:safety-12}} {\color{muted}\scriptsize Safety, Security, Interpretability and Governance}\par\vspace{2pt}
\hyperref[c:safety-13]{Government AI safety institutes}\,{\color{muted}\scriptsize p.\,\pageref{c:safety-13}} {\color{muted}\scriptsize Safety, Security, Interpretability and Governance}\par\vspace{2pt}
\hyperref[c:practice-1]{AI Engineer}\,{\color{muted}\scriptsize p.\,\pageref{c:practice-1}} {\color{muted}\scriptsize Roles, Deployment and AI Engineering Practice}\par\vspace{2pt}
\hyperref[c:practice-5]{LLMOps and observability}\,{\color{muted}\scriptsize p.\,\pageref{c:practice-5}} {\color{muted}\scriptsize Roles, Deployment and AI Engineering Practice}\par\vspace{2pt}
\hyperref[c:practice-7]{Open weights vs. open source AI}\,{\color{muted}\scriptsize p.\,\pageref{c:practice-7}} {\color{muted}\scriptsize Roles, Deployment and AI Engineering Practice}\par\vspace{2pt}
\end{multicols}
\Needspace{5\baselineskip}\section{2022 (12)}
\begin{multicols}{2}\small\raggedright
\hyperref[c:agents-3]{Agentic loop / ReAct}\,{\color{muted}\scriptsize p.\,\pageref{c:agents-3}} {\color{muted}\scriptsize Agents and Harnesses}\par\vspace{2pt}
\hyperref[c:reasoning-6]{RLHF, DPO and preference optimization}\,{\color{muted}\scriptsize p.\,\pageref{c:reasoning-6}} {\color{muted}\scriptsize Reasoning and Post-Training}\par\vspace{2pt}
\hyperref[c:reasoning-7]{Constitutional AI and model specs}\,{\color{muted}\scriptsize p.\,\pageref{c:reasoning-7}} {\color{muted}\scriptsize Reasoning and Post-Training}\par\vspace{2pt}
\hyperref[c:inference-2]{Continuous batching and prefix caching}\,{\color{muted}\scriptsize p.\,\pageref{c:inference-2}} {\color{muted}\scriptsize Inference and Serving}\par\vspace{2pt}
\hyperref[c:inference-5]{Attention kernels}\,{\color{muted}\scriptsize p.\,\pageref{c:inference-5}} {\color{muted}\scriptsize Inference and Serving}\par\vspace{2pt}
\hyperref[c:quant-2]{FP8 training and inference}\,{\color{muted}\scriptsize p.\,\pageref{c:quant-2}} {\color{muted}\scriptsize Quantization and Number Formats}\par\vspace{2pt}
\hyperref[c:quant-3]{Post-training quantization}\,{\color{muted}\scriptsize p.\,\pageref{c:quant-3}} {\color{muted}\scriptsize Quantization and Number Formats}\par\vspace{2pt}
\hyperref[c:embodied-3]{JEPA}\,{\color{muted}\scriptsize p.\,\pageref{c:embodied-3}} {\color{muted}\scriptsize Embodied AI, World Models and Generative Media}\par\vspace{2pt}
\hyperref[c:embodied-5]{Diffusion transformers and flow matching}\,{\color{muted}\scriptsize p.\,\pageref{c:embodied-5}} {\color{muted}\scriptsize Embodied AI, World Models and Generative Media}\par\vspace{2pt}
\hyperref[c:safety-3]{Prompt injection and the lethal trifecta}\,{\color{muted}\scriptsize p.\,\pageref{c:safety-3}} {\color{muted}\scriptsize Safety, Security, Interpretability and Governance}\par\vspace{2pt}
\hyperref[c:practice-3]{Prompt engineering and prompt optimization}\,{\color{muted}\scriptsize p.\,\pageref{c:practice-3}} {\color{muted}\scriptsize Roles, Deployment and AI Engineering Practice}\par\vspace{2pt}
\hyperref[c:learn-1]{Courses}\,{\color{muted}\scriptsize p.\,\pageref{c:learn-1}} {\color{muted}\scriptsize Keep Learning: Curated Resources}\par\vspace{2pt}
\end{multicols}
\Needspace{5\baselineskip}\section{2021 (4)}
\begin{multicols}{2}\small\raggedright
\hyperref[c:arch-7]{Long-context methods}\,{\color{muted}\scriptsize p.\,\pageref{c:arch-7}} {\color{muted}\scriptsize Model Architectures}\par\vspace{2pt}
\hyperref[c:training-5]{LoRA and parameter-efficient fine-tuning}\,{\color{muted}\scriptsize p.\,\pageref{c:training-5}} {\color{muted}\scriptsize Training at Scale and Fine-Tuning}\par\vspace{2pt}
\hyperref[c:evals-0]{Eval harnesses}\,{\color{muted}\scriptsize p.\,\pageref{c:evals-0}} {\color{muted}\scriptsize Evaluation and Benchmarks}\par\vspace{2pt}
\hyperref[c:safety-14]{Content provenance and watermarking}\,{\color{muted}\scriptsize p.\,\pageref{c:safety-14}} {\color{muted}\scriptsize Safety, Security, Interpretability and Governance}\par\vspace{2pt}
\end{multicols}
\Needspace{5\baselineskip}\section{2020 (5)}
\begin{multicols}{2}\small\raggedright
\hyperref[c:context-5]{Retrieval-augmented generation}\,{\color{muted}\scriptsize p.\,\pageref{c:context-5}} {\color{muted}\scriptsize Context, Memory and Retrieval}\par\vspace{2pt}
\hyperref[c:context-8]{Embeddings, late interaction and Matryoshka}\,{\color{muted}\scriptsize p.\,\pageref{c:context-8}} {\color{muted}\scriptsize Context, Memory and Retrieval}\par\vspace{2pt}
\hyperref[c:training-1]{Scaling laws}\,{\color{muted}\scriptsize p.\,\pageref{c:training-1}} {\color{muted}\scriptsize Training at Scale and Fine-Tuning}\par\vspace{2pt}
\hyperref[c:embodied-7]{AI for science and algorithm discovery}\,{\color{muted}\scriptsize p.\,\pageref{c:embodied-7}} {\color{muted}\scriptsize Embodied AI, World Models and Generative Media}\par\vspace{2pt}
\hyperref[c:safety-6]{Mechanistic interpretability}\,{\color{muted}\scriptsize p.\,\pageref{c:safety-6}} {\color{muted}\scriptsize Safety, Security, Interpretability and Governance}\par\vspace{2pt}
\end{multicols}
\Needspace{5\baselineskip}\section{2019 and earlier (17)}
\begin{multicols}{2}\small\raggedright
\hyperref[c:reasoning-9]{Reward hacking}\,{\color{muted}\scriptsize p.\,\pageref{c:reasoning-9}} {\color{muted}\scriptsize Reasoning and Post-Training}\par\vspace{2pt}
\hyperref[c:arch-0]{Mixture of Experts}\,{\color{muted}\scriptsize p.\,\pageref{c:arch-0}} {\color{muted}\scriptsize Model Architectures}\par\vspace{2pt}
\hyperref[c:training-0]{Parallelism strategies}\,{\color{muted}\scriptsize p.\,\pageref{c:training-0}} {\color{muted}\scriptsize Training at Scale and Fine-Tuning}\par\vspace{2pt}
\hyperref[c:hardware-6]{NVLink and NVLink Fusion}\,{\color{muted}\scriptsize p.\,\pageref{c:hardware-6}} {\color{muted}\scriptsize AI Hardware, Interconnect and Datacenter Standards}\par\vspace{2pt}
\hyperref[c:hardware-7]{CXL, PCIe and UCIe}\,{\color{muted}\scriptsize p.\,\pageref{c:hardware-7}} {\color{muted}\scriptsize AI Hardware, Interconnect and Datacenter Standards}\par\vspace{2pt}
\hyperref[c:hardware-8]{High-bandwidth memory}\,{\color{muted}\scriptsize p.\,\pageref{c:hardware-8}} {\color{muted}\scriptsize AI Hardware, Interconnect and Datacenter Standards}\par\vspace{2pt}
\hyperref[c:hardware-9]{Advanced packaging}\,{\color{muted}\scriptsize p.\,\pageref{c:hardware-9}} {\color{muted}\scriptsize AI Hardware, Interconnect and Datacenter Standards}\par\vspace{2pt}
\hyperref[c:hardware-11]{Custom accelerators}\,{\color{muted}\scriptsize p.\,\pageref{c:hardware-11}} {\color{muted}\scriptsize AI Hardware, Interconnect and Datacenter Standards}\par\vspace{2pt}
\hyperref[c:hardware-13]{Kernel languages and compilers}\,{\color{muted}\scriptsize p.\,\pageref{c:hardware-13}} {\color{muted}\scriptsize AI Hardware, Interconnect and Datacenter Standards}\par\vspace{2pt}
\hyperref[c:hardware-14]{Collective communication}\,{\color{muted}\scriptsize p.\,\pageref{c:hardware-14}} {\color{muted}\scriptsize AI Hardware, Interconnect and Datacenter Standards}\par\vspace{2pt}
\hyperref[c:hardware-15]{Model interchange formats}\,{\color{muted}\scriptsize p.\,\pageref{c:hardware-15}} {\color{muted}\scriptsize AI Hardware, Interconnect and Datacenter Standards}\par\vspace{2pt}
\hyperref[c:hardware-16]{Hardware benchmarks}\,{\color{muted}\scriptsize p.\,\pageref{c:hardware-16}} {\color{muted}\scriptsize AI Hardware, Interconnect and Datacenter Standards}\par\vspace{2pt}
\hyperref[c:embodied-2]{World models}\,{\color{muted}\scriptsize p.\,\pageref{c:embodied-2}} {\color{muted}\scriptsize Embodied AI, World Models and Generative Media}\par\vspace{2pt}
\hyperref[c:embodied-4]{Robot simulation and sim-to-real}\,{\color{muted}\scriptsize p.\,\pageref{c:embodied-4}} {\color{muted}\scriptsize Embodied AI, World Models and Generative Media}\par\vspace{2pt}
\hyperref[c:practice-0]{Forward Deployed Engineer}\,{\color{muted}\scriptsize p.\,\pageref{c:practice-0}} {\color{muted}\scriptsize Roles, Deployment and AI Engineering Practice}\par\vspace{2pt}
\hyperref[c:learn-0]{Deep technical blogs}\,{\color{muted}\scriptsize p.\,\pageref{c:learn-0}} {\color{muted}\scriptsize Keep Learning: Curated Resources}\par\vspace{2pt}
\hyperref[c:learn-2]{Annual reports and data}\,{\color{muted}\scriptsize p.\,\pageref{c:learn-2}} {\color{muted}\scriptsize Keep Learning: Curated Resources}\par\vspace{2pt}
\end{multicols}
\chapter{Concepts by level}
\Needspace{5\baselineskip}\section{Beginner (52)}
\begin{multicols}{2}\small\raggedright
\hyperref[c:landscape-0]{Mythos-class tier and trusted access}\,{\color{muted}\scriptsize p.\,\pageref{c:landscape-0}}\par\vspace{1pt}
\hyperref[c:landscape-1]{GPT-6 Astra and the Critical cyber threshold}\,{\color{muted}\scriptsize p.\,\pageref{c:landscape-1}}\par\vspace{1pt}
\hyperref[c:landscape-3]{Open-weight frontier models}\,{\color{muted}\scriptsize p.\,\pageref{c:landscape-3}}\par\vspace{1pt}
\hyperref[c:landscape-4]{Model release trackers}\,{\color{muted}\scriptsize p.\,\pageref{c:landscape-4}}\par\vspace{1pt}
\hyperref[c:agents-0]{Agent harness}\,{\color{muted}\scriptsize p.\,\pageref{c:agents-0}}\par\vspace{1pt}
\hyperref[c:agents-3]{Agentic loop / ReAct}\,{\color{muted}\scriptsize p.\,\pageref{c:agents-3}}\par\vspace{1pt}
\hyperref[c:agents-4]{Workflows vs. agents}\,{\color{muted}\scriptsize p.\,\pageref{c:agents-4}}\par\vspace{1pt}
\hyperref[c:agents-5]{Tool use / function calling}\,{\color{muted}\scriptsize p.\,\pageref{c:agents-5}}\par\vspace{1pt}
\hyperref[c:agents-10]{Personal agents}\,{\color{muted}\scriptsize p.\,\pageref{c:agents-10}}\par\vspace{1pt}
\hyperref[c:agents-12]{Agent frameworks and SDKs}\,{\color{muted}\scriptsize p.\,\pageref{c:agents-12}}\par\vspace{1pt}
\hyperref[c:protocols-0]{Model Context Protocol}\,{\color{muted}\scriptsize p.\,\pageref{c:protocols-0}}\par\vspace{1pt}
\hyperref[c:protocols-4]{Agentic AI Foundation}\,{\color{muted}\scriptsize p.\,\pageref{c:protocols-4}}\par\vspace{1pt}
\hyperref[c:protocols-5]{Agent Skills}\,{\color{muted}\scriptsize p.\,\pageref{c:protocols-5}}\par\vspace{1pt}
\hyperref[c:protocols-6]{AGENTS.md}\,{\color{muted}\scriptsize p.\,\pageref{c:protocols-6}}\par\vspace{1pt}
\hyperref[c:protocols-10]{llms.txt}\,{\color{muted}\scriptsize p.\,\pageref{c:protocols-10}}\par\vspace{1pt}
\hyperref[c:coding-0]{Coding agents}\,{\color{muted}\scriptsize p.\,\pageref{c:coding-0}}\par\vspace{1pt}
\hyperref[c:coding-1]{Vibe coding}\,{\color{muted}\scriptsize p.\,\pageref{c:coding-1}}\par\vspace{1pt}
\hyperref[c:context-0]{Context engineering}\,{\color{muted}\scriptsize p.\,\pageref{c:context-0}}\par\vspace{1pt}
\hyperref[c:context-1]{Context rot}\,{\color{muted}\scriptsize p.\,\pageref{c:context-1}}\par\vspace{1pt}
\hyperref[c:context-5]{Retrieval-augmented generation}\,{\color{muted}\scriptsize p.\,\pageref{c:context-5}}\par\vspace{1pt}
\hyperref[c:reasoning-0]{Reasoning models and test-time compute}\,{\color{muted}\scriptsize p.\,\pageref{c:reasoning-0}}\par\vspace{1pt}
\hyperref[c:arch-0]{Mixture of Experts}\,{\color{muted}\scriptsize p.\,\pageref{c:arch-0}}\par\vspace{1pt}
\hyperref[c:arch-11]{Small language models}\,{\color{muted}\scriptsize p.\,\pageref{c:arch-11}}\par\vspace{1pt}
\hyperref[c:arch-12]{Architecture overviews}\,{\color{muted}\scriptsize p.\,\pageref{c:arch-12}}\par\vspace{1pt}
\hyperref[c:training-6]{Fine-tuning toolkits and APIs}\,{\color{muted}\scriptsize p.\,\pageref{c:training-6}}\par\vspace{1pt}
\hyperref[c:training-9]{Build-it-yourself references}\,{\color{muted}\scriptsize p.\,\pageref{c:training-9}}\par\vspace{1pt}
\hyperref[c:inference-0]{Inference engines}\,{\color{muted}\scriptsize p.\,\pageref{c:inference-0}}\par\vspace{1pt}
\hyperref[c:inference-6]{On-device and edge inference}\,{\color{muted}\scriptsize p.\,\pageref{c:inference-6}}\par\vspace{1pt}
\hyperref[c:inference-7]{LLM gateways and routers}\,{\color{muted}\scriptsize p.\,\pageref{c:inference-7}}\par\vspace{1pt}
\hyperref[c:quant-5]{GGUF and local model formats}\,{\color{muted}\scriptsize p.\,\pageref{c:quant-5}}\par\vspace{1pt}
\hyperref[c:hardware-15]{Model interchange formats}\,{\color{muted}\scriptsize p.\,\pageref{c:hardware-15}}\par\vspace{1pt}
\hyperref[c:hardware-16]{Hardware benchmarks}\,{\color{muted}\scriptsize p.\,\pageref{c:hardware-16}}\par\vspace{1pt}
\hyperref[c:vlm-0]{VLM architecture}\,{\color{muted}\scriptsize p.\,\pageref{c:vlm-0}}\par\vspace{1pt}
\hyperref[c:vlm-1]{Open VLM families}\,{\color{muted}\scriptsize p.\,\pageref{c:vlm-1}}\par\vspace{1pt}
\hyperref[c:vlm-8]{Native image generation and editing}\,{\color{muted}\scriptsize p.\,\pageref{c:vlm-8}}\par\vspace{1pt}
\hyperref[c:embodied-6]{Video generation}\,{\color{muted}\scriptsize p.\,\pageref{c:embodied-6}}\par\vspace{1pt}
\hyperref[c:evals-2]{LLM-as-a-judge}\,{\color{muted}\scriptsize p.\,\pageref{c:evals-2}}\par\vspace{1pt}
\hyperref[c:evals-3]{Agentic coding benchmarks}\,{\color{muted}\scriptsize p.\,\pageref{c:evals-3}}\par\vspace{1pt}
\hyperref[c:evals-7]{Economic value evals}\,{\color{muted}\scriptsize p.\,\pageref{c:evals-7}}\par\vspace{1pt}
\hyperref[c:evals-9]{Leaderboards and trackers}\,{\color{muted}\scriptsize p.\,\pageref{c:evals-9}}\par\vspace{1pt}
\hyperref[c:safety-2]{Export controls on frontier models}\,{\color{muted}\scriptsize p.\,\pageref{c:safety-2}}\par\vspace{1pt}
\hyperref[c:safety-3]{Prompt injection and the lethal trifecta}\,{\color{muted}\scriptsize p.\,\pageref{c:safety-3}}\par\vspace{1pt}
\hyperref[c:safety-13]{Government AI safety institutes}\,{\color{muted}\scriptsize p.\,\pageref{c:safety-13}}\par\vspace{1pt}
\hyperref[c:safety-14]{Content provenance and watermarking}\,{\color{muted}\scriptsize p.\,\pageref{c:safety-14}}\par\vspace{1pt}
\hyperref[c:practice-0]{Forward Deployed Engineer}\,{\color{muted}\scriptsize p.\,\pageref{c:practice-0}}\par\vspace{1pt}
\hyperref[c:practice-1]{AI Engineer}\,{\color{muted}\scriptsize p.\,\pageref{c:practice-1}}\par\vspace{1pt}
\hyperref[c:practice-3]{Prompt engineering and prompt optimization}\,{\color{muted}\scriptsize p.\,\pageref{c:practice-3}}\par\vspace{1pt}
\hyperref[c:practice-7]{Open weights vs. open source AI}\,{\color{muted}\scriptsize p.\,\pageref{c:practice-7}}\par\vspace{1pt}
\hyperref[c:learn-0]{Deep technical blogs}\,{\color{muted}\scriptsize p.\,\pageref{c:learn-0}}\par\vspace{1pt}
\hyperref[c:learn-1]{Courses}\,{\color{muted}\scriptsize p.\,\pageref{c:learn-1}}\par\vspace{1pt}
\hyperref[c:learn-2]{Annual reports and data}\,{\color{muted}\scriptsize p.\,\pageref{c:learn-2}}\par\vspace{1pt}
\hyperref[c:learn-3]{Curated lists}\,{\color{muted}\scriptsize p.\,\pageref{c:learn-3}}\par\vspace{1pt}
\end{multicols}
\Needspace{5\baselineskip}\section{Intermediate (81)}
\begin{multicols}{2}\small\raggedright
\hyperref[c:landscape-2]{Tiered access for dual-use models}\,{\color{muted}\scriptsize p.\,\pageref{c:landscape-2}}\par\vspace{1pt}
\hyperref[c:agents-1]{Harness engineering}\,{\color{muted}\scriptsize p.\,\pageref{c:agents-1}}\par\vspace{1pt}
\hyperref[c:agents-6]{Tool search and programmatic tool calling}\,{\color{muted}\scriptsize p.\,\pageref{c:agents-6}}\par\vspace{1pt}
\hyperref[c:agents-7]{Multi-agent systems and subagents}\,{\color{muted}\scriptsize p.\,\pageref{c:agents-7}}\par\vspace{1pt}
\hyperref[c:agents-8]{Long-running and background agents}\,{\color{muted}\scriptsize p.\,\pageref{c:agents-8}}\par\vspace{1pt}
\hyperref[c:agents-9]{Computer use and browser agents}\,{\color{muted}\scriptsize p.\,\pageref{c:agents-9}}\par\vspace{1pt}
\hyperref[c:agents-11]{Agent sandboxes}\,{\color{muted}\scriptsize p.\,\pageref{c:agents-11}}\par\vspace{1pt}
\hyperref[c:protocols-1]{MCP Apps}\,{\color{muted}\scriptsize p.\,\pageref{c:protocols-1}}\par\vspace{1pt}
\hyperref[c:protocols-2]{MCP Tasks}\,{\color{muted}\scriptsize p.\,\pageref{c:protocols-2}}\par\vspace{1pt}
\hyperref[c:protocols-3]{Agent2Agent protocol}\,{\color{muted}\scriptsize p.\,\pageref{c:protocols-3}}\par\vspace{1pt}
\hyperref[c:protocols-7]{Agent Client Protocol}\,{\color{muted}\scriptsize p.\,\pageref{c:protocols-7}}\par\vspace{1pt}
\hyperref[c:protocols-8]{AG-UI}\,{\color{muted}\scriptsize p.\,\pageref{c:protocols-8}}\par\vspace{1pt}
\hyperref[c:protocols-9]{Agentic commerce protocols}\,{\color{muted}\scriptsize p.\,\pageref{c:protocols-9}}\par\vspace{1pt}
\hyperref[c:protocols-11]{OpenTelemetry GenAI conventions}\,{\color{muted}\scriptsize p.\,\pageref{c:protocols-11}}\par\vspace{1pt}
\hyperref[c:coding-2]{Agentic engineering / Software 3.0}\,{\color{muted}\scriptsize p.\,\pageref{c:coding-2}}\par\vspace{1pt}
\hyperref[c:coding-3]{Spec-driven development}\,{\color{muted}\scriptsize p.\,\pageref{c:coding-3}}\par\vspace{1pt}
\hyperref[c:coding-4]{Ralph loop}\,{\color{muted}\scriptsize p.\,\pageref{c:coding-4}}\par\vspace{1pt}
\hyperref[c:coding-5]{Hooks, plugins and slash commands}\,{\color{muted}\scriptsize p.\,\pageref{c:coding-5}}\par\vspace{1pt}
\hyperref[c:coding-6]{Parallel agents with git worktrees}\,{\color{muted}\scriptsize p.\,\pageref{c:coding-6}}\par\vspace{1pt}
\hyperref[c:coding-7]{The review bottleneck}\,{\color{muted}\scriptsize p.\,\pageref{c:coding-7}}\par\vspace{1pt}
\hyperref[c:context-2]{Compaction and context editing}\,{\color{muted}\scriptsize p.\,\pageref{c:context-2}}\par\vspace{1pt}
\hyperref[c:context-3]{Prompt caching}\,{\color{muted}\scriptsize p.\,\pageref{c:context-3}}\par\vspace{1pt}
\hyperref[c:context-4]{Agent memory}\,{\color{muted}\scriptsize p.\,\pageref{c:context-4}}\par\vspace{1pt}
\hyperref[c:context-6]{Agentic search and deep research}\,{\color{muted}\scriptsize p.\,\pageref{c:context-6}}\par\vspace{1pt}
\hyperref[c:context-7]{GraphRAG}\,{\color{muted}\scriptsize p.\,\pageref{c:context-7}}\par\vspace{1pt}
\hyperref[c:context-8]{Embeddings, late interaction and Matryoshka}\,{\color{muted}\scriptsize p.\,\pageref{c:context-8}}\par\vspace{1pt}
\hyperref[c:context-9]{Structured outputs and constrained decoding}\,{\color{muted}\scriptsize p.\,\pageref{c:context-9}}\par\vspace{1pt}
\hyperref[c:reasoning-1]{Thinking budgets, adaptive thinking and effort}\,{\color{muted}\scriptsize p.\,\pageref{c:reasoning-1}}\par\vspace{1pt}
\hyperref[c:reasoning-2]{RL with verifiable rewards}\,{\color{muted}\scriptsize p.\,\pageref{c:reasoning-2}}\par\vspace{1pt}
\hyperref[c:reasoning-6]{RLHF, DPO and preference optimization}\,{\color{muted}\scriptsize p.\,\pageref{c:reasoning-6}}\par\vspace{1pt}
\hyperref[c:reasoning-7]{Constitutional AI and model specs}\,{\color{muted}\scriptsize p.\,\pageref{c:reasoning-7}}\par\vspace{1pt}
\hyperref[c:reasoning-9]{Reward hacking}\,{\color{muted}\scriptsize p.\,\pageref{c:reasoning-9}}\par\vspace{1pt}
\hyperref[c:arch-1]{GQA and Multi-head Latent Attention}\,{\color{muted}\scriptsize p.\,\pageref{c:arch-1}}\par\vspace{1pt}
\hyperref[c:arch-6]{Multi-token prediction and speculative decoding}\,{\color{muted}\scriptsize p.\,\pageref{c:arch-6}}\par\vspace{1pt}
\hyperref[c:arch-7]{Long-context methods}\,{\color{muted}\scriptsize p.\,\pageref{c:arch-7}}\par\vspace{1pt}
\hyperref[c:training-1]{Scaling laws}\,{\color{muted}\scriptsize p.\,\pageref{c:training-1}}\par\vspace{1pt}
\hyperref[c:training-3]{Data curation and synthetic data}\,{\color{muted}\scriptsize p.\,\pageref{c:training-3}}\par\vspace{1pt}
\hyperref[c:training-4]{Mid-training}\,{\color{muted}\scriptsize p.\,\pageref{c:training-4}}\par\vspace{1pt}
\hyperref[c:training-5]{LoRA and parameter-efficient fine-tuning}\,{\color{muted}\scriptsize p.\,\pageref{c:training-5}}\par\vspace{1pt}
\hyperref[c:training-7]{Model merging}\,{\color{muted}\scriptsize p.\,\pageref{c:training-7}}\par\vspace{1pt}
\hyperref[c:inference-1]{KV cache and PagedAttention}\,{\color{muted}\scriptsize p.\,\pageref{c:inference-1}}\par\vspace{1pt}
\hyperref[c:inference-2]{Continuous batching and prefix caching}\,{\color{muted}\scriptsize p.\,\pageref{c:inference-2}}\par\vspace{1pt}
\hyperref[c:inference-8]{Inference benchmarking}\,{\color{muted}\scriptsize p.\,\pageref{c:inference-8}}\par\vspace{1pt}
\hyperref[c:quant-2]{FP8 training and inference}\,{\color{muted}\scriptsize p.\,\pageref{c:quant-2}}\par\vspace{1pt}
\hyperref[c:quant-3]{Post-training quantization}\,{\color{muted}\scriptsize p.\,\pageref{c:quant-3}}\par\vspace{1pt}
\hyperref[c:hardware-0]{Rack-scale systems}\,{\color{muted}\scriptsize p.\,\pageref{c:hardware-0}}\par\vspace{1pt}
\hyperref[c:hardware-1]{AMD Helios and Open Rack Wide}\,{\color{muted}\scriptsize p.\,\pageref{c:hardware-1}}\par\vspace{1pt}
\hyperref[c:hardware-2]{Scale-up vs. scale-out networking}\,{\color{muted}\scriptsize p.\,\pageref{c:hardware-2}}\par\vspace{1pt}
\hyperref[c:hardware-3]{UALink}\,{\color{muted}\scriptsize p.\,\pageref{c:hardware-3}}\par\vspace{1pt}
\hyperref[c:hardware-6]{NVLink and NVLink Fusion}\,{\color{muted}\scriptsize p.\,\pageref{c:hardware-6}}\par\vspace{1pt}
\hyperref[c:hardware-8]{High-bandwidth memory}\,{\color{muted}\scriptsize p.\,\pageref{c:hardware-8}}\par\vspace{1pt}
\hyperref[c:hardware-11]{Custom accelerators}\,{\color{muted}\scriptsize p.\,\pageref{c:hardware-11}}\par\vspace{1pt}
\hyperref[c:hardware-12]{AI factories, power and cooling}\,{\color{muted}\scriptsize p.\,\pageref{c:hardware-12}}\par\vspace{1pt}
\hyperref[c:vlm-2]{Native and dynamic resolution}\,{\color{muted}\scriptsize p.\,\pageref{c:vlm-2}}\par\vspace{1pt}
\hyperref[c:vlm-3]{Thinking with images}\,{\color{muted}\scriptsize p.\,\pageref{c:vlm-3}}\par\vspace{1pt}
\hyperref[c:vlm-4]{Document AI and VLM-based OCR}\,{\color{muted}\scriptsize p.\,\pageref{c:vlm-4}}\par\vspace{1pt}
\hyperref[c:vlm-5]{Visual document retrieval}\,{\color{muted}\scriptsize p.\,\pageref{c:vlm-5}}\par\vspace{1pt}
\hyperref[c:vlm-7]{Omni and speech-to-speech models}\,{\color{muted}\scriptsize p.\,\pageref{c:vlm-7}}\par\vspace{1pt}
\hyperref[c:vlm-9]{Multimodal evaluation}\,{\color{muted}\scriptsize p.\,\pageref{c:vlm-9}}\par\vspace{1pt}
\hyperref[c:embodied-0]{Vision-language-action models}\,{\color{muted}\scriptsize p.\,\pageref{c:embodied-0}}\par\vspace{1pt}
\hyperref[c:embodied-2]{World models}\,{\color{muted}\scriptsize p.\,\pageref{c:embodied-2}}\par\vspace{1pt}
\hyperref[c:embodied-4]{Robot simulation and sim-to-real}\,{\color{muted}\scriptsize p.\,\pageref{c:embodied-4}}\par\vspace{1pt}
\hyperref[c:embodied-7]{AI for science and algorithm discovery}\,{\color{muted}\scriptsize p.\,\pageref{c:embodied-7}}\par\vspace{1pt}
\hyperref[c:evals-0]{Eval harnesses}\,{\color{muted}\scriptsize p.\,\pageref{c:evals-0}}\par\vspace{1pt}
\hyperref[c:evals-1]{Harness sensitivity of benchmarks}\,{\color{muted}\scriptsize p.\,\pageref{c:evals-1}}\par\vspace{1pt}
\hyperref[c:evals-4]{ARC-AGI-3 and interactive reasoning}\,{\color{muted}\scriptsize p.\,\pageref{c:evals-4}}\par\vspace{1pt}
\hyperref[c:evals-5]{Agent and tool-use benchmarks}\,{\color{muted}\scriptsize p.\,\pageref{c:evals-5}}\par\vspace{1pt}
\hyperref[c:evals-6]{Frontier knowledge benchmarks}\,{\color{muted}\scriptsize p.\,\pageref{c:evals-6}}\par\vspace{1pt}
\hyperref[c:evals-8]{Task time horizon}\,{\color{muted}\scriptsize p.\,\pageref{c:evals-8}}\par\vspace{1pt}
\hyperref[c:evals-10]{Contamination and saturation}\,{\color{muted}\scriptsize p.\,\pageref{c:evals-10}}\par\vspace{1pt}
\hyperref[c:safety-0]{Capability thresholds and frontier safety frameworks}\,{\color{muted}\scriptsize p.\,\pageref{c:safety-0}}\par\vspace{1pt}
\hyperref[c:safety-1]{AI for cyber defense}\,{\color{muted}\scriptsize p.\,\pageref{c:safety-1}}\par\vspace{1pt}
\hyperref[c:safety-4]{MCP and tool security}\,{\color{muted}\scriptsize p.\,\pageref{c:safety-4}}\par\vspace{1pt}
\hyperref[c:safety-5]{Guardrails and classifiers}\,{\color{muted}\scriptsize p.\,\pageref{c:safety-5}}\par\vspace{1pt}
\hyperref[c:safety-10]{Model welfare}\,{\color{muted}\scriptsize p.\,\pageref{c:safety-10}}\par\vspace{1pt}
\hyperref[c:safety-11]{EU AI Act and the Digital Omnibus}\,{\color{muted}\scriptsize p.\,\pageref{c:safety-11}}\par\vspace{1pt}
\hyperref[c:safety-12]{Other regulation and standards}\,{\color{muted}\scriptsize p.\,\pageref{c:safety-12}}\par\vspace{1pt}
\hyperref[c:practice-2]{Services as software}\,{\color{muted}\scriptsize p.\,\pageref{c:practice-2}}\par\vspace{1pt}
\hyperref[c:practice-4]{Eval-driven development}\,{\color{muted}\scriptsize p.\,\pageref{c:practice-4}}\par\vspace{1pt}
\hyperref[c:practice-5]{LLMOps and observability}\,{\color{muted}\scriptsize p.\,\pageref{c:practice-5}}\par\vspace{1pt}
\hyperref[c:practice-6]{Human-in-the-loop and approval design}\,{\color{muted}\scriptsize p.\,\pageref{c:practice-6}}\par\vspace{1pt}
\end{multicols}
\Needspace{5\baselineskip}\section{Advanced (40)}
\begin{multicols}{2}\small\raggedright
\hyperref[c:agents-2]{Managed agents: decoupling brain, hands and session}\,{\color{muted}\scriptsize p.\,\pageref{c:agents-2}}\par\vspace{1pt}
\hyperref[c:agents-13]{Durable execution and agent trajectories}\,{\color{muted}\scriptsize p.\,\pageref{c:agents-13}}\par\vspace{1pt}
\hyperref[c:reasoning-3]{Multi-teacher on-policy distillation}\,{\color{muted}\scriptsize p.\,\pageref{c:reasoning-3}}\par\vspace{1pt}
\hyperref[c:reasoning-4]{GRPO and its successors}\,{\color{muted}\scriptsize p.\,\pageref{c:reasoning-4}}\par\vspace{1pt}
\hyperref[c:reasoning-5]{RL environments and verifiers}\,{\color{muted}\scriptsize p.\,\pageref{c:reasoning-5}}\par\vspace{1pt}
\hyperref[c:reasoning-8]{Process reward models}\,{\color{muted}\scriptsize p.\,\pageref{c:reasoning-8}}\par\vspace{1pt}
\hyperref[c:reasoning-10]{Self-play and self-generated curricula}\,{\color{muted}\scriptsize p.\,\pageref{c:reasoning-10}}\par\vspace{1pt}
\hyperref[c:reasoning-11]{Continual learning}\,{\color{muted}\scriptsize p.\,\pageref{c:reasoning-11}}\par\vspace{1pt}
\hyperref[c:arch-2]{Sparse and compressed attention}\,{\color{muted}\scriptsize p.\,\pageref{c:arch-2}}\par\vspace{1pt}
\hyperref[c:arch-3]{Manifold-constrained hyper-connections}\,{\color{muted}\scriptsize p.\,\pageref{c:arch-3}}\par\vspace{1pt}
\hyperref[c:arch-4]{Hybrid linear-attention and state space models}\,{\color{muted}\scriptsize p.\,\pageref{c:arch-4}}\par\vspace{1pt}
\hyperref[c:arch-5]{Diffusion language models}\,{\color{muted}\scriptsize p.\,\pageref{c:arch-5}}\par\vspace{1pt}
\hyperref[c:arch-8]{Attention sinks}\,{\color{muted}\scriptsize p.\,\pageref{c:arch-8}}\par\vspace{1pt}
\hyperref[c:arch-9]{Tokenizer-free and byte-level models}\,{\color{muted}\scriptsize p.\,\pageref{c:arch-9}}\par\vspace{1pt}
\hyperref[c:arch-10]{Tiny recursive reasoning models}\,{\color{muted}\scriptsize p.\,\pageref{c:arch-10}}\par\vspace{1pt}
\hyperref[c:training-0]{Parallelism strategies}\,{\color{muted}\scriptsize p.\,\pageref{c:training-0}}\par\vspace{1pt}
\hyperref[c:training-2]{Muon and modern optimizers}\,{\color{muted}\scriptsize p.\,\pageref{c:training-2}}\par\vspace{1pt}
\hyperref[c:training-8]{Decentralized training}\,{\color{muted}\scriptsize p.\,\pageref{c:training-8}}\par\vspace{1pt}
\hyperref[c:inference-3]{Prefill/decode disaggregation}\,{\color{muted}\scriptsize p.\,\pageref{c:inference-3}}\par\vspace{1pt}
\hyperref[c:inference-4]{KV cache offloading}\,{\color{muted}\scriptsize p.\,\pageref{c:inference-4}}\par\vspace{1pt}
\hyperref[c:inference-5]{Attention kernels}\,{\color{muted}\scriptsize p.\,\pageref{c:inference-5}}\par\vspace{1pt}
\hyperref[c:quant-0]{Microscaling formats}\,{\color{muted}\scriptsize p.\,\pageref{c:quant-0}}\par\vspace{1pt}
\hyperref[c:quant-1]{NVFP4}\,{\color{muted}\scriptsize p.\,\pageref{c:quant-1}}\par\vspace{1pt}
\hyperref[c:quant-4]{Quantization-aware training}\,{\color{muted}\scriptsize p.\,\pageref{c:quant-4}}\par\vspace{1pt}
\hyperref[c:quant-6]{Ternary and 1-bit models}\,{\color{muted}\scriptsize p.\,\pageref{c:quant-6}}\par\vspace{1pt}
\hyperref[c:hardware-4]{ESUN and SUE}\,{\color{muted}\scriptsize p.\,\pageref{c:hardware-4}}\par\vspace{1pt}
\hyperref[c:hardware-5]{Ultra Ethernet}\,{\color{muted}\scriptsize p.\,\pageref{c:hardware-5}}\par\vspace{1pt}
\hyperref[c:hardware-7]{CXL, PCIe and UCIe}\,{\color{muted}\scriptsize p.\,\pageref{c:hardware-7}}\par\vspace{1pt}
\hyperref[c:hardware-9]{Advanced packaging}\,{\color{muted}\scriptsize p.\,\pageref{c:hardware-9}}\par\vspace{1pt}
\hyperref[c:hardware-10]{Co-packaged optics}\,{\color{muted}\scriptsize p.\,\pageref{c:hardware-10}}\par\vspace{1pt}
\hyperref[c:hardware-13]{Kernel languages and compilers}\,{\color{muted}\scriptsize p.\,\pageref{c:hardware-13}}\par\vspace{1pt}
\hyperref[c:hardware-14]{Collective communication}\,{\color{muted}\scriptsize p.\,\pageref{c:hardware-14}}\par\vspace{1pt}
\hyperref[c:vlm-6]{GUI grounding agents}\,{\color{muted}\scriptsize p.\,\pageref{c:vlm-6}}\par\vspace{1pt}
\hyperref[c:embodied-1]{World action models}\,{\color{muted}\scriptsize p.\,\pageref{c:embodied-1}}\par\vspace{1pt}
\hyperref[c:embodied-3]{JEPA}\,{\color{muted}\scriptsize p.\,\pageref{c:embodied-3}}\par\vspace{1pt}
\hyperref[c:embodied-5]{Diffusion transformers and flow matching}\,{\color{muted}\scriptsize p.\,\pageref{c:embodied-5}}\par\vspace{1pt}
\hyperref[c:safety-6]{Mechanistic interpretability}\,{\color{muted}\scriptsize p.\,\pageref{c:safety-6}}\par\vspace{1pt}
\hyperref[c:safety-7]{Persona vectors and introspection}\,{\color{muted}\scriptsize p.\,\pageref{c:safety-7}}\par\vspace{1pt}
\hyperref[c:safety-8]{Chain-of-thought monitorability}\,{\color{muted}\scriptsize p.\,\pageref{c:safety-8}}\par\vspace{1pt}
\hyperref[c:safety-9]{Misalignment research}\,{\color{muted}\scriptsize p.\,\pageref{c:safety-9}}\par\vspace{1pt}
\end{multicols}
\chapter{Resources by type and year}
Every resource in the book, grouped by type and sorted newest first. Duplicates that support several concepts are listed once, with the first concept that cites them.

\section{GitHub (157)}
{\small\begin{longtable}{@{}p{0.07\linewidth}>{\raggedright}p{0.55\linewidth}>{\raggedright\arraybackslash}p{0.34\linewidth}@{}}
\textbf{\sffamily Year} & \textbf{\sffamily Resource} & \textbf{\sffamily Concept}\\ \hline\endhead
2026 & \href{https://github.com/AutoJunjie/awesome-agent-harness}{awesome-agent-harness} & \hyperref[c:agents-1]{Harness engineering}\,{\color{muted}\scriptsize p.\,\pageref{c:agents-1}}\\
2026 & \href{https://github.com/ai-boost/awesome-harness-engineering}{awesome-harness-engineering} & \hyperref[c:agents-0]{Agent harness}\,{\color{muted}\scriptsize p.\,\pageref{c:agents-0}}\\
2026 & \href{https://github.com/openclaw/openclaw}{OpenClaw} & \hyperref[c:agents-10]{Personal agents}\,{\color{muted}\scriptsize p.\,\pageref{c:agents-10}}\\
2025 & \href{https://github.com/humanlayer/12-factor-agents}{12-factor agents} & \hyperref[c:agents-4]{Workflows vs. agents}\,{\color{muted}\scriptsize p.\,\pageref{c:agents-4}}\\
2025 & \href{https://github.com/a2aproject/A2A}{A2A repository} & \hyperref[c:protocols-3]{Agent2Agent protocol}\,{\color{muted}\scriptsize p.\,\pageref{c:protocols-3}}\\
2025 & \href{https://github.com/ag-ui-protocol/ag-ui}{ag-ui} & \hyperref[c:protocols-8]{AG-UI}\,{\color{muted}\scriptsize p.\,\pageref{c:protocols-8}}\\
2025 & \href{https://github.com/zed-industries/agent-client-protocol}{agent-client-protocol} & \hyperref[c:protocols-7]{Agent Client Protocol}\,{\color{muted}\scriptsize p.\,\pageref{c:protocols-7}}\\
2025 & \href{https://github.com/agentic-commerce-protocol/agentic-commerce-protocol}{Agentic Commerce Protocol} & \hyperref[c:protocols-9]{Agentic commerce protocols}\,{\color{muted}\scriptsize p.\,\pageref{c:protocols-9}}\\
2025 & \href{https://github.com/agentsmd/agents.md}{agentsmd/agents.md} & \hyperref[c:protocols-6]{AGENTS.md}\,{\color{muted}\scriptsize p.\,\pageref{c:protocols-6}}\\
2025 & \href{https://github.com/anthropic-experimental/sandbox-runtime}{Anthropic sandbox-runtime} & \hyperref[c:agents-11]{Agent sandboxes}\,{\color{muted}\scriptsize p.\,\pageref{c:agents-11}}\\
2025 & \href{https://github.com/anthropics/skills}{anthropics/skills} & \hyperref[c:protocols-5]{Agent Skills}\,{\color{muted}\scriptsize p.\,\pageref{c:protocols-5}}\\
2025 & \href{https://github.com/google-agentic-commerce/AP2}{AP2} & \hyperref[c:protocols-9]{Agentic commerce protocols}\,{\color{muted}\scriptsize p.\,\pageref{c:protocols-9}}\\
2025 & \href{https://github.com/safety-research/circuit-tracer}{circuit-tracer} & \hyperref[c:safety-6]{Mechanistic interpretability}\,{\color{muted}\scriptsize p.\,\pageref{c:safety-6}}\\
2025 & \href{https://github.com/anthropics/claude-agent-sdk-python}{Claude Agent SDK (Python)} & \hyperref[c:agents-12]{Agent frameworks and SDKs}\,{\color{muted}\scriptsize p.\,\pageref{c:agents-12}}\\
2025 & \href{https://github.com/anthropics/claude-code}{Claude Code} & \hyperref[c:coding-0]{Coding agents}\,{\color{muted}\scriptsize p.\,\pageref{c:coding-0}}\\
2025 & \href{https://github.com/dagger/container-use}{container-use} & \hyperref[c:agents-11]{Agent sandboxes}\,{\color{muted}\scriptsize p.\,\pageref{c:agents-11}}\\
2025 & \href{https://github.com/deepseek-ai/DeepEP}{DeepEP} & \hyperref[c:arch-0]{Mixture of Experts}\,{\color{muted}\scriptsize p.\,\pageref{c:arch-0}}\\
2025 & \href{https://github.com/deepseek-ai/DeepSeek-OCR}{DeepSeek-OCR} & \hyperref[c:vlm-4]{Document AI and VLM-based OCR}\,{\color{muted}\scriptsize p.\,\pageref{c:vlm-4}}\\
2025 & \href{https://github.com/deepseek-ai/DeepSeek-V3.2-Exp}{DeepSeek-V3.2-Exp (DSA)} & \hyperref[c:arch-2]{Sparse and compressed attention}\,{\color{muted}\scriptsize p.\,\pageref{c:arch-2}}\\
2025 & \href{https://github.com/google-gemini/gemini-cli}{Gemini CLI} & \hyperref[c:coding-0]{Coding agents}\,{\color{muted}\scriptsize p.\,\pageref{c:coding-0}}\\
2025 & \href{https://github.com/gepa-ai/gepa}{GEPA} & \hyperref[c:practice-3]{Prompt engineering and prompt optimization}\,{\color{muted}\scriptsize p.\,\pageref{c:practice-3}}\\
2025 & \href{https://github.com/github/spec-kit}{GitHub Spec Kit} & \hyperref[c:coding-3]{Spec-driven development}\,{\color{muted}\scriptsize p.\,\pageref{c:coding-3}}\\
2025 & \href{https://github.com/google/adk-python}{Google ADK} & \hyperref[c:agents-12]{Agent frameworks and SDKs}\,{\color{muted}\scriptsize p.\,\pageref{c:agents-12}}\\
2025 & \href{https://github.com/block/goose}{goose} & \hyperref[c:protocols-4]{Agentic AI Foundation}\,{\color{muted}\scriptsize p.\,\pageref{c:protocols-4}}\\
2025 & \href{https://github.com/sapientinc/HRM}{Hierarchical Reasoning Model} & \hyperref[c:arch-10]{Tiny recursive reasoning models}\,{\color{muted}\scriptsize p.\,\pageref{c:arch-10}}\\
2025 & \href{https://github.com/InferenceMAX/InferenceMAX}{InferenceMAX} & \hyperref[c:inference-8]{Inference benchmarking}\,{\color{muted}\scriptsize p.\,\pageref{c:inference-8}}\\
2025 & \href{https://github.com/MoonshotAI/Kimi-K2}{Kimi K2} & \hyperref[c:training-2]{Muon and modern optimizers}\,{\color{muted}\scriptsize p.\,\pageref{c:training-2}}\\
2025 & \href{https://github.com/MoonshotAI/Kimi-Linear}{Kimi Linear} & \hyperref[c:arch-4]{Hybrid linear-attention and state space models}\,{\color{muted}\scriptsize p.\,\pageref{c:arch-4}}\\
2025 & \href{https://github.com/langchain-ai/deepagents}{LangChain deepagents} & \hyperref[c:agents-12]{Agent frameworks and SDKs}\,{\color{muted}\scriptsize p.\,\pageref{c:agents-12}}\\
2025 & \href{https://github.com/ML-GSAI/LLaDA}{LLaDA} & \hyperref[c:arch-5]{Diffusion language models}\,{\color{muted}\scriptsize p.\,\pageref{c:arch-5}}\\
2025 & \href{https://github.com/llm-d/llm-d}{llm-d} & \hyperref[c:inference-3]{Prefill/decode disaggregation}\,{\color{muted}\scriptsize p.\,\pageref{c:inference-3}}\\
2025 & \href{https://github.com/modelcontextprotocol/ext-apps}{MCP ext-apps} & \hyperref[c:protocols-1]{MCP Apps}\,{\color{muted}\scriptsize p.\,\pageref{c:protocols-1}}\\
2025 & \href{https://github.com/microsoft/agent-framework}{Microsoft Agent Framework} & \hyperref[c:agents-12]{Agent frameworks and SDKs}\,{\color{muted}\scriptsize p.\,\pageref{c:agents-12}}\\
2025 & \href{https://github.com/SWE-agent/mini-swe-agent}{mini-swe-agent} & \hyperref[c:agents-3]{Agentic loop / ReAct}\,{\color{muted}\scriptsize p.\,\pageref{c:agents-3}}\\
2025 & \href{https://github.com/karpathy/nanochat}{nanochat} & \hyperref[c:training-9]{Build-it-yourself references}\,{\color{muted}\scriptsize p.\,\pageref{c:training-9}}\\
2025 & \href{https://github.com/nvidia-cosmos}{NVIDIA Cosmos} & \hyperref[c:embodied-2]{World models}\,{\color{muted}\scriptsize p.\,\pageref{c:embodied-2}}\\
2025 & \href{https://github.com/ai-dynamo/dynamo}{NVIDIA Dynamo} & \hyperref[c:inference-3]{Prefill/decode disaggregation}\,{\color{muted}\scriptsize p.\,\pageref{c:inference-3}}\\
2025 & \href{https://github.com/NVIDIA/Isaac-GR00T}{NVIDIA Isaac GR00T} & \hyperref[c:embodied-0]{Vision-language-action models}\,{\color{muted}\scriptsize p.\,\pageref{c:embodied-0}}\\
2025 & \href{https://github.com/allenai/olmocr}{olmOCR} & \hyperref[c:vlm-4]{Document AI and VLM-based OCR}\,{\color{muted}\scriptsize p.\,\pageref{c:vlm-4}}\\
2025 & \href{https://github.com/openai/openai-agents-python}{OpenAI Agents SDK} & \hyperref[c:agents-12]{Agent frameworks and SDKs}\,{\color{muted}\scriptsize p.\,\pageref{c:agents-12}}\\
2025 & \href{https://github.com/openai/codex}{OpenAI Codex CLI} & \hyperref[c:coding-0]{Coding agents}\,{\color{muted}\scriptsize p.\,\pageref{c:coding-0}}\\
2025 & \href{https://github.com/openai/gpt-oss}{openai/gpt-oss (MXFP4)} & \hyperref[c:quant-0]{Microscaling formats}\,{\color{muted}\scriptsize p.\,\pageref{c:quant-0}}\\
2025 & \href{https://github.com/sst/opencode}{OpenCode} & \hyperref[c:coding-0]{Coding agents}\,{\color{muted}\scriptsize p.\,\pageref{c:coding-0}}\\
2025 & \href{https://github.com/meta-pytorch/OpenEnv}{OpenEnv} & \hyperref[c:reasoning-5]{RL environments and verifiers}\,{\color{muted}\scriptsize p.\,\pageref{c:reasoning-5}}\\
2025 & \href{https://github.com/codelion/openevolve}{OpenEvolve} & \hyperref[c:embodied-7]{AI for science and algorithm discovery}\,{\color{muted}\scriptsize p.\,\pageref{c:embodied-7}}\\
2025 & \href{https://github.com/Physical-Intelligence/openpi}{openpi (Physical Intelligence)} & \hyperref[c:embodied-0]{Vision-language-action models}\,{\color{muted}\scriptsize p.\,\pageref{c:embodied-0}}\\
2025 & \href{https://github.com/microsoft/playwright-mcp}{Playwright MCP} & \hyperref[c:agents-9]{Computer use and browser agents}\,{\color{muted}\scriptsize p.\,\pageref{c:agents-9}}\\
2025 & \href{https://github.com/PrimeIntellect-ai/verifiers}{Prime Intellect verifiers} & \hyperref[c:reasoning-5]{RL environments and verifiers}\,{\color{muted}\scriptsize p.\,\pageref{c:reasoning-5}}\\
2025 & \href{https://github.com/PrimeIntellect-ai/prime-rl}{prime-rl} & \hyperref[c:training-8]{Decentralized training}\,{\color{muted}\scriptsize p.\,\pageref{c:training-8}}\\
2025 & \href{https://github.com/QwenLM/Qwen3-Omni}{Qwen3-Omni} & \hyperref[c:vlm-7]{Omni and speech-to-speech models}\,{\color{muted}\scriptsize p.\,\pageref{c:vlm-7}}\\
2025 & \href{https://github.com/QwenLM/Qwen3-VL}{Qwen3-VL} & \hyperref[c:vlm-1]{Open VLM families}\,{\color{muted}\scriptsize p.\,\pageref{c:vlm-1}}\\
2025 & \href{https://github.com/NovaSky-AI/SkyRL}{SkyRL} & \hyperref[c:reasoning-5]{RL environments and verifiers}\,{\color{muted}\scriptsize p.\,\pageref{c:reasoning-5}}\\
2025 & \href{https://github.com/sierra-research/tau2-bench}{tau2-bench} & \hyperref[c:evals-5]{Agent and tool-use benchmarks}\,{\color{muted}\scriptsize p.\,\pageref{c:evals-5}}\\
2025 & \href{https://github.com/tile-ai/tilelang}{TileLang} & \hyperref[c:hardware-13]{Kernel languages and compilers}\,{\color{muted}\scriptsize p.\,\pageref{c:hardware-13}}\\
2025 & \href{https://github.com/SamsungSAILMontreal/TinyRecursiveModels}{Tiny Recursive Models} & \hyperref[c:arch-10]{Tiny recursive reasoning models}\,{\color{muted}\scriptsize p.\,\pageref{c:arch-10}}\\
2025 & \href{https://github.com/bytedance/UI-TARS}{UI-TARS} & \hyperref[c:vlm-6]{GUI grounding agents}\,{\color{muted}\scriptsize p.\,\pageref{c:vlm-6}}\\
2025 & \href{https://github.com/facebookresearch/vjepa2}{V-JEPA 2} & \hyperref[c:embodied-3]{JEPA}\,{\color{muted}\scriptsize p.\,\pageref{c:embodied-3}}\\
2025 & \href{https://github.com/Wan-Video/Wan2.2}{Wan 2.2} & \hyperref[c:embodied-6]{Video generation}\,{\color{muted}\scriptsize p.\,\pageref{c:embodied-6}}\\
2025 & \href{https://github.com/coinbase/x402}{x402} & \hyperref[c:protocols-9]{Agentic commerce protocols}\,{\color{muted}\scriptsize p.\,\pageref{c:protocols-9}}\\
2024 & \href{https://github.com/google-deepmind/alphafold3}{AlphaFold 3} & \hyperref[c:embodied-7]{AI for science and algorithm discovery}\,{\color{muted}\scriptsize p.\,\pageref{c:embodied-7}}\\
2024 & \href{https://github.com/browser-use/browser-use}{browser-use} & \hyperref[c:agents-9]{Computer use and browser agents}\,{\color{muted}\scriptsize p.\,\pageref{c:agents-9}}\\
2024 & \href{https://github.com/facebookresearch/blt}{Byte Latent Transformer} & \hyperref[c:arch-9]{Tokenizer-free and byte-level models}\,{\color{muted}\scriptsize p.\,\pageref{c:arch-9}}\\
2024 & \href{https://github.com/chiphuyen/aie-book}{Chip Huyen: AI Engineering book resources} & \hyperref[c:practice-1]{AI Engineer}\,{\color{muted}\scriptsize p.\,\pageref{c:practice-1}}\\
2024 & \href{https://github.com/illuin-tech/colpali}{ColPali} & \hyperref[c:vlm-5]{Visual document retrieval}\,{\color{muted}\scriptsize p.\,\pageref{c:vlm-5}}\\
2024 & \href{https://github.com/daytonaio/daytona}{Daytona} & \hyperref[c:agents-11]{Agent sandboxes}\,{\color{muted}\scriptsize p.\,\pageref{c:agents-11}}\\
2024 & \href{https://github.com/docling-project/docling}{Docling} & \hyperref[c:vlm-4]{Document AI and VLM-based OCR}\,{\color{muted}\scriptsize p.\,\pageref{c:vlm-4}}\\
2024 & \href{https://github.com/SafeAILab/EAGLE}{EAGLE} & \hyperref[c:arch-6]{Multi-token prediction and speculative decoding}\,{\color{muted}\scriptsize p.\,\pageref{c:arch-6}}\\
2024 & \href{https://github.com/black-forest-labs/flux}{FLUX} & \hyperref[c:embodied-5]{Diffusion transformers and flow matching}\,{\color{muted}\scriptsize p.\,\pageref{c:embodied-5}}\\
2024 & \href{https://github.com/NVlabs/GatedDeltaNet}{Gated DeltaNet} & \hyperref[c:arch-4]{Hybrid linear-attention and state space models}\,{\color{muted}\scriptsize p.\,\pageref{c:arch-4}}\\
2024 & \href{https://github.com/getzep/graphiti}{Graphiti} & \hyperref[c:context-4]{Agent memory}\,{\color{muted}\scriptsize p.\,\pageref{c:context-4}}\\
2024 & \href{https://github.com/UKGovernmentBEIS/inspect_ai}{Inspect AI (UK AISI)} & \hyperref[c:evals-0]{Eval harnesses}\,{\color{muted}\scriptsize p.\,\pageref{c:evals-0}}\\
2024 & \href{https://github.com/UKGovernmentBEIS/inspect_evals}{inspect\_evals} & \hyperref[c:evals-0]{Eval harnesses}\,{\color{muted}\scriptsize p.\,\pageref{c:evals-0}}\\
2024 & \href{https://github.com/isaac-sim/IsaacLab}{Isaac Lab} & \hyperref[c:embodied-4]{Robot simulation and sim-to-real}\,{\color{muted}\scriptsize p.\,\pageref{c:embodied-4}}\\
2024 & \href{https://github.com/langchain-ai/langgraph}{LangGraph} & \hyperref[c:agents-7]{Multi-agent systems and subagents}\,{\color{muted}\scriptsize p.\,\pageref{c:agents-7}}\\
2024 & \href{https://github.com/huggingface/lerobot}{LeRobot} & \hyperref[c:embodied-0]{Vision-language-action models}\,{\color{muted}\scriptsize p.\,\pageref{c:embodied-0}}\\
2024 & \href{https://github.com/huggingface/lighteval}{lighteval} & \hyperref[c:evals-0]{Eval harnesses}\,{\color{muted}\scriptsize p.\,\pageref{c:evals-0}}\\
2024 & \href{https://github.com/HKUDS/LightRAG}{LightRAG} & \hyperref[c:context-7]{GraphRAG}\,{\color{muted}\scriptsize p.\,\pageref{c:context-7}}\\
2024 & \href{https://github.com/livekit/agents}{LiveKit Agents} & \hyperref[c:vlm-7]{Omni and speech-to-speech models}\,{\color{muted}\scriptsize p.\,\pageref{c:vlm-7}}\\
2024 & \href{https://github.com/vllm-project/llm-compressor}{llm-compressor} & \hyperref[c:quant-3]{Post-training quantization}\,{\color{muted}\scriptsize p.\,\pageref{c:quant-3}}\\
2024 & \href{https://github.com/LMCache/LMCache}{LMCache} & \hyperref[c:inference-4]{KV cache offloading}\,{\color{muted}\scriptsize p.\,\pageref{c:inference-4}}\\
2024 & \href{https://github.com/EvolvingLMMs-Lab/lmms-eval}{lmms-eval} & \hyperref[c:vlm-9]{Multimodal evaluation}\,{\color{muted}\scriptsize p.\,\pageref{c:vlm-9}}\\
2024 & \href{https://github.com/modelcontextprotocol}{MCP org (SDKs, servers)} & \hyperref[c:protocols-0]{Model Context Protocol}\,{\color{muted}\scriptsize p.\,\pageref{c:protocols-0}}\\
2024 & \href{https://github.com/modelcontextprotocol/servers}{MCP servers} & \hyperref[c:learn-3]{Curated lists}\,{\color{muted}\scriptsize p.\,\pageref{c:learn-3}}\\
2024 & \href{https://github.com/mem0ai/mem0}{mem0} & \hyperref[c:context-4]{Agent memory}\,{\color{muted}\scriptsize p.\,\pageref{c:context-4}}\\
2024 & \href{https://github.com/microsoft/BitNet}{Microsoft BitNet} & \hyperref[c:quant-6]{Ternary and 1-bit models}\,{\color{muted}\scriptsize p.\,\pageref{c:quant-6}}\\
2024 & \href{https://github.com/microsoft/graphrag}{Microsoft GraphRAG} & \hyperref[c:context-7]{GraphRAG}\,{\color{muted}\scriptsize p.\,\pageref{c:context-7}}\\
2024 & \href{https://github.com/KellerJordan/modded-nanogpt}{modded-nanogpt} & \hyperref[c:training-2]{Muon and modern optimizers}\,{\color{muted}\scriptsize p.\,\pageref{c:training-2}}\\
2024 & \href{https://github.com/kvcache-ai/Mooncake}{Mooncake} & \hyperref[c:inference-3]{Prefill/decode disaggregation}\,{\color{muted}\scriptsize p.\,\pageref{c:inference-3}}\\
2024 & \href{https://github.com/kyutai-labs/moshi}{Moshi} & \hyperref[c:vlm-7]{Omni and speech-to-speech models}\,{\color{muted}\scriptsize p.\,\pageref{c:vlm-7}}\\
2024 & \href{https://github.com/microsoft/OmniParser}{OmniParser} & \hyperref[c:vlm-6]{GUI grounding agents}\,{\color{muted}\scriptsize p.\,\pageref{c:vlm-6}}\\
2024 & \href{https://github.com/OpenHands/OpenHands}{OpenHands} & \hyperref[c:coding-0]{Coding agents}\,{\color{muted}\scriptsize p.\,\pageref{c:coding-0}}\\
2024 & \href{https://github.com/openvla/openvla}{OpenVLA} & \hyperref[c:embodied-0]{Vision-language-action models}\,{\color{muted}\scriptsize p.\,\pageref{c:embodied-0}}\\
2024 & \href{https://github.com/pipecat-ai/pipecat}{Pipecat} & \hyperref[c:vlm-7]{Omni and speech-to-speech models}\,{\color{muted}\scriptsize p.\,\pageref{c:vlm-7}}\\
2024 & \href{https://github.com/pydantic/pydantic-ai}{Pydantic AI} & \hyperref[c:agents-12]{Agent frameworks and SDKs}\,{\color{muted}\scriptsize p.\,\pageref{c:agents-12}}\\
2024 & \href{https://github.com/sgl-project/sglang}{SGLang} & \hyperref[c:inference-0]{Inference engines}\,{\color{muted}\scriptsize p.\,\pageref{c:inference-0}}\\
2024 & \href{https://github.com/huggingface/smolagents}{smolagents} & \hyperref[c:agents-12]{Agent frameworks and SDKs}\,{\color{muted}\scriptsize p.\,\pageref{c:agents-12}}\\
2024 & \href{https://github.com/HazyResearch/ThunderKittens}{ThunderKittens} & \hyperref[c:hardware-13]{Kernel languages and compilers}\,{\color{muted}\scriptsize p.\,\pageref{c:hardware-13}}\\
2024 & \href{https://github.com/pytorch/torchtitan}{torchtitan} & \hyperref[c:training-0]{Parallelism strategies}\,{\color{muted}\scriptsize p.\,\pageref{c:training-0}}\\
2024 & \href{https://github.com/volcengine/verl}{verl} & \hyperref[c:reasoning-5]{RL environments and verifiers}\,{\color{muted}\scriptsize p.\,\pageref{c:reasoning-5}}\\
2024 & \href{https://github.com/mlc-ai/xgrammar}{XGrammar} & \hyperref[c:context-9]{Structured outputs and constrained decoding}\,{\color{muted}\scriptsize p.\,\pageref{c:context-9}}\\
2023 & \href{https://github.com/Aider-AI/aider}{Aider} & \hyperref[c:coding-0]{Coding agents}\,{\color{muted}\scriptsize p.\,\pageref{c:coding-0}}\\
2023 & \href{https://github.com/Arize-ai/phoenix}{Arize Phoenix} & \hyperref[c:practice-5]{LLMOps and observability}\,{\color{muted}\scriptsize p.\,\pageref{c:practice-5}}\\
2023 & \href{https://github.com/mit-han-lab/llm-awq}{AWQ} & \hyperref[c:quant-3]{Post-training quantization}\,{\color{muted}\scriptsize p.\,\pageref{c:quant-3}}\\
2023 & \href{https://github.com/axolotl-ai-cloud/axolotl}{Axolotl} & \hyperref[c:training-6]{Fine-tuning toolkits and APIs}\,{\color{muted}\scriptsize p.\,\pageref{c:training-6}}\\
2023 & \href{https://github.com/comfyanonymous/ComfyUI}{ComfyUI} & \hyperref[c:embodied-5]{Diffusion transformers and flow matching}\,{\color{muted}\scriptsize p.\,\pageref{c:embodied-5}}\\
2023 & \href{https://github.com/crewAIInc/crewAI}{CrewAI} & \hyperref[c:agents-12]{Agent frameworks and SDKs}\,{\color{muted}\scriptsize p.\,\pageref{c:agents-12}}\\
2023 & \href{https://github.com/stanfordnlp/dspy}{DSPy} & \hyperref[c:practice-3]{Prompt engineering and prompt optimization}\,{\color{muted}\scriptsize p.\,\pageref{c:practice-3}}\\
2023 & \href{https://github.com/e2b-dev/E2B}{E2B} & \hyperref[c:agents-11]{Agent sandboxes}\,{\color{muted}\scriptsize p.\,\pageref{c:agents-11}}\\
2023 & \href{https://github.com/pytorch/executorch}{ExecuTorch} & \hyperref[c:inference-6]{On-device and edge inference}\,{\color{muted}\scriptsize p.\,\pageref{c:inference-6}}\\
2023 & \href{https://github.com/fla-org/flash-linear-attention}{flash-linear-attention} & \hyperref[c:arch-4]{Hybrid linear-attention and state space models}\,{\color{muted}\scriptsize p.\,\pageref{c:arch-4}}\\
2023 & \href{https://github.com/flashinfer-ai/flashinfer}{FlashInfer} & \hyperref[c:inference-5]{Attention kernels}\,{\color{muted}\scriptsize p.\,\pageref{c:inference-5}}\\
2023 & \href{https://github.com/idavidrein/gpqa}{GPQA} & \hyperref[c:evals-6]{Frontier knowledge benchmarks}\,{\color{muted}\scriptsize p.\,\pageref{c:evals-6}}\\
2023 & \href{https://github.com/assafelovic/gpt-researcher}{GPT Researcher} & \hyperref[c:context-6]{Agentic search and deep research}\,{\color{muted}\scriptsize p.\,\pageref{c:context-6}}\\
2023 & \href{https://github.com/guardrails-ai/guardrails}{Guardrails AI} & \hyperref[c:safety-5]{Guardrails and classifiers}\,{\color{muted}\scriptsize p.\,\pageref{c:safety-5}}\\
2023 & \href{https://github.com/huggingface/peft}{HF PEFT} & \hyperref[c:training-5]{LoRA and parameter-efficient fine-tuning}\,{\color{muted}\scriptsize p.\,\pageref{c:training-5}}\\
2023 & \href{https://github.com/OpenGVLab/InternVL}{InternVL} & \hyperref[c:vlm-1]{Open VLM families}\,{\color{muted}\scriptsize p.\,\pageref{c:vlm-1}}\\
2023 & \href{https://github.com/langfuse/langfuse}{Langfuse} & \hyperref[c:practice-5]{LLMOps and observability}\,{\color{muted}\scriptsize p.\,\pageref{c:practice-5}}\\
2023 & \href{https://github.com/letta-ai/letta}{Letta (MemGPT)} & \hyperref[c:context-4]{Agent memory}\,{\color{muted}\scriptsize p.\,\pageref{c:context-4}}\\
2023 & \href{https://github.com/BerriAI/litellm}{LiteLLM} & \hyperref[c:inference-7]{LLM gateways and routers}\,{\color{muted}\scriptsize p.\,\pageref{c:inference-7}}\\
2023 & \href{https://github.com/hiyouga/LLaMA-Factory}{LLaMA-Factory} & \hyperref[c:training-6]{Fine-tuning toolkits and APIs}\,{\color{muted}\scriptsize p.\,\pageref{c:training-6}}\\
2023 & \href{https://github.com/ggml-org/llama.cpp}{llama.cpp} & \hyperref[c:inference-0]{Inference engines}\,{\color{muted}\scriptsize p.\,\pageref{c:inference-0}}\\
2023 & \href{https://github.com/haotian-liu/LLaVA}{LLaVA} & \hyperref[c:vlm-0]{VLM architecture}\,{\color{muted}\scriptsize p.\,\pageref{c:vlm-0}}\\
2023 & \href{https://github.com/state-spaces/mamba}{Mamba} & \hyperref[c:arch-4]{Hybrid linear-attention and state space models}\,{\color{muted}\scriptsize p.\,\pageref{c:arch-4}}\\
2023 & \href{https://github.com/arcee-ai/mergekit}{mergekit} & \hyperref[c:training-7]{Model merging}\,{\color{muted}\scriptsize p.\,\pageref{c:training-7}}\\
2023 & \href{https://github.com/ml-explore/mlx}{MLX} & \hyperref[c:inference-0]{Inference engines}\,{\color{muted}\scriptsize p.\,\pageref{c:inference-0}}\\
2023 & \href{https://github.com/modular/modular}{Modular (Mojo, MAX)} & \hyperref[c:hardware-13]{Kernel languages and compilers}\,{\color{muted}\scriptsize p.\,\pageref{c:hardware-13}}\\
2023 & \href{https://github.com/NVIDIA/NeMo-Guardrails}{NeMo Guardrails} & \hyperref[c:safety-5]{Guardrails and classifiers}\,{\color{muted}\scriptsize p.\,\pageref{c:safety-5}}\\
2023 & \href{https://github.com/ndif-team/nnsight}{nnsight} & \hyperref[c:safety-6]{Mechanistic interpretability}\,{\color{muted}\scriptsize p.\,\pageref{c:safety-6}}\\
2023 & \href{https://github.com/ollama/ollama}{Ollama} & \hyperref[c:inference-0]{Inference engines}\,{\color{muted}\scriptsize p.\,\pageref{c:inference-0}}\\
2023 & \href{https://github.com/dottxt-ai/outlines}{Outlines} & \hyperref[c:context-9]{Structured outputs and constrained decoding}\,{\color{muted}\scriptsize p.\,\pageref{c:context-9}}\\
2023 & \href{https://github.com/promptfoo/promptfoo}{promptfoo} & \hyperref[c:evals-0]{Eval harnesses}\,{\color{muted}\scriptsize p.\,\pageref{c:evals-0}}\\
2023 & \href{https://github.com/NVIDIA/TensorRT-LLM}{TensorRT-LLM} & \hyperref[c:inference-0]{Inference engines}\,{\color{muted}\scriptsize p.\,\pageref{c:inference-0}}\\
2023 & \href{https://github.com/pytorch/ao}{torchao} & \hyperref[c:quant-4]{Quantization-aware training}\,{\color{muted}\scriptsize p.\,\pageref{c:quant-4}}\\
2023 & \href{https://github.com/unslothai/unsloth}{Unsloth} & \hyperref[c:training-6]{Fine-tuning toolkits and APIs}\,{\color{muted}\scriptsize p.\,\pageref{c:training-6}}\\
2023 & \href{https://github.com/vllm-project/vllm}{vLLM} & \hyperref[c:inference-0]{Inference engines}\,{\color{muted}\scriptsize p.\,\pageref{c:inference-0}}\\
2023 & \href{https://github.com/open-compass/VLMEvalKit}{VLMEvalKit} & \hyperref[c:vlm-9]{Multimodal evaluation}\,{\color{muted}\scriptsize p.\,\pageref{c:vlm-9}}\\
2023 & \href{https://github.com/jquesnelle/yarn}{YaRN} & \hyperref[c:arch-7]{Long-context methods}\,{\color{muted}\scriptsize p.\,\pageref{c:arch-7}}\\
2022 & \href{https://github.com/Dao-AILab/flash-attention}{FlashAttention} & \hyperref[c:inference-5]{Attention kernels}\,{\color{muted}\scriptsize p.\,\pageref{c:inference-5}}\\
2022 & \href{https://github.com/run-llama/llama_index}{LlamaIndex} & \hyperref[c:context-5]{Retrieval-augmented generation}\,{\color{muted}\scriptsize p.\,\pageref{c:context-5}}\\
2022 & \href{https://github.com/openxla/xla}{OpenXLA} & \hyperref[c:hardware-13]{Kernel languages and compilers}\,{\color{muted}\scriptsize p.\,\pageref{c:hardware-13}}\\
2022 & \href{https://github.com/huggingface/safetensors}{safetensors} & \hyperref[c:hardware-15]{Model interchange formats}\,{\color{muted}\scriptsize p.\,\pageref{c:hardware-15}}\\
2022 & \href{https://github.com/TransformerLensOrg/TransformerLens}{TransformerLens} & \hyperref[c:safety-6]{Mechanistic interpretability}\,{\color{muted}\scriptsize p.\,\pageref{c:safety-6}}\\
2022 & \href{https://github.com/openai/whisper}{Whisper} & \hyperref[c:vlm-7]{Omni and speech-to-speech models}\,{\color{muted}\scriptsize p.\,\pageref{c:vlm-7}}\\
2021 & \href{https://github.com/bitsandbytes-foundation/bitsandbytes}{bitsandbytes} & \hyperref[c:quant-5]{GGUF and local model formats}\,{\color{muted}\scriptsize p.\,\pageref{c:quant-5}}\\
2021 & \href{https://github.com/EleutherAI/lm-evaluation-harness}{lm-evaluation-harness} & \hyperref[c:evals-0]{Eval harnesses}\,{\color{muted}\scriptsize p.\,\pageref{c:evals-0}}\\
2021 & \href{https://github.com/google-deepmind/mujoco}{MuJoCo} & \hyperref[c:embodied-4]{Robot simulation and sim-to-real}\,{\color{muted}\scriptsize p.\,\pageref{c:embodied-4}}\\
2020 & \href{https://github.com/stanford-futuredata/ColBERT}{ColBERT} & \hyperref[c:context-8]{Embeddings, late interaction and Matryoshka}\,{\color{muted}\scriptsize p.\,\pageref{c:context-8}}\\
2020 & \href{https://github.com/deepspeedai/DeepSpeed}{DeepSpeed} & \hyperref[c:training-0]{Parallelism strategies}\,{\color{muted}\scriptsize p.\,\pageref{c:training-0}}\\
2020 & \href{https://github.com/huggingface/trl}{Hugging Face TRL} & \hyperref[c:reasoning-6]{RLHF, DPO and preference optimization}\,{\color{muted}\scriptsize p.\,\pageref{c:reasoning-6}}\\
2020 & \href{https://github.com/PaddlePaddle/PaddleOCR}{PaddleOCR} & \hyperref[c:vlm-4]{Document AI and VLM-based OCR}\,{\color{muted}\scriptsize p.\,\pageref{c:vlm-4}}\\
2019 & \href{https://github.com/NVIDIA/Megatron-LM}{Megatron-LM} & \hyperref[c:training-0]{Parallelism strategies}\,{\color{muted}\scriptsize p.\,\pageref{c:training-0}}\\
2019 & \href{https://github.com/triton-lang/triton}{Triton} & \hyperref[c:hardware-13]{Kernel languages and compilers}\,{\color{muted}\scriptsize p.\,\pageref{c:hardware-13}}\\
2018 & \href{https://github.com/microsoft/onnxruntime}{ONNX Runtime} & \hyperref[c:inference-6]{On-device and edge inference}\,{\color{muted}\scriptsize p.\,\pageref{c:inference-6}}\\
2017 & \href{https://github.com/NVIDIA/cutlass}{CUTLASS / CuTe} & \hyperref[c:hardware-13]{Kernel languages and compilers}\,{\color{muted}\scriptsize p.\,\pageref{c:hardware-13}}\\
2017 & \href{https://github.com/onnx/onnx}{ONNX} & \hyperref[c:hardware-15]{Model interchange formats}\,{\color{muted}\scriptsize p.\,\pageref{c:hardware-15}}\\
2016 & \href{https://github.com/NVIDIA/nccl}{NCCL} & \hyperref[c:hardware-14]{Collective communication}\,{\color{muted}\scriptsize p.\,\pageref{c:hardware-14}}\\
2016 & \href{https://github.com/ROCm/ROCm}{ROCm} & \hyperref[c:hardware-13]{Kernel languages and compilers}\,{\color{muted}\scriptsize p.\,\pageref{c:hardware-13}}\\
\end{longtable}}
\section{Blog (103)}
{\small\begin{longtable}{@{}p{0.07\linewidth}>{\raggedright}p{0.55\linewidth}>{\raggedright\arraybackslash}p{0.34\linewidth}@{}}
\textbf{\sffamily Year} & \textbf{\sffamily Resource} & \textbf{\sffamily Concept}\\ \hline\endhead
2026 & \href{https://650group.com/blog/in-the-ai-era-ethernet-set-to-surge-in-scale-out-and-ramp-in-scale-up/}{650 Group: Ethernet in scale-up} & \hyperref[c:hardware-4]{ESUN and SUE}\,{\color{muted}\scriptsize p.\,\pageref{c:hardware-4}}\\
2026 & \href{https://aishwaryasrinivasan.substack.com/p/the-hottest-role-in-2026-forward}{Aishwarya Srinivasan: The hottest role in 2026} & \hyperref[c:practice-0]{Forward Deployed Engineer}\,{\color{muted}\scriptsize p.\,\pageref{c:practice-0}}\\
2026 & \href{https://alexeyondata.substack.com/p/what-ai-forward-deployed-engineers}{Alexey Grigorev: What AI FDEs do} & \hyperref[c:practice-0]{Forward Deployed Engineer}\,{\color{muted}\scriptsize p.\,\pageref{c:practice-0}}\\
2026 & \href{https://anthropic.com/engineering/managed-agents}{Anthropic: Scaling Managed Agents} & \hyperref[c:agents-2]{Managed agents: decoupling brain, hands and session}\,{\color{muted}\scriptsize p.\,\pageref{c:agents-2}}\\
2026 & \href{https://arcprize.org/blog/arc-agi-3-launch}{ARC Prize: Announcing ARC-AGI-3} & \hyperref[c:evals-1]{Harness sensitivity of benchmarks}\,{\color{muted}\scriptsize p.\,\pageref{c:evals-1}}\\
2026 & \href{https://www.bentoml.com/blog/multimodal-ai-a-guide-to-open-source-vision-language-models}{BentoML: Open VLM guide} & \hyperref[c:vlm-1]{Open VLM families}\,{\color{muted}\scriptsize p.\,\pageref{c:vlm-1}}\\
2026 & \href{https://www.csoonline.com/article/4218679/openai-launches-gpt-6-astra-its-first-model-to-cross-a-critical-cybersecurity-threshold.html}{CSO Online: Astra crosses the Critical cyber threshold} & \hyperref[c:landscape-2]{Tiered access for dual-use models}\,{\color{muted}\scriptsize p.\,\pageref{c:landscape-2}}\\
2026 & \href{https://www.datacamp.com/blog/arc-agi-3}{DataCamp: ARC-AGI-3 explained} & \hyperref[c:evals-4]{ARC-AGI-3 and interactive reasoning}\,{\color{muted}\scriptsize p.\,\pageref{c:evals-4}}\\
2026 & \href{https://datanorth.ai/news/anthropic-releases-claude-fable-5-and-mythos-5}{DataNorth: Claude Fable 5 and Mythos 5 explained} & \hyperref[c:landscape-0]{Mythos-class tier and trusted access}\,{\color{muted}\scriptsize p.\,\pageref{c:landscape-0}}\\
2026 & \href{https://www.datacenterdynamics.com/en/news/nvidia-announces-vera-rubin-superchip-for-late-2026/}{DCD: Vera Rubin superchip} & \hyperref[c:hardware-0]{Rack-scale systems}\,{\color{muted}\scriptsize p.\,\pageref{c:hardware-0}}\\
2026 & \href{https://www.faros.ai/blog/harness-engineering}{Faros: Harness engineering guide} & \hyperref[c:agents-1]{Harness engineering}\,{\color{muted}\scriptsize p.\,\pageref{c:agents-1}}\\
2026 & \href{https://gpuinsights.net/nvidia-rubin-vs-amd-helios-mi450-2026/}{GPU Insights: Rubin vs Helios} & \hyperref[c:hardware-1]{AMD Helios and Open Rack Wide}\,{\color{muted}\scriptsize p.\,\pageref{c:hardware-1}}\\
2026 & \href{https://huggingface.co/blog/ResterChed/deepseek-v4-ga-architecture}{HF blog: DeepSeek V4 architecture analysis} & \hyperref[c:landscape-3]{Open-weight frontier models}\,{\color{muted}\scriptsize p.\,\pageref{c:landscape-3}}\\
2026 & \href{https://huggingface.co/blog/sergiopaniego/distillation-2026}{Hugging Face: Distillation in 2026} & \hyperref[c:reasoning-3]{Multi-teacher on-policy distillation}\,{\color{muted}\scriptsize p.\,\pageref{c:reasoning-3}}\\
2026 & \href{https://ijonis.com/en/forward-deployed-engineer-explained}{IJONIS: FDE and services-as-software} & \hyperref[c:practice-2]{Services as software}\,{\color{muted}\scriptsize p.\,\pageref{c:practice-2}}\\
2026 & \href{https://www.infoq.com/news/2026/09/gpt-6-astra-critical-cyber/}{InfoQ: Astra is first model rated Critical for cyber} & \hyperref[c:landscape-1]{GPT-6 Astra and the Critical cyber threshold}\,{\color{muted}\scriptsize p.\,\pageref{c:landscape-1}}\\
2026 & \href{https://www.marktechpost.com/2026/05/20/what-is-a-forward-deployed-engineer-the-ai-role-openai-anthropic-and-google-are-hiring-in-2026/}{MarkTechPost: What is an FDE} & \hyperref[c:practice-0]{Forward Deployed Engineer}\,{\color{muted}\scriptsize p.\,\pageref{c:practice-0}}\\
2026 & \href{https://martinfowler.com/articles/harness-engineering.html}{Martin Fowler: Harness engineering for coding agent users} & \hyperref[c:agents-0]{Agent harness}\,{\color{muted}\scriptsize p.\,\pageref{c:agents-0}}\\
2026 & \href{https://blog.modelcontextprotocol.io/posts/2026-07-28/}{MCP blog: The 2026-07-28 specification} & \hyperref[c:protocols-0]{Model Context Protocol}\,{\color{muted}\scriptsize p.\,\pageref{c:protocols-0}}\\
2026 & \href{https://blog.modelcontextprotocol.io/posts/2026-07-28-release-candidate/}{MCP release candidate notes} & \hyperref[c:protocols-2]{MCP Tasks}\,{\color{muted}\scriptsize p.\,\pageref{c:protocols-2}}\\
2026 & \href{https://blog.modelcontextprotocol.io/posts/mcp-roadmap/}{MCP roadmap (Aug 2026)} & \hyperref[c:protocols-0]{Model Context Protocol}\,{\color{muted}\scriptsize p.\,\pageref{c:protocols-0}}\\
2026 & \href{https://metavert.io/world-models-for-robotics}{Metavert: World models for robotics} & \hyperref[c:embodied-1]{World action models}\,{\color{muted}\scriptsize p.\,\pageref{c:embodied-1}}\\
2026 & \href{https://www.modulos.ai/blog/eu-ai-act-omnibus-deal}{Modulos: The Omnibus deal} & \hyperref[c:safety-11]{EU AI Act and the Digital Omnibus}\,{\color{muted}\scriptsize p.\,\pageref{c:safety-11}}\\
2026 & \href{https://openai.com/index/harness-engineering/}{OpenAI: Harness engineering} & \hyperref[c:agents-1]{Harness engineering}\,{\color{muted}\scriptsize p.\,\pageref{c:agents-1}}\\
2026 & \href{https://openai.com/index/path-to-astra/}{OpenAI: Path to Astra} & \hyperref[c:landscape-1]{GPT-6 Astra and the Critical cyber threshold}\,{\color{muted}\scriptsize p.\,\pageref{c:landscape-1}}\\
2026 & \href{https://openai.com/index/safety-overview-gpt-6-astra/}{OpenAI: Safety overview of GPT-6 Astra} & \hyperref[c:landscape-1]{GPT-6 Astra and the Critical cyber threshold}\,{\color{muted}\scriptsize p.\,\pageref{c:landscape-1}}\\
2026 & \href{https://www.prolific.com/resources/rl-environments-for-agentic-ai-the-2026-guide-to-building-and-deploying-enterprise-rl-environments}{Prolific: RL environments for agentic AI (2026 guide)} & \hyperref[c:reasoning-5]{RL environments and verifiers}\,{\color{muted}\scriptsize p.\,\pageref{c:reasoning-5}}\\
2026 & \href{https://magazine.sebastianraschka.com/p/llm-research-papers-2026-part1}{Raschka: LLM research papers 2026} & \hyperref[c:arch-12]{Architecture overviews}\,{\color{muted}\scriptsize p.\,\pageref{c:arch-12}}\\
2026 & \href{https://www.servethehome.com/nvidia-launches-next-generation-rubin-ai-compute-platform-at-ces-2026/}{ServeTheHome: Rubin platform at CES 2026} & \hyperref[c:hardware-0]{Rack-scale systems}\,{\color{muted}\scriptsize p.\,\pageref{c:hardware-0}}\\
2026 & \href{https://www.synopsys.com/blogs/chip-design/esun-ethernet-ai-scale-up-networking.html}{Synopsys: ESUN explained} & \hyperref[c:hardware-4]{ESUN and SUE}\,{\color{muted}\scriptsize p.\,\pageref{c:hardware-4}}\\
2026 & \href{https://www.synopsys.com/articles/ethernet-standards-scale-up-ai.html}{Synopsys: Ethernet standards for scale-up AI} & \hyperref[c:hardware-2]{Scale-up vs. scale-out networking}\,{\color{muted}\scriptsize p.\,\pageref{c:hardware-2}}\\
2026 & \href{https://theaiengineer.substack.com/p/what-is-a-forward-deployed-engineer}{The AI Engineer: What is an FDE} & \hyperref[c:practice-0]{Forward Deployed Engineer}\,{\color{muted}\scriptsize p.\,\pageref{c:practice-0}}\\
2026 & \href{https://ualinkconsortium.org/wp-content/uploads/2026/01/UALink_White_Paper_Publication_Candidate_FINAL_VERSION.pdf}{UALink white paper (2026)} & \hyperref[c:hardware-3]{UALink}\,{\color{muted}\scriptsize p.\,\pageref{c:hardware-3}}\\
2026 & \href{https://venturebeat.com/technology/four-ai-research-trends-enterprise-teams-should-watch-in-2026}{VentureBeat: AI research trends for 2026} & \hyperref[c:embodied-3]{JEPA}\,{\color{muted}\scriptsize p.\,\pageref{c:embodied-3}}\\
2026 & \href{https://workos.com/blog/mcp-2026-spec-agent-authentication}{WorkOS: What changes for agent authentication} & \hyperref[c:protocols-0]{Model Context Protocol}\,{\color{muted}\scriptsize p.\,\pageref{c:protocols-0}}\\
2025 & \href{https://www.anthropic.com/engineering/advanced-tool-use}{Anthropic: Advanced tool use} & \hyperref[c:agents-6]{Tool search and programmatic tool calling}\,{\color{muted}\scriptsize p.\,\pageref{c:agents-6}}\\
2025 & \href{https://www.anthropic.com/research/agentic-misalignment}{Anthropic: Agentic misalignment} & \hyperref[c:safety-9]{Misalignment research}\,{\color{muted}\scriptsize p.\,\pageref{c:safety-9}}\\
2025 & \href{https://www.anthropic.com/engineering/claude-code-best-practices}{Anthropic: Claude Code best practices} & \hyperref[c:coding-2]{Agentic engineering / Software 3.0}\,{\color{muted}\scriptsize p.\,\pageref{c:coding-2}}\\
2025 & \href{https://www.anthropic.com/engineering/code-execution-with-mcp}{Anthropic: Code execution with MCP} & \hyperref[c:agents-6]{Tool search and programmatic tool calling}\,{\color{muted}\scriptsize p.\,\pageref{c:agents-6}}\\
2025 & \href{https://www.anthropic.com/research/constitutional-classifiers}{Anthropic: Constitutional classifiers} & \hyperref[c:safety-5]{Guardrails and classifiers}\,{\color{muted}\scriptsize p.\,\pageref{c:safety-5}}\\
2025 & \href{https://www.anthropic.com/engineering/effective-context-engineering-for-ai-agents}{Anthropic: Effective context engineering for AI agents} & \hyperref[c:context-0]{Context engineering}\,{\color{muted}\scriptsize p.\,\pageref{c:context-0}}\\
2025 & \href{https://www.anthropic.com/engineering/effective-harnesses-for-long-running-agents}{Anthropic: Effective harnesses for long-running agents} & \hyperref[c:agents-0]{Agent harness}\,{\color{muted}\scriptsize p.\,\pageref{c:agents-0}}\\
2025 & \href{https://www.anthropic.com/research/emergent-misalignment-reward-hacking}{Anthropic: Emergent misalignment from reward hacking} & \hyperref[c:reasoning-9]{Reward hacking}\,{\color{muted}\scriptsize p.\,\pageref{c:reasoning-9}}\\
2025 & \href{https://www.anthropic.com/engineering/equipping-agents-for-the-real-world-with-agent-skills}{Anthropic: Equipping agents with Agent Skills} & \hyperref[c:protocols-5]{Agent Skills}\,{\color{muted}\scriptsize p.\,\pageref{c:protocols-5}}\\
2025 & \href{https://www.anthropic.com/research/exploring-model-welfare}{Anthropic: Exploring model welfare} & \hyperref[c:safety-10]{Model welfare}\,{\color{muted}\scriptsize p.\,\pageref{c:safety-10}}\\
2025 & \href{https://www.anthropic.com/research/introspection}{Anthropic: Introspection} & \hyperref[c:safety-7]{Persona vectors and introspection}\,{\color{muted}\scriptsize p.\,\pageref{c:safety-7}}\\
2025 & \href{https://www.anthropic.com/news/context-management}{Anthropic: Managing context on the Claude Developer Platform} & \hyperref[c:context-2]{Compaction and context editing}\,{\color{muted}\scriptsize p.\,\pageref{c:context-2}}\\
2025 & \href{https://www.anthropic.com/engineering/multi-agent-research-system}{Anthropic: Multi-agent research system} & \hyperref[c:agents-7]{Multi-agent systems and subagents}\,{\color{muted}\scriptsize p.\,\pageref{c:agents-7}}\\
2025 & \href{https://www.anthropic.com/research/persona-vectors}{Anthropic: Persona vectors} & \hyperref[c:safety-7]{Persona vectors and introspection}\,{\color{muted}\scriptsize p.\,\pageref{c:safety-7}}\\
2025 & \href{https://www.anthropic.com/engineering/writing-tools-for-agents}{Anthropic: Writing effective tools for agents} & \hyperref[c:agents-5]{Tool use / function calling}\,{\color{muted}\scriptsize p.\,\pageref{c:agents-5}}\\
2025 & \href{https://blogs.arista.com/blog/the-sun-rises-on-scale-up-ethernet}{Arista: The sun rises on scale-up Ethernet} & \hyperref[c:hardware-4]{ESUN and SUE}\,{\color{muted}\scriptsize p.\,\pageref{c:hardware-4}}\\
2025 & \href{https://openai.com/index/browsecomp/}{BrowseComp} & \hyperref[c:evals-5]{Agent and tool-use benchmarks}\,{\color{muted}\scriptsize p.\,\pageref{c:evals-5}}\\
2025 & \href{https://research.trychroma.com/context-rot}{Chroma Research: Context rot} & \hyperref[c:context-1]{Context rot}\,{\color{muted}\scriptsize p.\,\pageref{c:context-1}}\\
2025 & \href{https://transformer-circuits.pub/2025/attribution-graphs/methods.html}{Circuit tracing} & \hyperref[c:safety-6]{Mechanistic interpretability}\,{\color{muted}\scriptsize p.\,\pageref{c:safety-6}}\\
2025 & \href{https://blog.cloudflare.com/code-mode/}{Cloudflare: Code Mode} & \hyperref[c:agents-6]{Tool search and programmatic tool calling}\,{\color{muted}\scriptsize p.\,\pageref{c:agents-6}}\\
2025 & \href{https://deepmind.google/discover/blog/alphaevolve-a-gemini-powered-coding-agent-for-designing-advanced-algorithms/}{DeepMind: AlphaEvolve} & \hyperref[c:embodied-7]{AI for science and algorithm discovery}\,{\color{muted}\scriptsize p.\,\pageref{c:embodied-7}}\\
2025 & \href{https://deepmind.google/discover/blog/genie-3-a-new-frontier-for-world-models/}{DeepMind: Genie 3} & \hyperref[c:embodied-2]{World models}\,{\color{muted}\scriptsize p.\,\pageref{c:embodied-2}}\\
2025 & \href{https://ghuntley.com/ralph/}{Geoffrey Huntley: Ralph} & \hyperref[c:coding-4]{Ralph loop}\,{\color{muted}\scriptsize p.\,\pageref{c:coding-4}}\\
2025 & \href{https://research.google/blog/introducing-nested-learning-a-new-ml-paradigm-for-continual-learning/}{Google Research: Nested learning} & \hyperref[c:reasoning-11]{Continual learning}\,{\color{muted}\scriptsize p.\,\pageref{c:reasoning-11}}\\
2025 & \href{https://developers.googleblog.com/en/introducing-gemini-2-5-flash-image/}{Google: Gemini 2.5 Flash Image} & \hyperref[c:vlm-8]{Native image generation and editing}\,{\color{muted}\scriptsize p.\,\pageref{c:vlm-8}}\\
2025 & \href{https://blog.google/products/google-cloud/ironwood-tpu-age-of-inference/}{Google: Ironwood TPU} & \hyperref[c:hardware-11]{Custom accelerators}\,{\color{muted}\scriptsize p.\,\pageref{c:hardware-11}}\\
2025 & \href{https://huggingface.co/spaces/nanotron/ultrascale-playbook}{HF: The Ultra-Scale Playbook} & \hyperref[c:training-0]{Parallelism strategies}\,{\color{muted}\scriptsize p.\,\pageref{c:training-0}}\\
2025 & \href{https://jax-ml.github.io/scaling-book/}{How to Scale Your Model (JAX)} & \hyperref[c:training-0]{Parallelism strategies}\,{\color{muted}\scriptsize p.\,\pageref{c:training-0}}\\
2025 & \href{https://huggingface.co/blog/smollm3}{Hugging Face: SmolLM3} & \hyperref[c:arch-11]{Small language models}\,{\color{muted}\scriptsize p.\,\pageref{c:arch-11}}\\
2025 & \href{https://huggingface.co/blog/vlms-2025}{Hugging Face: Vision language models 2025} & \hyperref[c:vlm-0]{VLM architecture}\,{\color{muted}\scriptsize p.\,\pageref{c:vlm-0}}\\
2025 & \href{https://invariantlabs.ai/blog/mcp-security-notification-tool-poisoning-attacks}{Invariant Labs: MCP tool poisoning} & \hyperref[c:safety-4]{MCP and tool security}\,{\color{muted}\scriptsize p.\,\pageref{c:safety-4}}\\
2025 & \href{https://x.com/karpathy/status/1886192184808149383}{Karpathy's original post} & \hyperref[c:coding-1]{Vibe coding}\,{\color{muted}\scriptsize p.\,\pageref{c:coding-1}}\\
2025 & \href{https://blog.langchain.com/context-engineering-for-agents/}{LangChain: Context engineering for agents} & \hyperref[c:context-0]{Context engineering}\,{\color{muted}\scriptsize p.\,\pageref{c:context-0}}\\
2025 & \href{https://metr.org/blog/2025-03-19-measuring-ai-ability-to-complete-long-tasks/}{METR: Measuring AI ability to complete long tasks} & \hyperref[c:agents-8]{Long-running and background agents}\,{\color{muted}\scriptsize p.\,\pageref{c:agents-8}}\\
2025 & \href{https://developer.nvidia.com/blog/nvidia-800-v-hvdc-architecture-will-power-the-next-generation-of-ai-factories/}{NVIDIA: 800 V HVDC architecture} & \hyperref[c:hardware-12]{AI factories, power and cooling}\,{\color{muted}\scriptsize p.\,\pageref{c:hardware-12}}\\
2025 & \href{https://developer.nvidia.com/blog/introducing-nvfp4-for-efficient-and-accurate-low-precision-inference/}{NVIDIA: Introducing NVFP4} & \hyperref[c:quant-1]{NVFP4}\,{\color{muted}\scriptsize p.\,\pageref{c:quant-1}}\\
2025 & \href{https://openai.com/index/updating-our-preparedness-framework/}{OpenAI Preparedness Framework} & \hyperref[c:safety-0]{Capability thresholds and frontier safety frameworks}\,{\color{muted}\scriptsize p.\,\pageref{c:safety-0}}\\
2025 & \href{https://openai.com/index/introducing-4o-image-generation/}{OpenAI: 4o image generation} & \hyperref[c:vlm-8]{Native image generation and editing}\,{\color{muted}\scriptsize p.\,\pageref{c:vlm-8}}\\
2025 & \href{https://openai.com/index/computer-using-agent/}{OpenAI: Computer-Using Agent} & \hyperref[c:agents-9]{Computer use and browser agents}\,{\color{muted}\scriptsize p.\,\pageref{c:agents-9}}\\
2025 & \href{https://openai.com/index/gdpval/}{OpenAI: GDPval} & \hyperref[c:evals-7]{Economic value evals}\,{\color{muted}\scriptsize p.\,\pageref{c:evals-7}}\\
2025 & \href{https://openai.com/index/introducing-deep-research/}{OpenAI: Introducing deep research} & \hyperref[c:context-6]{Agentic search and deep research}\,{\color{muted}\scriptsize p.\,\pageref{c:context-6}}\\
2025 & \href{https://openai.com/index/sora-2/}{OpenAI: Sora 2} & \hyperref[c:embodied-6]{Video generation}\,{\color{muted}\scriptsize p.\,\pageref{c:embodied-6}}\\
2025 & \href{https://openai.com/index/announcing-the-stargate-project/}{OpenAI: Stargate} & \hyperref[c:hardware-12]{AI factories, power and cooling}\,{\color{muted}\scriptsize p.\,\pageref{c:hardware-12}}\\
2025 & \href{https://openai.com/index/thinking-with-images/}{OpenAI: Thinking with images} & \hyperref[c:vlm-3]{Thinking with images}\,{\color{muted}\scriptsize p.\,\pageref{c:vlm-3}}\\
2025 & \href{https://magazine.sebastianraschka.com/p/the-big-llm-architecture-comparison}{Raschka: The big LLM architecture comparison} & \hyperref[c:arch-12]{Architecture overviews}\,{\color{muted}\scriptsize p.\,\pageref{c:arch-12}}\\
2025 & \href{https://simonwillison.net/2025/Mar/19/vibe-coding/}{Simon Willison: Not all AI-assisted programming is vibe coding} & \hyperref[c:coding-1]{Vibe coding}\,{\color{muted}\scriptsize p.\,\pageref{c:coding-1}}\\
2025 & \href{https://simonwillison.net/2025/Jun/16/the-lethal-trifecta/}{Simon Willison: The lethal trifecta} & \hyperref[c:safety-3]{Prompt injection and the lethal trifecta}\,{\color{muted}\scriptsize p.\,\pageref{c:safety-3}}\\
2025 & \href{https://thinkingmachines.ai/blog/lora/}{Thinking Machines: LoRA without regret} & \hyperref[c:training-5]{LoRA and parameter-efficient fine-tuning}\,{\color{muted}\scriptsize p.\,\pageref{c:training-5}}\\
2025 & \href{https://thinkingmachines.ai/blog/on-policy-distillation/}{Thinking Machines: On-policy distillation} & \hyperref[c:reasoning-3]{Multi-teacher on-policy distillation}\,{\color{muted}\scriptsize p.\,\pageref{c:reasoning-3}}\\
2024 & \href{https://www.anthropic.com/engineering}{Anthropic Engineering} & \hyperref[c:learn-0]{Deep technical blogs}\,{\color{muted}\scriptsize p.\,\pageref{c:learn-0}}\\
2024 & \href{https://www.anthropic.com/research/alignment-faking}{Anthropic: Alignment faking} & \hyperref[c:safety-9]{Misalignment research}\,{\color{muted}\scriptsize p.\,\pageref{c:safety-9}}\\
2024 & \href{https://www.anthropic.com/engineering/building-effective-agents}{Anthropic: Building effective agents} & \hyperref[c:agents-3]{Agentic loop / ReAct}\,{\color{muted}\scriptsize p.\,\pageref{c:agents-3}}\\
2024 & \href{https://www.anthropic.com/news/contextual-retrieval}{Anthropic: Contextual retrieval} & \hyperref[c:context-5]{Retrieval-augmented generation}\,{\color{muted}\scriptsize p.\,\pageref{c:context-5}}\\
2024 & \href{https://huggingface.co/blog/cosmopedia}{Cosmopedia} & \hyperref[c:training-3]{Data curation and synthetic data}\,{\color{muted}\scriptsize p.\,\pageref{c:training-3}}\\
2024 & \href{https://huggingface.co/spaces/HuggingFaceFW/blogpost-fineweb-v1}{FineWeb} & \hyperref[c:training-3]{Data curation and synthetic data}\,{\color{muted}\scriptsize p.\,\pageref{c:training-3}}\\
2024 & \href{https://hamel.dev/blog/posts/evals/}{Hamel Husain: Your AI product needs evals} & \hyperref[c:practice-4]{Eval-driven development}\,{\color{muted}\scriptsize p.\,\pageref{c:practice-4}}\\
2024 & \href{https://kellerjordan.github.io/posts/muon/}{Keller Jordan: Muon} & \hyperref[c:training-2]{Muon and modern optimizers}\,{\color{muted}\scriptsize p.\,\pageref{c:training-2}}\\
2024 & \href{https://lilianweng.github.io/posts/2024-11-28-reward-hacking/}{Lilian Weng: Reward hacking} & \hyperref[c:reasoning-9]{Reward hacking}\,{\color{muted}\scriptsize p.\,\pageref{c:reasoning-9}}\\
2024 & \href{https://openai.com/index/learning-to-reason-with-llms/}{OpenAI: Learning to reason with LLMs} & \hyperref[c:reasoning-0]{Reasoning models and test-time compute}\,{\color{muted}\scriptsize p.\,\pageref{c:reasoning-0}}\\
2024 & \href{https://huggingface.co/blog/smolvlm}{SmolVLM} & \hyperref[c:vlm-1]{Open VLM families}\,{\color{muted}\scriptsize p.\,\pageref{c:vlm-1}}\\
2023 & \href{https://huggingface.co/blog/moe}{Hugging Face: MoE explained} & \hyperref[c:arch-0]{Mixture of Experts}\,{\color{muted}\scriptsize p.\,\pageref{c:arch-0}}\\
2023 & \href{https://www.latent.space}{Latent Space} & \hyperref[c:learn-0]{Deep technical blogs}\,{\color{muted}\scriptsize p.\,\pageref{c:learn-0}}\\
2023 & \href{https://www.latent.space/p/ai-engineer}{Latent Space: The rise of the AI engineer} & \hyperref[c:practice-1]{AI Engineer}\,{\color{muted}\scriptsize p.\,\pageref{c:practice-1}}\\
2022 & \href{https://magazine.sebastianraschka.com}{Ahead of AI (Sebastian Raschka)} & \hyperref[c:learn-0]{Deep technical blogs}\,{\color{muted}\scriptsize p.\,\pageref{c:learn-0}}\\
2022 & \href{https://www.interconnects.ai}{Interconnects (Nathan Lambert)} & \hyperref[c:learn-0]{Deep technical blogs}\,{\color{muted}\scriptsize p.\,\pageref{c:learn-0}}\\
2022 & \href{https://simonwillison.net/series/prompt-injection/}{Simon Willison: prompt injection series} & \hyperref[c:safety-3]{Prompt injection and the lethal trifecta}\,{\color{muted}\scriptsize p.\,\pageref{c:safety-3}}\\
2017 & \href{https://lilianweng.github.io}{Lil'Log (Lilian Weng)} & \hyperref[c:learn-0]{Deep technical blogs}\,{\color{muted}\scriptsize p.\,\pageref{c:learn-0}}\\
2002 & \href{https://simonwillison.net}{Simon Willison's weblog} & \hyperref[c:learn-0]{Deep technical blogs}\,{\color{muted}\scriptsize p.\,\pageref{c:learn-0}}\\
\end{longtable}}
\section{Paper (45)}
{\small\begin{longtable}{@{}p{0.07\linewidth}>{\raggedright}p{0.55\linewidth}>{\raggedright\arraybackslash}p{0.34\linewidth}@{}}
\textbf{\sffamily Year} & \textbf{\sffamily Resource} & \textbf{\sffamily Concept}\\ \hline\endhead
2026 & \href{https://arxiv.org/abs/2606.19348}{DeepSeek-V4 technical report} & \hyperref[c:landscape-3]{Open-weight frontier models}\,{\color{muted}\scriptsize p.\,\pageref{c:landscape-3}}\\
2026 & \href{https://arxiv.org/abs/2606.10402}{EinsteinArena: collective agent discovery} & \hyperref[c:embodied-7]{AI for science and algorithm discovery}\,{\color{muted}\scriptsize p.\,\pageref{c:embodied-7}}\\
2026 & \href{https://arxiv.org/abs/2604.11548}{SemaClaw: personal agents via harness engineering} & \hyperref[c:agents-10]{Personal agents}\,{\color{muted}\scriptsize p.\,\pageref{c:agents-10}}\\
2026 & \href{https://arxiv.org/abs/2603.12658}{Survey: Continual learning in LLMs (2026)} & \hyperref[c:reasoning-11]{Continual learning}\,{\color{muted}\scriptsize p.\,\pageref{c:reasoning-11}}\\
2025 & \href{https://arxiv.org/abs/2505.03335}{Absolute Zero} & \hyperref[c:reasoning-10]{Self-play and self-generated curricula}\,{\color{muted}\scriptsize p.\,\pageref{c:reasoning-10}}\\
2025 & \href{https://arxiv.org/abs/2503.18813}{CaMeL} & \hyperref[c:safety-3]{Prompt injection and the lethal trifecta}\,{\color{muted}\scriptsize p.\,\pageref{c:safety-3}}\\
2025 & \href{https://arxiv.org/abs/2507.11473}{CoT monitorability position paper} & \hyperref[c:safety-8]{Chain-of-thought monitorability}\,{\color{muted}\scriptsize p.\,\pageref{c:safety-8}}\\
2025 & \href{https://arxiv.org/abs/2503.14476}{DAPO} & \hyperref[c:reasoning-4]{GRPO and its successors}\,{\color{muted}\scriptsize p.\,\pageref{c:reasoning-4}}\\
2025 & \href{https://arxiv.org/abs/2501.12948}{DeepSeek-R1} & \hyperref[c:reasoning-0]{Reasoning models and test-time compute}\,{\color{muted}\scriptsize p.\,\pageref{c:reasoning-0}}\\
2025 & \href{https://arxiv.org/abs/2503.20783}{Dr. GRPO} & \hyperref[c:reasoning-4]{GRPO and its successors}\,{\color{muted}\scriptsize p.\,\pageref{c:reasoning-4}}\\
2025 & \href{https://arxiv.org/abs/2502.17424}{Emergent misalignment} & \hyperref[c:safety-9]{Misalignment research}\,{\color{muted}\scriptsize p.\,\pageref{c:safety-9}}\\
2025 & \href{https://arxiv.org/abs/2507.18071}{GSPO} & \hyperref[c:reasoning-4]{GRPO and its successors}\,{\color{muted}\scriptsize p.\,\pageref{c:reasoning-4}}\\
2025 & \href{https://arxiv.org/abs/2502.11089}{Native Sparse Attention} & \hyperref[c:arch-2]{Sparse and compressed attention}\,{\color{muted}\scriptsize p.\,\pageref{c:arch-2}}\\
2025 & \href{https://arxiv.org/abs/2501.19393}{s1: simple test-time scaling} & \hyperref[c:reasoning-1]{Thinking budgets, adaptive thinking and effort}\,{\color{muted}\scriptsize p.\,\pageref{c:reasoning-1}}\\
2025 & \href{https://arxiv.org/abs/2502.14786}{SigLIP 2} & \hyperref[c:vlm-0]{VLM architecture}\,{\color{muted}\scriptsize p.\,\pageref{c:vlm-0}}\\
2024 & \href{https://arxiv.org/abs/2405.04434}{DeepSeek-V2 (MLA)} & \hyperref[c:arch-1]{GQA and Multi-head Latent Attention}\,{\color{muted}\scriptsize p.\,\pageref{c:arch-1}}\\
2024 & \href{https://arxiv.org/abs/2412.19437}{DeepSeek-V3} & \hyperref[c:arch-0]{Mixture of Experts}\,{\color{muted}\scriptsize p.\,\pageref{c:arch-0}}\\
2024 & \href{https://arxiv.org/abs/2402.03300}{DeepSeekMath (GRPO)} & \hyperref[c:reasoning-4]{GRPO and its successors}\,{\color{muted}\scriptsize p.\,\pageref{c:reasoning-4}}\\
2024 & \href{https://arxiv.org/abs/2401.09670}{DistServe} & \hyperref[c:inference-3]{Prefill/decode disaggregation}\,{\color{muted}\scriptsize p.\,\pageref{c:inference-3}}\\
2024 & \href{https://arxiv.org/abs/2404.19737}{Multi-token prediction} & \hyperref[c:arch-6]{Multi-token prediction and speculative decoding}\,{\color{muted}\scriptsize p.\,\pageref{c:arch-6}}\\
2024 & \href{https://arxiv.org/abs/2408.03314}{Scaling test-time compute optimally} & \hyperref[c:reasoning-0]{Reasoning models and test-time compute}\,{\color{muted}\scriptsize p.\,\pageref{c:reasoning-0}}\\
2024 & \href{https://arxiv.org/abs/2411.15124}{Tulu 3} & \hyperref[c:reasoning-2]{RL with verifiable rewards}\,{\color{muted}\scriptsize p.\,\pageref{c:reasoning-2}}\\
2023 & \href{https://arxiv.org/abs/2311.08105}{DiLoCo} & \hyperref[c:training-8]{Decentralized training}\,{\color{muted}\scriptsize p.\,\pageref{c:training-8}}\\
2023 & \href{https://arxiv.org/abs/2305.18290}{DPO} & \hyperref[c:reasoning-6]{RLHF, DPO and preference optimization}\,{\color{muted}\scriptsize p.\,\pageref{c:reasoning-6}}\\
2023 & \href{https://arxiv.org/abs/2305.13245}{GQA} & \hyperref[c:arch-1]{GQA and Multi-head Latent Attention}\,{\color{muted}\scriptsize p.\,\pageref{c:arch-1}}\\
2023 & \href{https://arxiv.org/abs/2306.05685}{Judging LLM-as-a-judge} & \hyperref[c:evals-2]{LLM-as-a-judge}\,{\color{muted}\scriptsize p.\,\pageref{c:evals-2}}\\
2023 & \href{https://arxiv.org/abs/2305.20050}{Let's verify step by step} & \hyperref[c:reasoning-8]{Process reward models}\,{\color{muted}\scriptsize p.\,\pageref{c:reasoning-8}}\\
2023 & \href{https://arxiv.org/abs/2310.10537}{Microscaling data formats} & \hyperref[c:quant-0]{Microscaling formats}\,{\color{muted}\scriptsize p.\,\pageref{c:quant-0}}\\
2023 & \href{https://arxiv.org/abs/2307.06304}{NaViT} & \hyperref[c:vlm-2]{Native and dynamic resolution}\,{\color{muted}\scriptsize p.\,\pageref{c:vlm-2}}\\
2023 & \href{https://arxiv.org/abs/2309.06180}{PagedAttention} & \hyperref[c:inference-1]{KV cache and PagedAttention}\,{\color{muted}\scriptsize p.\,\pageref{c:inference-1}}\\
2023 & \href{https://arxiv.org/abs/2305.14314}{QLoRA} & \hyperref[c:training-5]{LoRA and parameter-efficient fine-tuning}\,{\color{muted}\scriptsize p.\,\pageref{c:training-5}}\\
2023 & \href{https://arxiv.org/abs/2309.17453}{StreamingLLM} & \hyperref[c:arch-8]{Attention sinks}\,{\color{muted}\scriptsize p.\,\pageref{c:arch-8}}\\
2022 & \href{https://arxiv.org/abs/2203.15556}{Chinchilla} & \hyperref[c:training-1]{Scaling laws}\,{\color{muted}\scriptsize p.\,\pageref{c:training-1}}\\
2022 & \href{https://arxiv.org/abs/2212.08073}{Constitutional AI} & \hyperref[c:reasoning-7]{Constitutional AI and model specs}\,{\color{muted}\scriptsize p.\,\pageref{c:reasoning-7}}\\
2022 & \href{https://arxiv.org/abs/2212.09748}{DiT} & \hyperref[c:embodied-5]{Diffusion transformers and flow matching}\,{\color{muted}\scriptsize p.\,\pageref{c:embodied-5}}\\
2022 & \href{https://arxiv.org/abs/2210.02747}{Flow matching} & \hyperref[c:embodied-5]{Diffusion transformers and flow matching}\,{\color{muted}\scriptsize p.\,\pageref{c:embodied-5}}\\
2022 & \href{https://arxiv.org/abs/2209.05433}{FP8 formats for deep learning} & \hyperref[c:quant-2]{FP8 training and inference}\,{\color{muted}\scriptsize p.\,\pageref{c:quant-2}}\\
2022 & \href{https://arxiv.org/abs/2210.17323}{GPTQ} & \hyperref[c:quant-3]{Post-training quantization}\,{\color{muted}\scriptsize p.\,\pageref{c:quant-3}}\\
2022 & \href{https://arxiv.org/abs/2203.02155}{InstructGPT} & \hyperref[c:reasoning-6]{RLHF, DPO and preference optimization}\,{\color{muted}\scriptsize p.\,\pageref{c:reasoning-6}}\\
2022 & \href{https://arxiv.org/abs/2205.13147}{Matryoshka representation learning} & \hyperref[c:context-8]{Embeddings, late interaction and Matryoshka}\,{\color{muted}\scriptsize p.\,\pageref{c:context-8}}\\
2022 & \href{https://arxiv.org/abs/2210.03629}{ReAct paper} & \hyperref[c:agents-3]{Agentic loop / ReAct}\,{\color{muted}\scriptsize p.\,\pageref{c:agents-3}}\\
2022 & \href{https://arxiv.org/abs/2211.17192}{Speculative decoding} & \hyperref[c:arch-6]{Multi-token prediction and speculative decoding}\,{\color{muted}\scriptsize p.\,\pageref{c:arch-6}}\\
2021 & \href{https://arxiv.org/abs/2106.09685}{LoRA} & \hyperref[c:training-5]{LoRA and parameter-efficient fine-tuning}\,{\color{muted}\scriptsize p.\,\pageref{c:training-5}}\\
2021 & \href{https://arxiv.org/abs/2104.09864}{RoFormer (RoPE)} & \hyperref[c:arch-7]{Long-context methods}\,{\color{muted}\scriptsize p.\,\pageref{c:arch-7}}\\
2020 & \href{https://arxiv.org/abs/2005.11401}{Original RAG paper} & \hyperref[c:context-5]{Retrieval-augmented generation}\,{\color{muted}\scriptsize p.\,\pageref{c:context-5}}\\
\end{longtable}}
\section{Docs (77)}
{\small\begin{longtable}{@{}p{0.07\linewidth}>{\raggedright}p{0.55\linewidth}>{\raggedright\arraybackslash}p{0.34\linewidth}@{}}
\textbf{\sffamily Year} & \textbf{\sffamily Resource} & \textbf{\sffamily Concept}\\ \hline\endhead
2026 & \href{https://www.anthropic.com/news/fable-mythos-access}{Anthropic statement on Fable and Mythos access} & \hyperref[c:landscape-0]{Mythos-class tier and trusted access}\,{\color{muted}\scriptsize p.\,\pageref{c:landscape-0}}\\
2026 & \href{https://www.anthropic.com/glasswing}{Anthropic: Project Glasswing} & \hyperref[c:landscape-0]{Mythos-class tier and trusted access}\,{\color{muted}\scriptsize p.\,\pageref{c:landscape-0}}\\
2026 & \href{https://www.anthropic.com/constitution}{Claude's constitution} & \hyperref[c:reasoning-7]{Constitutional AI and model specs}\,{\color{muted}\scriptsize p.\,\pageref{c:reasoning-7}}\\
2026 & \href{https://deploymentsafety.openai.com/gpt-6-astra}{GPT-6 Astra system card} & \hyperref[c:landscape-1]{GPT-6 Astra and the Critical cyber threshold}\,{\color{muted}\scriptsize p.\,\pageref{c:landscape-1}}\\
2026 & \href{https://rlhfbook.com/teach/course/conversation-01/}{RLHF Book course: Frontier post-training recipe survey} & \hyperref[c:reasoning-3]{Multi-teacher on-policy distillation}\,{\color{muted}\scriptsize p.\,\pageref{c:reasoning-3}}\\
2026 & \href{https://en.wikipedia.org/wiki/Forward_Deployed_Engineer}{Wikipedia: Forward Deployed Engineer} & \hyperref[c:practice-0]{Forward Deployed Engineer}\,{\color{muted}\scriptsize p.\,\pageref{c:practice-0}}\\
2025 & \href{https://aaif.io}{aaif.io} & \hyperref[c:protocols-4]{Agentic AI Foundation}\,{\color{muted}\scriptsize p.\,\pageref{c:protocols-4}}\\
2025 & \href{https://docs.ag-ui.com}{AG-UI docs} & \hyperref[c:protocols-8]{AG-UI}\,{\color{muted}\scriptsize p.\,\pageref{c:protocols-8}}\\
2025 & \href{https://developer.apple.com/documentation/foundationmodels}{Apple Foundation Models framework} & \hyperref[c:arch-11]{Small language models}\,{\color{muted}\scriptsize p.\,\pageref{c:arch-11}}\\
2025 & \href{https://leginfo.legislature.ca.gov/faces/billNavClient.xhtml?bill_id=202520260SB53}{California SB 53} & \hyperref[c:safety-12]{Other regulation and standards}\,{\color{muted}\scriptsize p.\,\pageref{c:safety-12}}\\
2025 & \href{https://docs.claude.com/en/docs/claude-code/hooks}{Claude Code hooks} & \hyperref[c:coding-5]{Hooks, plugins and slash commands}\,{\color{muted}\scriptsize p.\,\pageref{c:coding-5}}\\
2025 & \href{https://docs.claude.com/en/docs/claude-code/plugins}{Claude Code plugins} & \hyperref[c:coding-5]{Hooks, plugins and slash commands}\,{\color{muted}\scriptsize p.\,\pageref{c:coding-5}}\\
2025 & \href{https://docs.claude.com/en/docs/claude-code/sub-agents}{Claude Code subagents} & \hyperref[c:agents-7]{Multi-agent systems and subagents}\,{\color{muted}\scriptsize p.\,\pageref{c:agents-7}}\\
2025 & \href{https://docs.claude.com/en/docs/build-with-claude/extended-thinking}{Claude extended thinking docs} & \hyperref[c:reasoning-1]{Thinking budgets, adaptive thinking and effort}\,{\color{muted}\scriptsize p.\,\pageref{c:reasoning-1}}\\
2025 & \href{https://digital-strategy.ec.europa.eu/en/policies/contents-code-gpai}{EU GPAI Code of Practice} & \hyperref[c:safety-11]{EU AI Act and the Digital Omnibus}\,{\color{muted}\scriptsize p.\,\pageref{c:safety-11}}\\
2025 & \href{https://deepmind.google/models/gemini-diffusion/}{Gemini Diffusion} & \hyperref[c:arch-5]{Diffusion language models}\,{\color{muted}\scriptsize p.\,\pageref{c:arch-5}}\\
2025 & \href{https://deepmind.google/models/gemini-robotics/}{Gemini Robotics} & \hyperref[c:embodied-0]{Vision-language-action models}\,{\color{muted}\scriptsize p.\,\pageref{c:embodied-0}}\\
2025 & \href{https://lastexam.ai}{Humanity's Last Exam} & \hyperref[c:evals-6]{Frontier knowledge benchmarks}\,{\color{muted}\scriptsize p.\,\pageref{c:evals-6}}\\
2025 & \href{https://www.inceptionlabs.ai}{Inception Labs (Mercury)} & \hyperref[c:arch-5]{Diffusion language models}\,{\color{muted}\scriptsize p.\,\pageref{c:arch-5}}\\
2025 & \href{https://kiro.dev}{Kiro} & \hyperref[c:coding-3]{Spec-driven development}\,{\color{muted}\scriptsize p.\,\pageref{c:coding-3}}\\
2025 & \href{https://www.nvidia.com/en-us/networking/products/silicon-photonics/}{NVIDIA silicon photonics} & \hyperref[c:hardware-10]{Co-packaged optics}\,{\color{muted}\scriptsize p.\,\pageref{c:hardware-10}}\\
2025 & \href{https://www.nvidia.com/en-us/data-center/nvlink-fusion/}{NVLink Fusion} & \hyperref[c:hardware-6]{NVLink and NVLink Fusion}\,{\color{muted}\scriptsize p.\,\pageref{c:hardware-6}}\\
2025 & \href{https://stanford-cs336.github.io}{Stanford CS336} & \hyperref[c:training-9]{Build-it-yourself references}\,{\color{muted}\scriptsize p.\,\pageref{c:training-9}}\\
2025 & \href{https://scale.com/leaderboard/swe_bench_pro_public}{SWE-bench Pro} & \hyperref[c:evals-3]{Agentic coding benchmarks}\,{\color{muted}\scriptsize p.\,\pageref{c:evals-3}}\\
2025 & \href{https://www.tbench.ai}{Terminal-Bench} & \hyperref[c:evals-3]{Agentic coding benchmarks}\,{\color{muted}\scriptsize p.\,\pageref{c:evals-3}}\\
2025 & \href{https://thinkingmachines.ai/tinker/}{Tinker} & \hyperref[c:training-6]{Fine-tuning toolkits and APIs}\,{\color{muted}\scriptsize p.\,\pageref{c:training-6}}\\
2025 & \href{https://www.nist.gov/caisi}{US CAISI (NIST)} & \hyperref[c:safety-13]{Government AI safety institutes}\,{\color{muted}\scriptsize p.\,\pageref{c:safety-13}}\\
2025 & \href{https://marble.worldlabs.ai}{World Labs Marble} & \hyperref[c:embodied-2]{World models}\,{\color{muted}\scriptsize p.\,\pageref{c:embodied-2}}\\
2024 & \href{https://allenai.org/olmo}{Allen AI OLMo} & \hyperref[c:training-4]{Mid-training}\,{\color{muted}\scriptsize p.\,\pageref{c:training-4}}\\
2024 & \href{https://docs.claude.com/en/docs/build-with-claude/prompt-engineering/overview}{Anthropic prompt engineering guide} & \hyperref[c:practice-3]{Prompt engineering and prompt optimization}\,{\color{muted}\scriptsize p.\,\pageref{c:practice-3}}\\
2024 & \href{https://artificialanalysis.ai}{Artificial Analysis} & \hyperref[c:landscape-4]{Model release trackers}\,{\color{muted}\scriptsize p.\,\pageref{c:landscape-4}}\\
2024 & \href{https://docs.claude.com/en/docs/agents-and-tools/tool-use/computer-use-tool}{Claude computer use tool} & \hyperref[c:agents-9]{Computer use and browser agents}\,{\color{muted}\scriptsize p.\,\pageref{c:agents-9}}\\
2024 & \href{https://docs.claude.com/en/docs/build-with-claude/prompt-caching}{Claude prompt caching} & \hyperref[c:context-3]{Prompt caching}\,{\color{muted}\scriptsize p.\,\pageref{c:context-3}}\\
2024 & \href{https://docs.claude.com/en/docs/agents-and-tools/tool-use/overview}{Claude tool use docs} & \hyperref[c:agents-5]{Tool use / function calling}\,{\color{muted}\scriptsize p.\,\pageref{c:agents-5}}\\
2024 & \href{https://epoch.ai/benchmarks}{Epoch AI benchmarking hub} & \hyperref[c:landscape-4]{Model release trackers}\,{\color{muted}\scriptsize p.\,\pageref{c:landscape-4}}\\
2024 & \href{https://artificialintelligenceact.eu}{EU AI Act explorer} & \hyperref[c:safety-11]{EU AI Act and the Digital Omnibus}\,{\color{muted}\scriptsize p.\,\pageref{c:safety-11}}\\
2024 & \href{https://epoch.ai/frontiermath}{FrontierMath} & \hyperref[c:evals-6]{Frontier knowledge benchmarks}\,{\color{muted}\scriptsize p.\,\pageref{c:evals-6}}\\
2024 & \href{https://deepmind.google/models/veo/}{Google Veo} & \hyperref[c:embodied-6]{Video generation}\,{\color{muted}\scriptsize p.\,\pageref{c:embodied-6}}\\
2024 & \href{https://metr.org/faisc}{METR: frontier AI safety policies} & \hyperref[c:safety-0]{Capability thresholds and frontier safety frameworks}\,{\color{muted}\scriptsize p.\,\pageref{c:safety-0}}\\
2024 & \href{https://modelcontextprotocol.io}{modelcontextprotocol.io} & \hyperref[c:protocols-0]{Model Context Protocol}\,{\color{muted}\scriptsize p.\,\pageref{c:protocols-0}}\\
2024 & \href{https://www.nvidia.com/en-us/data-center/gb200-nvl72/}{NVIDIA GB200 NVL72} & \hyperref[c:hardware-0]{Rack-scale systems}\,{\color{muted}\scriptsize p.\,\pageref{c:hardware-0}}\\
2024 & \href{https://platform.openai.com/docs/guides/realtime}{OpenAI Realtime API} & \hyperref[c:vlm-7]{Omni and speech-to-speech models}\,{\color{muted}\scriptsize p.\,\pageref{c:vlm-7}}\\
2024 & \href{https://platform.openai.com/docs/guides/structured-outputs}{OpenAI structured outputs} & \hyperref[c:context-9]{Structured outputs and constrained decoding}\,{\color{muted}\scriptsize p.\,\pageref{c:context-9}}\\
2024 & \href{https://os-world.github.io}{OSWorld} & \hyperref[c:evals-5]{Agent and tool-use benchmarks}\,{\color{muted}\scriptsize p.\,\pageref{c:evals-5}}\\
2024 & \href{https://rlhfbook.com}{RLHF Book (Nathan Lambert)} & \hyperref[c:reasoning-2]{RL with verifiable rewards}\,{\color{muted}\scriptsize p.\,\pageref{c:reasoning-2}}\\
2024 & \href{https://en.wikipedia.org/wiki/UALink}{UALink overview} & \hyperref[c:hardware-3]{UALink}\,{\color{muted}\scriptsize p.\,\pageref{c:hardware-3}}\\
2023 & \href{https://www.amd.com/en/products/accelerators/instinct.html}{AMD Instinct} & \hyperref[c:hardware-1]{AMD Helios and Open Rack Wide}\,{\color{muted}\scriptsize p.\,\pageref{c:hardware-1}}\\
2023 & \href{https://www.anthropic.com/responsible-scaling-policy}{Anthropic Responsible Scaling Policy} & \hyperref[c:safety-0]{Capability thresholds and frontier safety frameworks}\,{\color{muted}\scriptsize p.\,\pageref{c:safety-0}}\\
2023 & \href{https://www.apolloresearch.ai}{Apollo Research} & \hyperref[c:safety-9]{Misalignment research}\,{\color{muted}\scriptsize p.\,\pageref{c:safety-9}}\\
2023 & \href{https://epoch.ai/data/machine-learning-hardware}{Epoch AI: ML hardware data} & \hyperref[c:hardware-16]{Hardware benchmarks}\,{\color{muted}\scriptsize p.\,\pageref{c:hardware-16}}\\
2023 & \href{https://deepmind.google/models/synthid/}{Google SynthID} & \hyperref[c:safety-14]{Content provenance and watermarking}\,{\color{muted}\scriptsize p.\,\pageref{c:safety-14}}\\
2023 & \href{https://huggingface.co/learn/llm-course}{Hugging Face LLM course} & \hyperref[c:learn-1]{Courses}\,{\color{muted}\scriptsize p.\,\pageref{c:learn-1}}\\
2023 & \href{https://lmarena.ai}{LMArena} & \hyperref[c:landscape-4]{Model release trackers}\,{\color{muted}\scriptsize p.\,\pageref{c:landscape-4}}\\
2023 & \href{https://mmmu-benchmark.github.io}{MMMU} & \hyperref[c:vlm-9]{Multimodal evaluation}\,{\color{muted}\scriptsize p.\,\pageref{c:vlm-9}}\\
2023 & \href{https://www.neuronpedia.org}{Neuronpedia} & \hyperref[c:safety-6]{Mechanistic interpretability}\,{\color{muted}\scriptsize p.\,\pageref{c:safety-6}}\\
2023 & \href{https://openrouter.ai}{OpenRouter} & \hyperref[c:inference-7]{LLM gateways and routers}\,{\color{muted}\scriptsize p.\,\pageref{c:inference-7}}\\
2023 & \href{https://openrouter.ai/models}{OpenRouter model list} & \hyperref[c:landscape-4]{Model release trackers}\,{\color{muted}\scriptsize p.\,\pageref{c:landscape-4}}\\
2023 & \href{https://genai.owasp.org/llm-top-10/}{OWASP GenAI Security Project} & \hyperref[c:safety-4]{MCP and tool security}\,{\color{muted}\scriptsize p.\,\pageref{c:safety-4}}\\
2023 & \href{https://www.swebench.com}{SWE-bench} & \hyperref[c:evals-3]{Agentic coding benchmarks}\,{\color{muted}\scriptsize p.\,\pageref{c:evals-3}}\\
2023 & \href{https://www.aisi.gov.uk}{UK AI Security Institute} & \hyperref[c:safety-13]{Government AI safety institutes}\,{\color{muted}\scriptsize p.\,\pageref{c:safety-13}}\\
2022 & \href{https://epoch.ai}{Epoch AI} & \hyperref[c:learn-2]{Annual reports and data}\,{\color{muted}\scriptsize p.\,\pageref{c:learn-2}}\\
2022 & \href{https://epoch.ai/data}{Epoch AI data} & \hyperref[c:training-1]{Scaling laws}\,{\color{muted}\scriptsize p.\,\pageref{c:training-1}}\\
2022 & \href{https://karpathy.ai/zero-to-hero.html}{Karpathy: Neural Networks Zero to Hero} & \hyperref[c:learn-1]{Courses}\,{\color{muted}\scriptsize p.\,\pageref{c:learn-1}}\\
2022 & \href{https://huggingface.co/spaces/mteb/leaderboard}{MTEB leaderboard} & \hyperref[c:context-8]{Embeddings, late interaction and Matryoshka}\,{\color{muted}\scriptsize p.\,\pageref{c:context-8}}\\
2021 & \href{https://transformer-circuits.pub}{Transformer Circuits} & \hyperref[c:safety-6]{Mechanistic interpretability}\,{\color{muted}\scriptsize p.\,\pageref{c:safety-6}}\\
2020 & \href{https://aws.amazon.com/ai/machine-learning/trainium/}{AWS Trainium} & \hyperref[c:hardware-11]{Custom accelerators}\,{\color{muted}\scriptsize p.\,\pageref{c:hardware-11}}\\
2019 & \href{https://arcprize.org}{ARC Prize} & \hyperref[c:evals-4]{ARC-AGI-3 and interactive reasoning}\,{\color{muted}\scriptsize p.\,\pageref{c:evals-4}}\\
2019 & \href{https://temporal.io}{Temporal} & \hyperref[c:agents-13]{Durable execution and agent trajectories}\,{\color{muted}\scriptsize p.\,\pageref{c:agents-13}}\\
2018 & \href{https://mlcommons.org/benchmarks/}{MLCommons benchmarks} & \hyperref[c:hardware-16]{Hardware benchmarks}\,{\color{muted}\scriptsize p.\,\pageref{c:hardware-16}}\\
2018 & \href{https://www.stateof.ai}{State of AI Report} & \hyperref[c:learn-2]{Annual reports and data}\,{\color{muted}\scriptsize p.\,\pageref{c:learn-2}}\\
2017 & \href{https://hai.stanford.edu/ai-index}{Stanford AI Index} & \hyperref[c:learn-2]{Annual reports and data}\,{\color{muted}\scriptsize p.\,\pageref{c:learn-2}}\\
2016 & \href{https://www.cerebras.ai}{Cerebras} & \hyperref[c:hardware-11]{Custom accelerators}\,{\color{muted}\scriptsize p.\,\pageref{c:hardware-11}}\\
2016 & \href{https://groq.com}{Groq} & \hyperref[c:hardware-11]{Custom accelerators}\,{\color{muted}\scriptsize p.\,\pageref{c:hardware-11}}\\
2016 & \href{https://huggingface.co/models}{Hugging Face models hub} & \hyperref[c:landscape-3]{Open-weight frontier models}\,{\color{muted}\scriptsize p.\,\pageref{c:landscape-3}}\\
2015 & \href{https://git-scm.com/docs/git-worktree}{git worktree docs} & \hyperref[c:coding-6]{Parallel agents with git worktrees}\,{\color{muted}\scriptsize p.\,\pageref{c:coding-6}}\\
2013 & \href{https://en.wikipedia.org/wiki/High_Bandwidth_Memory}{High Bandwidth Memory overview} & \hyperref[c:hardware-8]{High-bandwidth memory}\,{\color{muted}\scriptsize p.\,\pageref{c:hardware-8}}\\
2012 & \href{https://3dfabric.tsmc.com/english/dedicatedFoundry/technology/cowos.htm}{TSMC CoWoS} & \hyperref[c:hardware-9]{Advanced packaging}\,{\color{muted}\scriptsize p.\,\pageref{c:hardware-9}}\\
\end{longtable}}
\section{Spec (19)}
{\small\begin{longtable}{@{}p{0.07\linewidth}>{\raggedright}p{0.55\linewidth}>{\raggedright\arraybackslash}p{0.34\linewidth}@{}}
\textbf{\sffamily Year} & \textbf{\sffamily Resource} & \textbf{\sffamily Concept}\\ \hline\endhead
2026 & \href{https://ucp.dev}{Universal Commerce Protocol} & \hyperref[c:protocols-9]{Agentic commerce protocols}\,{\color{muted}\scriptsize p.\,\pageref{c:protocols-9}}\\
2025 & \href{https://a2a-protocol.org}{A2A protocol} & \hyperref[c:protocols-3]{Agent2Agent protocol}\,{\color{muted}\scriptsize p.\,\pageref{c:protocols-3}}\\
2025 & \href{https://agentclientprotocol.com}{agentclientprotocol.com} & \hyperref[c:protocols-7]{Agent Client Protocol}\,{\color{muted}\scriptsize p.\,\pageref{c:protocols-7}}\\
2025 & \href{https://agents.md}{agents.md} & \hyperref[c:protocols-6]{AGENTS.md}\,{\color{muted}\scriptsize p.\,\pageref{c:protocols-6}}\\
2025 & \href{https://agentskills.io}{agentskills.io} & \hyperref[c:protocols-5]{Agent Skills}\,{\color{muted}\scriptsize p.\,\pageref{c:protocols-5}}\\
2024 & \href{https://llmstxt.org}{llmstxt.org} & \hyperref[c:protocols-10]{llms.txt}\,{\color{muted}\scriptsize p.\,\pageref{c:protocols-10}}\\
2024 & \href{https://opensource.org/ai/open-source-ai-definition}{OSI Open Source AI Definition} & \hyperref[c:practice-7]{Open weights vs. open source AI}\,{\color{muted}\scriptsize p.\,\pageref{c:practice-7}}\\
2024 & \href{https://opentelemetry.io/docs/specs/semconv/gen-ai/}{OTel GenAI semantic conventions} & \hyperref[c:protocols-11]{OpenTelemetry GenAI conventions}\,{\color{muted}\scriptsize p.\,\pageref{c:protocols-11}}\\
2024 & \href{https://ualinkconsortium.org}{UALink Consortium} & \hyperref[c:hardware-3]{UALink}\,{\color{muted}\scriptsize p.\,\pageref{c:hardware-3}}\\
2023 & \href{https://github.com/ggml-org/ggml/blob/master/docs/gguf.md}{GGUF specification} & \hyperref[c:quant-5]{GGUF and local model formats}\,{\color{muted}\scriptsize p.\,\pageref{c:quant-5}}\\
2023 & \href{https://www.iso.org/standard/81230.html}{ISO/IEC 42001} & \hyperref[c:safety-12]{Other regulation and standards}\,{\color{muted}\scriptsize p.\,\pageref{c:safety-12}}\\
2023 & \href{https://www.nist.gov/itl/ai-risk-management-framework}{NIST AI RMF} & \hyperref[c:safety-12]{Other regulation and standards}\,{\color{muted}\scriptsize p.\,\pageref{c:safety-12}}\\
2023 & \href{https://www.opencompute.org/documents/ocp-microscaling-formats-mx-v1-0-spec-final-pdf}{OCP Microscaling (MX) spec} & \hyperref[c:quant-0]{Microscaling formats}\,{\color{muted}\scriptsize p.\,\pageref{c:quant-0}}\\
2023 & \href{https://ultraethernet.org}{Ultra Ethernet Consortium} & \hyperref[c:hardware-5]{Ultra Ethernet}\,{\color{muted}\scriptsize p.\,\pageref{c:hardware-5}}\\
2022 & \href{https://www.uciexpress.org}{UCIe} & \hyperref[c:hardware-7]{CXL, PCIe and UCIe}\,{\color{muted}\scriptsize p.\,\pageref{c:hardware-7}}\\
2021 & \href{https://c2pa.org}{C2PA} & \hyperref[c:safety-14]{Content provenance and watermarking}\,{\color{muted}\scriptsize p.\,\pageref{c:safety-14}}\\
2019 & \href{https://computeexpresslink.org}{CXL Consortium} & \hyperref[c:hardware-7]{CXL, PCIe and UCIe}\,{\color{muted}\scriptsize p.\,\pageref{c:hardware-7}}\\
2011 & \href{https://www.opencompute.org}{Open Compute Project} & \hyperref[c:hardware-12]{AI factories, power and cooling}\,{\color{muted}\scriptsize p.\,\pageref{c:hardware-12}}\\
1992 & \href{https://pcisig.com}{PCI-SIG} & \hyperref[c:hardware-7]{CXL, PCIe and UCIe}\,{\color{muted}\scriptsize p.\,\pageref{c:hardware-7}}\\
\end{longtable}}
\section{Video (1)}
{\small\begin{longtable}{@{}p{0.07\linewidth}>{\raggedright}p{0.55\linewidth}>{\raggedright\arraybackslash}p{0.34\linewidth}@{}}
\textbf{\sffamily Year} & \textbf{\sffamily Resource} & \textbf{\sffamily Concept}\\ \hline\endhead
2025 & \href{https://www.youtube.com/watch?v=LCEmiRjPEtQ}{Karpathy: Software is changing (again)} & \hyperref[c:coding-2]{Agentic engineering / Software 3.0}\,{\color{muted}\scriptsize p.\,\pageref{c:coding-2}}\\
\end{longtable}}
\chapter{Acronyms and aliases}
\begin{multicols}{2}\small\raggedright
\textbf{\sffamily\color{accent}800 VDC}\ \ \hyperref[c:hardware-12]{AI factories, power and cooling}\,{\color{muted}\scriptsize p.\,\pageref{c:hardware-12}}\par\vspace{2pt}
\textbf{\sffamily\color{accent}A2A}\ \ \hyperref[c:protocols-3]{Agent2Agent protocol}\,{\color{muted}\scriptsize p.\,\pageref{c:protocols-3}}\par\vspace{2pt}
\textbf{\sffamily\color{accent}AAIF}\ \ \hyperref[c:protocols-4]{Agentic AI Foundation}\,{\color{muted}\scriptsize p.\,\pageref{c:protocols-4}}\par\vspace{2pt}
\textbf{\sffamily\color{accent}ACP}\ \ \hyperref[c:protocols-7]{Agent Client Protocol}\,{\color{muted}\scriptsize p.\,\pageref{c:protocols-7}}\par\vspace{2pt}
\textbf{\sffamily\color{accent}ACP, AP2, UCP, x402}\ \ \hyperref[c:protocols-9]{Agentic commerce protocols}\,{\color{muted}\scriptsize p.\,\pageref{c:protocols-9}}\par\vspace{2pt}
\textbf{\sffamily\color{accent}AISI, CAISI}\ \ \hyperref[c:safety-13]{Government AI safety institutes}\,{\color{muted}\scriptsize p.\,\pageref{c:safety-13}}\par\vspace{2pt}
\textbf{\sffamily\color{accent}alignment faking, scheming}\ \ \hyperref[c:safety-9]{Misalignment research}\,{\color{muted}\scriptsize p.\,\pageref{c:safety-9}}\par\vspace{2pt}
\textbf{\sffamily\color{accent}ARC-AGI-3}\ \ \hyperref[c:evals-4]{ARC-AGI-3 and interactive reasoning}\,{\color{muted}\scriptsize p.\,\pageref{c:evals-4}}\par\vspace{2pt}
\textbf{\sffamily\color{accent}BitNet}\ \ \hyperref[c:quant-6]{Ternary and 1-bit models}\,{\color{muted}\scriptsize p.\,\pageref{c:quant-6}}\par\vspace{2pt}
\textbf{\sffamily\color{accent}BLT}\ \ \hyperref[c:arch-9]{Tokenizer-free and byte-level models}\,{\color{muted}\scriptsize p.\,\pageref{c:arch-9}}\par\vspace{2pt}
\textbf{\sffamily\color{accent}C2PA, SynthID}\ \ \hyperref[c:safety-14]{Content provenance and watermarking}\,{\color{muted}\scriptsize p.\,\pageref{c:safety-14}}\par\vspace{2pt}
\textbf{\sffamily\color{accent}code mode}\ \ \hyperref[c:agents-6]{Tool search and programmatic tool calling}\,{\color{muted}\scriptsize p.\,\pageref{c:agents-6}}\par\vspace{2pt}
\textbf{\sffamily\color{accent}ColPali}\ \ \hyperref[c:vlm-5]{Visual document retrieval}\,{\color{muted}\scriptsize p.\,\pageref{c:vlm-5}}\par\vspace{2pt}
\textbf{\sffamily\color{accent}CoWoS}\ \ \hyperref[c:hardware-9]{Advanced packaging}\,{\color{muted}\scriptsize p.\,\pageref{c:hardware-9}}\par\vspace{2pt}
\textbf{\sffamily\color{accent}CPO}\ \ \hyperref[c:hardware-10]{Co-packaged optics}\,{\color{muted}\scriptsize p.\,\pageref{c:hardware-10}}\par\vspace{2pt}
\textbf{\sffamily\color{accent}CUA}\ \ \hyperref[c:agents-9]{Computer use and browser agents}\,{\color{muted}\scriptsize p.\,\pageref{c:agents-9}}\par\vspace{2pt}
\textbf{\sffamily\color{accent}DiLoCo}\ \ \hyperref[c:training-8]{Decentralized training}\,{\color{muted}\scriptsize p.\,\pageref{c:training-8}}\par\vspace{2pt}
\textbf{\sffamily\color{accent}DiT}\ \ \hyperref[c:embodied-5]{Diffusion transformers and flow matching}\,{\color{muted}\scriptsize p.\,\pageref{c:embodied-5}}\par\vspace{2pt}
\textbf{\sffamily\color{accent}dLLM}\ \ \hyperref[c:arch-5]{Diffusion language models}\,{\color{muted}\scriptsize p.\,\pageref{c:arch-5}}\par\vspace{2pt}
\textbf{\sffamily\color{accent}DP, FSDP, TP, PP, EP, CP}\ \ \hyperref[c:training-0]{Parallelism strategies}\,{\color{muted}\scriptsize p.\,\pageref{c:training-0}}\par\vspace{2pt}
\textbf{\sffamily\color{accent}DSPy, GEPA}\ \ \hyperref[c:practice-3]{Prompt engineering and prompt optimization}\,{\color{muted}\scriptsize p.\,\pageref{c:practice-3}}\par\vspace{2pt}
\textbf{\sffamily\color{accent}E4M3, E5M2}\ \ \hyperref[c:quant-2]{FP8 training and inference}\,{\color{muted}\scriptsize p.\,\pageref{c:quant-2}}\par\vspace{2pt}
\textbf{\sffamily\color{accent}Ethernet for Scale-Up Networking}\ \ \hyperref[c:hardware-4]{ESUN and SUE}\,{\color{muted}\scriptsize p.\,\pageref{c:hardware-4}}\par\vspace{2pt}
\textbf{\sffamily\color{accent}EU AI Act}\ \ \hyperref[c:safety-11]{EU AI Act and the Digital Omnibus}\,{\color{muted}\scriptsize p.\,\pageref{c:safety-11}}\par\vspace{2pt}
\textbf{\sffamily\color{accent}Fable / Mythos}\ \ \hyperref[c:landscape-0]{Mythos-class tier and trusted access}\,{\color{muted}\scriptsize p.\,\pageref{c:landscape-0}}\par\vspace{2pt}
\textbf{\sffamily\color{accent}FDE}\ \ \hyperref[c:practice-0]{Forward Deployed Engineer}\,{\color{muted}\scriptsize p.\,\pageref{c:practice-0}}\par\vspace{2pt}
\textbf{\sffamily\color{accent}FlashAttention}\ \ \hyperref[c:inference-5]{Attention kernels}\,{\color{muted}\scriptsize p.\,\pageref{c:inference-5}}\par\vspace{2pt}
\textbf{\sffamily\color{accent}GDPval}\ \ \hyperref[c:evals-7]{Economic value evals}\,{\color{muted}\scriptsize p.\,\pageref{c:evals-7}}\par\vspace{2pt}
\textbf{\sffamily\color{accent}GPT-6}\ \ \hyperref[c:landscape-1]{GPT-6 Astra and the Critical cyber threshold}\,{\color{muted}\scriptsize p.\,\pageref{c:landscape-1}}\par\vspace{2pt}
\textbf{\sffamily\color{accent}GPTQ, AWQ}\ \ \hyperref[c:quant-3]{Post-training quantization}\,{\color{muted}\scriptsize p.\,\pageref{c:quant-3}}\par\vspace{2pt}
\textbf{\sffamily\color{accent}GQA, MLA}\ \ \hyperref[c:arch-1]{GQA and Multi-head Latent Attention}\,{\color{muted}\scriptsize p.\,\pageref{c:arch-1}}\par\vspace{2pt}
\textbf{\sffamily\color{accent}GRPO, DAPO, GSPO}\ \ \hyperref[c:reasoning-4]{GRPO and its successors}\,{\color{muted}\scriptsize p.\,\pageref{c:reasoning-4}}\par\vspace{2pt}
\textbf{\sffamily\color{accent}HBM3E, HBM4}\ \ \hyperref[c:hardware-8]{High-bandwidth memory}\,{\color{muted}\scriptsize p.\,\pageref{c:hardware-8}}\par\vspace{2pt}
\textbf{\sffamily\color{accent}HLE, GPQA, FrontierMath}\ \ \hyperref[c:evals-6]{Frontier knowledge benchmarks}\,{\color{muted}\scriptsize p.\,\pageref{c:evals-6}}\par\vspace{2pt}
\textbf{\sffamily\color{accent}HRM, TRM}\ \ \hyperref[c:arch-10]{Tiny recursive reasoning models}\,{\color{muted}\scriptsize p.\,\pageref{c:arch-10}}\par\vspace{2pt}
\textbf{\sffamily\color{accent}Joint-Embedding Predictive Architecture}\ \ \hyperref[c:embodied-3]{JEPA}\,{\color{muted}\scriptsize p.\,\pageref{c:embodied-3}}\par\vspace{2pt}
\textbf{\sffamily\color{accent}Mamba, Gated DeltaNet}\ \ \hyperref[c:arch-4]{Hybrid linear-attention and state space models}\,{\color{muted}\scriptsize p.\,\pageref{c:arch-4}}\par\vspace{2pt}
\textbf{\sffamily\color{accent}MCP}\ \ \hyperref[c:protocols-0]{Model Context Protocol}\,{\color{muted}\scriptsize p.\,\pageref{c:protocols-0}}\par\vspace{2pt}
\textbf{\sffamily\color{accent}METR horizon}\ \ \hyperref[c:evals-8]{Task time horizon}\,{\color{muted}\scriptsize p.\,\pageref{c:evals-8}}\par\vspace{2pt}
\textbf{\sffamily\color{accent}mHC}\ \ \hyperref[c:arch-3]{Manifold-constrained hyper-connections}\,{\color{muted}\scriptsize p.\,\pageref{c:arch-3}}\par\vspace{2pt}
\textbf{\sffamily\color{accent}MI455X}\ \ \hyperref[c:hardware-1]{AMD Helios and Open Rack Wide}\,{\color{muted}\scriptsize p.\,\pageref{c:hardware-1}}\par\vspace{2pt}
\textbf{\sffamily\color{accent}MLPerf}\ \ \hyperref[c:hardware-16]{Hardware benchmarks}\,{\color{muted}\scriptsize p.\,\pageref{c:hardware-16}}\par\vspace{2pt}
\textbf{\sffamily\color{accent}MoE}\ \ \hyperref[c:arch-0]{Mixture of Experts}\,{\color{muted}\scriptsize p.\,\pageref{c:arch-0}}\par\vspace{2pt}
\textbf{\sffamily\color{accent}MOPD}\ \ \hyperref[c:reasoning-3]{Multi-teacher on-policy distillation}\,{\color{muted}\scriptsize p.\,\pageref{c:reasoning-3}}\par\vspace{2pt}
\textbf{\sffamily\color{accent}MTP}\ \ \hyperref[c:arch-6]{Multi-token prediction and speculative decoding}\,{\color{muted}\scriptsize p.\,\pageref{c:arch-6}}\par\vspace{2pt}
\textbf{\sffamily\color{accent}MX, MXFP8, MXFP4}\ \ \hyperref[c:quant-0]{Microscaling formats}\,{\color{muted}\scriptsize p.\,\pageref{c:quant-0}}\par\vspace{2pt}
\textbf{\sffamily\color{accent}NCCL}\ \ \hyperref[c:hardware-14]{Collective communication}\,{\color{muted}\scriptsize p.\,\pageref{c:hardware-14}}\par\vspace{2pt}
\textbf{\sffamily\color{accent}NIST, ISO 42001, SB 53}\ \ \hyperref[c:safety-12]{Other regulation and standards}\,{\color{muted}\scriptsize p.\,\pageref{c:safety-12}}\par\vspace{2pt}
\textbf{\sffamily\color{accent}NSA, DSA, CSA, HCA}\ \ \hyperref[c:arch-2]{Sparse and compressed attention}\,{\color{muted}\scriptsize p.\,\pageref{c:arch-2}}\par\vspace{2pt}
\textbf{\sffamily\color{accent}NVL72}\ \ \hyperref[c:hardware-0]{Rack-scale systems}\,{\color{muted}\scriptsize p.\,\pageref{c:hardware-0}}\par\vspace{2pt}
\textbf{\sffamily\color{accent}ONNX, safetensors}\ \ \hyperref[c:hardware-15]{Model interchange formats}\,{\color{muted}\scriptsize p.\,\pageref{c:hardware-15}}\par\vspace{2pt}
\textbf{\sffamily\color{accent}OpenClaw}\ \ \hyperref[c:agents-10]{Personal agents}\,{\color{muted}\scriptsize p.\,\pageref{c:agents-10}}\par\vspace{2pt}
\textbf{\sffamily\color{accent}orchestrator-worker}\ \ \hyperref[c:agents-7]{Multi-agent systems and subagents}\,{\color{muted}\scriptsize p.\,\pageref{c:agents-7}}\par\vspace{2pt}
\textbf{\sffamily\color{accent}OSAID}\ \ \hyperref[c:practice-7]{Open weights vs. open source AI}\,{\color{muted}\scriptsize p.\,\pageref{c:practice-7}}\par\vspace{2pt}
\textbf{\sffamily\color{accent}OTel GenAI}\ \ \hyperref[c:protocols-11]{OpenTelemetry GenAI conventions}\,{\color{muted}\scriptsize p.\,\pageref{c:protocols-11}}\par\vspace{2pt}
\textbf{\sffamily\color{accent}PD disaggregation}\ \ \hyperref[c:inference-3]{Prefill/decode disaggregation}\,{\color{muted}\scriptsize p.\,\pageref{c:inference-3}}\par\vspace{2pt}
\textbf{\sffamily\color{accent}PEFT, LoRA, QLoRA}\ \ \hyperref[c:training-5]{LoRA and parameter-efficient fine-tuning}\,{\color{muted}\scriptsize p.\,\pageref{c:training-5}}\par\vspace{2pt}
\textbf{\sffamily\color{accent}PRM}\ \ \hyperref[c:reasoning-8]{Process reward models}\,{\color{muted}\scriptsize p.\,\pageref{c:reasoning-8}}\par\vspace{2pt}
\textbf{\sffamily\color{accent}QAT}\ \ \hyperref[c:quant-4]{Quantization-aware training}\,{\color{muted}\scriptsize p.\,\pageref{c:quant-4}}\par\vspace{2pt}
\textbf{\sffamily\color{accent}RAG}\ \ \hyperref[c:context-5]{Retrieval-augmented generation}\,{\color{muted}\scriptsize p.\,\pageref{c:context-5}}\par\vspace{2pt}
\textbf{\sffamily\color{accent}Ralph Wiggum technique}\ \ \hyperref[c:coding-4]{Ralph loop}\,{\color{muted}\scriptsize p.\,\pageref{c:coding-4}}\par\vspace{2pt}
\textbf{\sffamily\color{accent}ReAct}\ \ \hyperref[c:agents-3]{Agentic loop / ReAct}\,{\color{muted}\scriptsize p.\,\pageref{c:agents-3}}\par\vspace{2pt}
\textbf{\sffamily\color{accent}RLVR}\ \ \hyperref[c:reasoning-2]{RL with verifiable rewards}\,{\color{muted}\scriptsize p.\,\pageref{c:reasoning-2}}\par\vspace{2pt}
\textbf{\sffamily\color{accent}RoPE, YaRN}\ \ \hyperref[c:arch-7]{Long-context methods}\,{\color{muted}\scriptsize p.\,\pageref{c:arch-7}}\par\vspace{2pt}
\textbf{\sffamily\color{accent}RSP, ASL, Preparedness}\ \ \hyperref[c:safety-0]{Capability thresholds and frontier safety frameworks}\,{\color{muted}\scriptsize p.\,\pageref{c:safety-0}}\par\vspace{2pt}
\textbf{\sffamily\color{accent}SAEs, attribution graphs}\ \ \hyperref[c:safety-6]{Mechanistic interpretability}\,{\color{muted}\scriptsize p.\,\pageref{c:safety-6}}\par\vspace{2pt}
\textbf{\sffamily\color{accent}SDD}\ \ \hyperref[c:coding-3]{Spec-driven development}\,{\color{muted}\scriptsize p.\,\pageref{c:coding-3}}\par\vspace{2pt}
\textbf{\sffamily\color{accent}SKILL.md}\ \ \hyperref[c:protocols-5]{Agent Skills}\,{\color{muted}\scriptsize p.\,\pageref{c:protocols-5}}\par\vspace{2pt}
\textbf{\sffamily\color{accent}SLM}\ \ \hyperref[c:arch-11]{Small language models}\,{\color{muted}\scriptsize p.\,\pageref{c:arch-11}}\par\vspace{2pt}
\textbf{\sffamily\color{accent}SWE-bench, Terminal-Bench}\ \ \hyperref[c:evals-3]{Agentic coding benchmarks}\,{\color{muted}\scriptsize p.\,\pageref{c:evals-3}}\par\vspace{2pt}
\textbf{\sffamily\color{accent}tool poisoning}\ \ \hyperref[c:safety-4]{MCP and tool security}\,{\color{muted}\scriptsize p.\,\pageref{c:safety-4}}\par\vspace{2pt}
\textbf{\sffamily\color{accent}TPU, Trainium, MTIA}\ \ \hyperref[c:hardware-11]{Custom accelerators}\,{\color{muted}\scriptsize p.\,\pageref{c:hardware-11}}\par\vspace{2pt}
\textbf{\sffamily\color{accent}Triton, CUTLASS, ROCm, XLA}\ \ \hyperref[c:hardware-13]{Kernel languages and compilers}\,{\color{muted}\scriptsize p.\,\pageref{c:hardware-13}}\par\vspace{2pt}
\textbf{\sffamily\color{accent}trusted access}\ \ \hyperref[c:landscape-2]{Tiered access for dual-use models}\,{\color{muted}\scriptsize p.\,\pageref{c:landscape-2}}\par\vspace{2pt}
\textbf{\sffamily\color{accent}UEC}\ \ \hyperref[c:hardware-5]{Ultra Ethernet}\,{\color{muted}\scriptsize p.\,\pageref{c:hardware-5}}\par\vspace{2pt}
\textbf{\sffamily\color{accent}Ultra Accelerator Link}\ \ \hyperref[c:hardware-3]{UALink}\,{\color{muted}\scriptsize p.\,\pageref{c:hardware-3}}\par\vspace{2pt}
\textbf{\sffamily\color{accent}VLA}\ \ \hyperref[c:embodied-0]{Vision-language-action models}\,{\color{muted}\scriptsize p.\,\pageref{c:embodied-0}}\par\vspace{2pt}
\textbf{\sffamily\color{accent}WAM}\ \ \hyperref[c:embodied-1]{World action models}\,{\color{muted}\scriptsize p.\,\pageref{c:embodied-1}}\par\vspace{2pt}
\end{multicols}
\end{document}
